% Fonction exponentielle népérienne — Mathématiques Tle Tech — Cours signé OnBuch+
% Programme officiel MINESEC. Compiler avec : tectonic N13-fonction-exponentielle-neperienne.tex
\newcommand{\DOCMATIERE}{Mathématiques}
\newcommand{\DOCNIVEAU}{Tle Tech}
\newcommand{\DOCMODULE}{Mathématiques – classe de Terminale technique}
\newcommand{\DOCLECON}{N13}
\newcommand{\DOCTITRE}{Fonction exponentielle népérienne}
% OnBuch+ — préambule « cours signés » ENRICHI (compilé avec tectonic / XeLaTeX)
% Système visuel « Pop » d'OnBuch+ : fond crème (jamais blanc), bordures encre
% épaisses, ombres « dures » décalées, titres Archivo Black en capitales.
% Version MATHÉMATIQUES : graphes (pgfplots/tikz), boîtes pédagogiques
% (définition, propriété, méthode, attention, le savais-tu, exercice résolu,
% exercice, TP, bilan), page de garde et signature OnBuch+ ancrée en bas.
%
% Macros attendues dans main.tex AVANT \input{preamble} :
%   \DOCMATIERE  \DOCNIVEAU  \DOCMODULE  \DOCLECON (numéro)  \DOCTITRE
\documentclass[11pt]{article}

\usepackage{fontspec}
\usepackage[french]{babel}
\usepackage[a4paper,margin=2cm,top=3.4cm,bottom=2.8cm,headheight=1.4cm,footskip=1.3cm]{geometry}
\usepackage{amsmath,amssymb,amsfonts}
\usepackage{mathtools} % \xmapsto, \coloneqq... souvent utilisés par les modèles
\usepackage{cancel} % simplifications barrées
% \abs{z} (module, valeur absolue) et \norm{u} (norme), utilisés par les modèles
\providecommand{\abs}{}\let\abs\relax\DeclarePairedDelimiter{\abs}{\lvert}{\rvert}
\providecommand{\norm}{}\let\norm\relax\DeclarePairedDelimiter{\norm}{\lVert}{\rVert}
\usepackage[table]{xcolor} % [table] : fournit aussi \rowcolors (alternance de lignes)
\usepackage{pagecolor}
\usepackage{tikz}
\usetikzlibrary{arrows.meta,positioning,calc,shapes.geometric,shapes.misc,shapes.arrows,decorations.pathmorphing,decorations.pathreplacing,decorations.markings,patterns,shadows,mindmap,trees,fit,backgrounds,matrix,intersections,angles,quotes}
\usepackage{pgfplots}
\usepackage{pgf-pie} % diagrammes circulaires \pie (statistiques)
\pgfplotsset{compat=1.18}
\usepgfplotslibrary{fillbetween,groupplots} % aires entre courbes (calcul intégral)
\usepackage{enumitem}
\usepackage{fancyhdr}
% [most] charge toutes les bibliothèques tcolorbox (listings, minted, raster,
% documentation...) et gonfle inutilement la mémoire TeX sur un document long
% (~300 boîtes) : on ne charge que ce qui est réellement utilisé.
\usepackage[skins,breakable]{tcolorbox}
\usepackage{graphicx}
\usepackage{colortbl}
\usepackage{booktabs}
\usepackage{tabularx}
\usepackage{diagbox} % case d'en-tête diagonale des tableaux à double entrée
% Colonnes extensibles centrée / gauche / droite (C, L, R), souvent utilisées par les modèles
\newcolumntype{C}{>{\centering\arraybackslash}X}
\newcolumntype{L}{>{\raggedright\arraybackslash}X}
\newcolumntype{R}{>{\raggedleft\arraybackslash}X}
\usepackage{array}
\usepackage{multicol}
\usepackage{siunitx}
% parse-numbers=false : accepte aussi \SI{2,5 \times 10^{-3}}{...}
\sisetup{locale=FR,output-decimal-marker={,},group-minimum-digits=4,parse-numbers=false}
\DeclareSIUnit\ohm{\ensuremath{\symup{\Omega}}} % U+2126 absent de la police
\usepackage{qrcode}
% \degree : commande fréquemment utilisée par les modèles pour un angle ou une
% température en dehors de \SI{}{} (non fournie par les paquets chargés).
\providecommand{\degree}{^{\circ}}
% \qed (fin de démonstration), utilisé par les modèles sans amsthm
\providecommand{\qed}{\hfill$\square$}
% \barre : double barre des valeurs interdites dans un tableau de variations (tkz-tab)
\providecommand{\barre}{\ensuremath{\|}}
% environnement proof (amsthm non chargé) utilisé par les modèles
\ifdefined\proof\else\newenvironment{proof}[1][Démonstration]{\par\noindent\textit{#1.}\ }{\qed\par}\fi
\usepackage{lastpage}
\usepackage{hyperref}
\hypersetup{hidelinks}

% ============================================================
% Palette « Pop » OnBuch+ — mêmes hex que la classe P (plus_ui.dart)
% ============================================================
\definecolor{poporange}{HTML}{F59321}
\definecolor{popdark}{HTML}{E07A0C}
\definecolor{popcream}{HTML}{FFF6E8}
\definecolor{popink}{HTML}{1C1714}
\definecolor{poppurple}{HTML}{7A5AE0}
\definecolor{popgreen}{HTML}{2E9E6A}
\definecolor{poppink}{HTML}{E0457B}
\definecolor{popblue}{HTML}{2F7DE1}
\definecolor{popgold}{HTML}{C98A12}
\definecolor{popmuted}{HTML}{8A7F76}
\definecolor{popline}{HTML}{EADFD2}
\definecolor{popgray}{HTML}{8A7F76}
\colorlet{popgrey}{popgray}
% Alias tolérants (le modèle nomme parfois les couleurs différemment)
\colorlet{popwhite}{white}
\colorlet{popblack}{popink}
\colorlet{popred}{poppink}
\colorlet{popyellow}{popgold}
% Teintes claires pour fonds de tableaux / zones
\colorlet{popblueL}{popblue!10!white}
\colorlet{poporangeL}{poporange!14!white}
\colorlet{popgreenL}{popgreen!12!white}
\colorlet{poppurpleL}{poppurple!12!white}
\colorlet{poppinkL}{poppink!12!white}
\colorlet{popgoldL}{popgold!14!white}

\pagecolor{popcream}

% ============================================================
% Typographie
% ============================================================
\defaultfontfeatures{Ligatures=TeX}
\setmainfont{PlusJakartaSans-Medium.ttf}[Path=../../../fonts/,BoldFont=PlusJakartaSans-Bold.ttf]
% Maths en sans-serif (Fira Math) avec chiffres et lettres droites de Plus
% Jakarta, pour que les formules se fondent dans le texte.
\usepackage[math-style=ISO,bold-style=ISO]{unicode-math}
\setmathfont{FiraMath-Regular.otf}[Path=../../../fonts/]
\setmathfont{PlusJakartaSans-Medium.ttf}[Path=../../../fonts/,range={up/num,up/{Latin,latin},\mathup}]
\usepackage{icomma} % virgule décimale sans espace en mode math
% Caractères mathématiques Unicode absents des polices de texte (Plus Jakarta,
% Archivo Black) : rendus par leur équivalent mathématique, en texte comme en
% formule — sinon ils s'affichent en carré vide (ex. « Arithmétique dans ℤ »).
\usepackage{newunicodechar}
\newunicodechar{ℝ}{\ensuremath{\mathbb{R}}}
\newunicodechar{ℤ}{\ensuremath{\mathbb{Z}}}
\newunicodechar{ℂ}{\ensuremath{\mathbb{C}}}
\newunicodechar{ℕ}{\ensuremath{\mathbb{N}}}
\newunicodechar{ℚ}{\ensuremath{\mathbb{Q}}}
\newunicodechar{𝔻}{\ensuremath{\mathbb{D}}}
\newunicodechar{α}{\ensuremath{\alpha}}
\newunicodechar{β}{\ensuremath{\beta}}
\newunicodechar{γ}{\ensuremath{\gamma}}
\newunicodechar{δ}{\ensuremath{\delta}}
\newunicodechar{ε}{\ensuremath{\varepsilon}}
\newunicodechar{θ}{\ensuremath{\theta}}
\newunicodechar{λ}{\ensuremath{\lambda}}
\newunicodechar{μ}{\ensuremath{\mu}}
\newunicodechar{σ}{\ensuremath{\sigma}}
\newunicodechar{τ}{\ensuremath{\tau}}
\newunicodechar{φ}{\ensuremath{\varphi}}
\newunicodechar{ψ}{\ensuremath{\psi}}
\newunicodechar{ω}{\ensuremath{\omega}}
\newunicodechar{Σ}{\ensuremath{\Sigma}}
% majuscules produites par \MakeUppercase dans les titres (α-aminés -> Α)
\newunicodechar{Α}{\ensuremath{\alpha}}
\newunicodechar{Β}{\ensuremath{\beta}}
\newunicodechar{Γ}{\ensuremath{\Gamma}}
\newunicodechar{Δ}{\ensuremath{\Delta}}
\newunicodechar{Ε}{\ensuremath{\varepsilon}}
\newunicodechar{Θ}{\ensuremath{\Theta}}
\newunicodechar{Λ}{\ensuremath{\Lambda}}
\newunicodechar{Μ}{\ensuremath{\mu}}
\newunicodechar{Τ}{\ensuremath{\tau}}
\newunicodechar{Φ}{\ensuremath{\Phi}}
\newunicodechar{Ψ}{\ensuremath{\Psi}}
\newunicodechar{Ω}{\ensuremath{\Omega}}
\newunicodechar{↦}{\ensuremath{\mapsto}}
\newunicodechar{∈}{\ensuremath{\in}}
\newunicodechar{∉}{\ensuremath{\notin}}
\newunicodechar{∧}{\ensuremath{\wedge}}
\newunicodechar{⇔}{\ensuremath{\Leftrightarrow}}
\newunicodechar{⟺}{\ensuremath{\Longleftrightarrow}}
\newunicodechar{⇒}{\ensuremath{\Rightarrow}}
\newunicodechar{⟹}{\ensuremath{\Longrightarrow}}
\newunicodechar{∘}{\ensuremath{\circ}}
\newunicodechar{∪}{\ensuremath{\cup}}
\newunicodechar{∩}{\ensuremath{\cap}}
\newunicodechar{≡}{\ensuremath{\equiv}}
\newunicodechar{⊂}{\ensuremath{\subset}}
\newunicodechar{⊥}{\ensuremath{\perp}}
\newunicodechar{∃}{\ensuremath{\exists}}
\newunicodechar{∀}{\ensuremath{\forall}}
\newunicodechar{∛}{\ensuremath{\sqrt[3]{\,}}}
\newunicodechar{∜}{\ensuremath{\sqrt[4]{\,}}}
\newunicodechar{⃗}{\ensuremath{{}^{\to}}}
\newunicodechar{ⁿ}{\ifmmode^{n}\else\textsuperscript{\ensuremath{n}}\fi}
\newunicodechar{ˣ}{\ifmmode^{x}\else\textsuperscript{\ensuremath{x}}\fi}
\newunicodechar{ᵇ}{\ifmmode^{b}\else\textsuperscript{\ensuremath{b}}\fi}
\newunicodechar{ʳ}{\ifmmode^{r}\else\textsuperscript{\ensuremath{r}}\fi}
\newunicodechar{ᵘ}{\ifmmode^{u}\else\textsuperscript{\ensuremath{u}}\fi}
\newunicodechar{ʸ}{\ifmmode^{y}\else\textsuperscript{\ensuremath{y}}\fi}
\newunicodechar{ᵃ}{\ifmmode^{a}\else\textsuperscript{\ensuremath{a}}\fi}
\newunicodechar{ᵉ}{\ifmmode^{e}\else\textsuperscript{\ensuremath{e}}\fi}
\newunicodechar{ᵏ}{\ifmmode^{k}\else\textsuperscript{\ensuremath{k}}\fi}
\newunicodechar{ᵖ}{\ifmmode^{p}\else\textsuperscript{\ensuremath{p}}\fi}
\newunicodechar{ᴺ}{\ifmmode^{N}\else\textsuperscript{\ensuremath{N}}\fi}
\newunicodechar{ᶜ}{\ifmmode^{c}\else\textsuperscript{\ensuremath{c}}\fi}
\newunicodechar{ʰ}{\ifmmode^{h}\else\textsuperscript{\ensuremath{h}}\fi}
\newunicodechar{ᵠ}{\ifmmode^{q}\else\textsuperscript{\ensuremath{q}}\fi}
\newunicodechar{ₙ}{\ifmmode_{n}\else\textsubscript{\ensuremath{n}}\fi}
\newunicodechar{ₐ}{\ifmmode_{a}\else\textsubscript{\ensuremath{a}}\fi}
\newunicodechar{ᵢ}{\ifmmode_{i}\else\textsubscript{\ensuremath{i}}\fi}
\newunicodechar{ᵣ}{\ifmmode_{r}\else\textsubscript{\ensuremath{r}}\fi}
\newunicodechar{ₖ}{\ifmmode_{k}\else\textsubscript{\ensuremath{k}}\fi}
\newunicodechar{ᵦ}{\ifmmode_{\beta}\else\textsubscript{\ensuremath{\beta}}\fi}
\newunicodechar{ₓ}{\ifmmode_{x}\else\textsubscript{\ensuremath{x}}\fi}
\newfontfamily\popdisplay{ArchivoBlack-Regular.ttf}[Path=../../../fonts/]
\newfontfamily\poplogo{SpaceGrotesk-Bold.ttf}[Path=../../../fonts/]
\newfontfamily\popsemi{PlusJakartaSans-SemiBold.ttf}[Path=../../../fonts/]
\newfontfamily\popxbold{PlusJakartaSans-ExtraBold.ttf}[Path=../../../fonts/]
\newfontfamily\popxboldls{PlusJakartaSans-ExtraBold.ttf}[Path=../../../fonts/,LetterSpace=3.0]

\setlist[enumerate]{leftmargin=1.7em,itemsep=5pt,topsep=4pt}
\setlist[itemize]{leftmargin=1.7em,itemsep=4pt,topsep=4pt,label={\color{poporange}\textbullet}}
\setlength{\parindent}{0pt}
\setlength{\parskip}{5pt}
\renewcommand{\arraystretch}{1.25}

% pgfplots : style Pop par défaut
\pgfplotsset{
  every axis/.append style={axis line style={popink,line width=1pt},tick style={popink},
    grid style={popline,line width=0.5pt},label style={font=\small},tick label style={font=\footnotesize}},
  popcurve/.style={poporange,line width=1.8pt,no markers,smooth},
}
% Même style disponible dans un \draw TikZ simple (hors pgfplots)
\tikzset{popcurve/.style={poporange,line width=1.8pt}}

% ============================================================
% Pill / squiggle
% ============================================================
\newcommand{\poppill}[3]{%
  \tikz[baseline=-0.55ex]\node[draw=popink,line width=1.2pt,rounded corners=2.6pt,%
    fill=#2,inner xsep=6.5pt,inner ysep=3pt]{%
    \popxboldls\fontsize{7}{7}\selectfont\color{#3}\MakeUppercase{#1}};%
}
\newcommand{\popsquiggle}[1]{%
  \tikz[baseline=-0.4ex]{
    \draw[#1,line width=2.6pt,line cap=round]
      plot [smooth] coordinates {(0,0.09) (0.34,0.01) (0.68,0.21) (1.02,0.01) (1.36,0.10)};
  }%
}
% Mot-clé surligné (marqueur orange sous le texte)
\newcommand{\cle}[1]{{\popxbold\color{popdark}#1}}

% ============================================================
% En-tête / pied
% ============================================================
\pagestyle{fancy}
\fancyhf{}
\renewcommand{\headrulewidth}{0pt}
\renewcommand{\footrulewidth}{0pt}
\lhead{\raisebox{-0.15ex}{\poplogo\fontsize{13}{13}\selectfont\color{popink}OnBuch\color{poporange}+}}
\rhead{\poppill{\DOCMATIERE\ \textbullet\ \DOCNIVEAU\ \textbullet\ Leçon \DOCLECON}{white}{popink}}
\lfoot{\popxbold\fontsize{7.6}{7.6}\selectfont\color{popmuted}\MakeUppercase{Cours signé OnBuch+}\ \textbullet\ {\popsemi\DOCTITRE}}
\rfoot{\tikz[baseline=-0.6ex]\node[rounded corners=8.5pt,fill=popink,minimum height=17pt,minimum width=17pt,inner xsep=5pt,inner sep=0]{\popxbold\fontsize{8}{8}\selectfont\color{popcream}\thepage\,/\,\pageref*{LastPage}};}

% ============================================================
% Titres
% ============================================================
\newcommand{\chaptitre}[2]{%
  \begin{tcolorbox}[enhanced,colback=poporange,colframe=popink,boxrule=2.6pt,arc=11pt,
      drop shadow=popink,left=15pt,right=15pt,top=12pt,bottom=11pt,before skip=3pt,after skip=18pt]
    \poppill{#2}{popink}{popcream}\par\vspace{9pt}
    {\raggedright\hyphenpenalty=10000\exhyphenpenalty=10000\popdisplay\fontsize{22}{24}\selectfont\color{popcream}\MakeUppercase{#1}\par}
  \end{tcolorbox}
}
% Section numérotée : pastille orange avec le numéro + titre Archivo
\newcounter{popsec}
\newcommand{\coursec}[1]{%
  \stepcounter{popsec}\setcounter{popsub}{0}%
  \par\vspace{16pt}\noindent
  \begin{minipage}{\linewidth}
  \tikz[baseline=-0.7ex]\node[draw=popink,line width=1.6pt,fill=poporange,rounded corners=4pt,minimum size=22pt,inner sep=0]{\popdisplay\fontsize{12}{12}\selectfont\color{popcream}\thepopsec};%
  \hspace{8pt}{\popdisplay\fontsize{15}{17}\selectfont\color{popink}\MakeUppercase{#1}}\par
  \vspace{4pt}\noindent{\color{poporange}\rule{36pt}{3.6pt}}
  \end{minipage}\par\nopagebreak\vspace{8pt}
}
\newcounter{popsub}
\newcommand{\courssub}[1]{%
  \stepcounter{popsub}%
  \par\vspace{9pt}\noindent{\popxbold\fontsize{11.5}{13}\selectfont\color{popink}{\color{poporange}\thepopsec.\thepopsub}\hspace{6pt}#1}\par\nopagebreak\vspace{4pt}
}

% ============================================================
% Boîtes pédagogiques — même recette : fond blanc, bordure épaisse colorée,
% ombre dure encre, badge plein. Titre optionnel [#1].
% ============================================================
\tcbset{popbase/.style={breakable,enhanced,colback=white,boxrule=2.3pt,arc=9pt,
  shadow={3pt}{-3pt}{0pt}{popink},left=14pt,right=14pt,top=11pt,bottom=11pt,before skip=11pt,after skip=15pt,
  fontupper=\popsemi\fontsize{10.6}{14.5}\selectfont\color{popink},
  fontlower=\popsemi\fontsize{10.6}{14.5}\selectfont\color{popink}}}
% Coche dessinée (le glyphe ✓ n'existe pas dans les polices du document)
\newcommand{\popcheck}[1]{\tikz[baseline=-0.5ex]\draw[#1,line width=1.6pt,line cap=round,line join=round](-0.13,0.0)--(-0.04,-0.09)--(0.13,0.1);}
\newcommand{\popboxhead}[3]{\poppill{#1}{#2}{white}\ifx\relax#3\relax\else\hspace{6pt}{\popxbold\fontsize{10.6}{12}\selectfont\color{popink}#3}\fi\par\vspace{7pt}}

\newtcolorbox{aretenir}[1][]{popbase,colframe=popgreen,before upper={\popboxhead{À retenir}{popgreen}{#1}}}
\newtcolorbox{exemplebox}[1][]{popbase,colframe=poporange,before upper={\popboxhead{Exemple}{poporange}{#1}}}
\newtcolorbox{definition}[1][]{popbase,colframe=popblue,before upper={\popboxhead{Définition}{popblue}{#1}}}
\newtcolorbox{propriete}[1][]{popbase,colframe=poppurple,before upper={\popboxhead{Propriété}{poppurple}{#1}}}
\newtcolorbox{methode}[1][]{popbase,colframe=popgold,before upper={\popboxhead{Méthode}{popgold}{#1}}}
\newtcolorbox{attention}[1][]{popbase,colframe=poppink,before upper={\popboxhead{Attention}{poppink}{#1}}}
\newtcolorbox{experience}[1][]{popbase,colframe=popblue,colback=popblueL,before upper={\popboxhead{Expérience}{popblue}{#1}}}
% « Le savais-tu ? » : carte violette pleine (contexte, culture, Cameroun)
\newtcolorbox{savaistu}[1][]{popbase,colback=poppurple,colframe=popink,
  fontupper=\popsemi\fontsize{10.4}{14.2}\selectfont\color{white},
  before upper={\poppill{Le savais-tu\,?}{white}{poppurple}\ifx\relax#1\relax\else\hspace{6pt}{\popxbold\color{white}#1}\fi\par\vspace{7pt}}}
% Exercice résolu : énoncé en haut, \tcblower puis la solution sur fond crème
\newcounter{popexo}\newcounter{popexr}
\newtcolorbox{exoresolu}[1][]{popbase,colframe=poporange,colbacklower=popcream,
  segmentation style={popink,line width=1.2pt,dashed},
  before upper={\stepcounter{popexr}\popboxhead{Exercice résolu \thepopexr}{poporange}{#1}},
  before lower={\poppill{Solution}{popink}{popcream}\par\vspace{6pt}}}
% Exercice (non résolu en place) — la correction est dans la partie Corrigés
\newtcolorbox{exercice}[1][]{popbase,colframe=popink,
  before upper={\stepcounter{popexo}\popboxhead{Exercice \thepopexo}{popink}{#1}}}
% Corrigé
\newtcolorbox{corrige}[1][]{popbase,colframe=popgreen,colback=popgreenL,
  before upper={\popboxhead{Corrigé}{popgreen}{#1}}}
% Objectifs / prérequis / bilan
\newtcolorbox{objectifs}{popbase,colframe=popink,colback=poporangeL,
  before upper={\popboxhead{Objectifs de la leçon}{popink}{}}}
\newtcolorbox{prerequis}{popbase,colframe=popink,colback=popblueL,
  before upper={\popboxhead{Prérequis}{popblue}{}}}
\newtcolorbox{bilan}{popbase,colframe=popink,colback=white,boxrule=2.8pt,
  before upper={\popboxhead{Fiche bilan}{poporange}{L'essentiel à connaître pour l'examen}}}
% Figure encadrée avec légende
\newtcolorbox{popfigure}[1][]{enhanced,breakable=false,colback=white,colframe=popline,boxrule=1.4pt,arc=8pt,
  left=10pt,right=10pt,top=10pt,bottom=8pt,before skip=10pt,after skip=12pt,halign=center,
  lower separated=false,halign lower=center,
  fontlower=\popsemi\fontsize{9}{11.5}\selectfont\color{popmuted},
  #1}
\newcommand{\legende}[1]{\par\vspace{6pt}{\popsemi\fontsize{9}{11.5}\selectfont\color{popmuted}#1}}

% ============================================================
% Signature OnBuch+
% ============================================================
\newcommand{\poplogoinline}[1]{{\poplogo\fontsize{#1}{#1}\selectfont\color{popink}OnBuch\color{poporange}+}}
\newcommand{\signatureonbuch}{%
  \begin{tcolorbox}[enhanced,colback=popink,colframe=popink,boxrule=2.2pt,arc=10pt,
      drop shadow=poporange,left=14pt,right=14pt,top=10pt,bottom=10pt,before skip=0pt,after skip=0pt]
    \tikz[baseline=-0.7ex]\node[circle,fill=popgreen,draw=popcream,line width=1.3pt,minimum size=22pt,inner sep=0]{\popcheck{white}};%
    \hspace{9pt}\begin{minipage}[c]{0.62\linewidth}
      {\popxbold\fontsize{12}{13}\selectfont\color{popcream}Signé\ \textbullet\ {\poplogo OnBuch\color{poporange}+}}\par\vspace{2pt}
      {\raggedright\popsemi\fontsize{8}{10.5}\selectfont\color{popline}\MakeUppercase{Équipe pédagogique OnBuch+}\\\MakeUppercase{Conforme au programme officiel MINESEC}\par}
    \end{minipage}\hfill
    {\poplogo\fontsize{22}{22}\selectfont\color{popcream}OnBuch\color{poporange}+}
  \end{tcolorbox}%
}

% ============================================================
% Page de garde
% #1 titre · #2 matière · #3 niveau · #4 module · #5 n° leçon · #6 durée ·
% #7 description / objectifs (texte ou itemize)
% ============================================================
\newcommand{\pagegarde}[7]{%
  \thispagestyle{empty}
  \newgeometry{a4paper,margin=1.7cm,top=1.6cm,bottom=1.4cm}%
  \begin{tikzpicture}[remember picture,overlay]
    \fill[poporange!18] ([xshift=-3.2cm,yshift=-3.6cm]current page.north east) circle (4.6cm);
    \fill[poppurple!14] ([xshift=2.4cm,yshift=7.6cm]current page.south west) circle (3.8cm);
  \end{tikzpicture}%
  {\poplogo\fontsize{30}{30}\selectfont\color{popink}OnBuch\color{poporange}+}\hfill
  \raisebox{0.6ex}{\poppill{Cours signé \textbullet\ Premium}{popink}{popcream}}\par
  \vspace{1pt}\popsquiggle{poppurple}\par
  \vspace{8pt}
  \begin{tcolorbox}[enhanced,colback=poporange,colframe=popink,boxrule=2.8pt,arc=13pt,
      drop shadow=popink,left=18pt,right=18pt,top=12pt,bottom=13pt,after skip=10pt]
    \poppill{#2 \textbullet\ #3}{popink}{popcream}\hspace{5pt}\poppill{#4}{white}{popink}\par\vspace{8pt}
    {\popdisplay\fontsize{14}{14}\selectfont\color{popink}LEÇON #5}\par\vspace{4pt}
    {\raggedright\hyphenpenalty=10000\exhyphenpenalty=10000\popdisplay\fontsize{25}{27}\selectfont\color{popcream}\MakeUppercase{#1}\par}
    \vspace{8pt}
    {\popxbold\fontsize{9.5}{9.5}\selectfont\color{popink}\MakeUppercase{Durée conseillée : #6}}
  \end{tcolorbox}
  \begin{tcolorbox}[enhanced,colback=white,colframe=popink,boxrule=2.2pt,arc=10pt,
      drop shadow=popink,left=16pt,right=16pt,top=10pt,bottom=10pt,after skip=8pt]
    \poppill{Dans cette leçon}{poppurple}{white}\par\vspace{5pt}
    \popsemi\fontsize{9.3}{12.6}\selectfont\color{popink}#7
  \end{tcolorbox}
  \noindent\begin{minipage}[t]{0.487\linewidth}
    \begin{tcolorbox}[enhanced,colback=popgreen,colframe=popink,boxrule=2.2pt,arc=10pt,
        drop shadow=popink,left=11pt,right=11pt,top=10pt,bottom=10pt,height=3.1cm,valign=center]
      \tcbox[colback=white,colframe=popink,boxrule=1.3pt,arc=5pt,boxsep=0pt,nobeforeafter,
        left=3pt,right=3pt,top=3pt,bottom=3pt,valign=center]{\qrcode[height=1.9cm]{https://play.google.com/store/apps/details?id=com.luvvix.onbuch&hl=fr}}%
      \hspace{8pt}\begin{minipage}[c]{0.5\linewidth}
        {\popdisplay\fontsize{11}{12.5}\selectfont\color{white}\MakeUppercase{Télécharge OnBuch}}\par\vspace{4pt}
        {\raggedright\popsemi\fontsize{8.4}{11}\selectfont\color{white}Gratuit \textbullet\ Google Play\\Tuteur IA, annales, concours, résultats\par}
      \end{minipage}
    \end{tcolorbox}
  \end{minipage}\hfill
  \begin{minipage}[t]{0.487\linewidth}
    \begin{tcolorbox}[enhanced,colback=poppurple,colframe=popink,boxrule=2.2pt,arc=10pt,
        drop shadow=popink,left=11pt,right=11pt,top=10pt,bottom=10pt,height=3.1cm,valign=center]
      \tikz[baseline=-0.7ex]\node[circle,fill=white,draw=popink,line width=1.3pt,minimum size=1.6cm,inner sep=0]{\popdisplay\fontsize{24}{24}\selectfont\color{poppurple}+};%
      \hspace{8pt}\begin{minipage}[c]{0.6\linewidth}
        {\popdisplay\fontsize{11}{12.5}\selectfont\color{white}\MakeUppercase{Passe à OnBuch+}}\par\vspace{4pt}
        {\raggedright\popsemi\fontsize{8.4}{11}\selectfont\color{white}Tous les cours signés, Tuteur IA illimité, fiches PDF — onglet OnBuch+ de l'app\par}
      \end{minipage}
    \end{tcolorbox}
  \end{minipage}
  \vspace{8pt}
  \popcontenu
  \vfill
  \signatureonbuch
  \restoregeometry
  \newpage
}

% Rangée « Ce cours contient » (page de garde)
\newcommand{\poptile}[3]{% #1 couleur #2 gros texte #3 légende
  \begin{tcolorbox}[enhanced,colback=#1,colframe=popink,boxrule=2pt,arc=9pt,shadow={2.5pt}{-2.5pt}{0pt}{popink},
      left=6pt,right=6pt,top=9pt,bottom=9pt,halign=center,valign=center,height=1.75cm,width=\linewidth,nobeforeafter]
    {\popdisplay\fontsize{12.5}{13}\selectfont\color{white}\MakeUppercase{#2}}\par\vspace{3pt}
    {\centering\popsemi\fontsize{7.8}{9.5}\selectfont\color{white}#3\par}
  \end{tcolorbox}}
\newcommand{\popcontenu}{%
  \par\noindent{\popxboldls\fontsize{7.5}{7.5}\selectfont\color{popmuted}CE COURS SIGNÉ CONTIENT}\par\vspace{7pt}
  \noindent\begin{minipage}[t]{0.235\linewidth}\poptile{popblue}{Cours}{complet, illustré, pas à pas}\end{minipage}\hfill
  \begin{minipage}[t]{0.235\linewidth}\poptile{poporange}{Méthodes}{et exercices résolus}\end{minipage}\hfill
  \begin{minipage}[t]{0.235\linewidth}\poptile{poppink}{Exercices}{type examen + corrigés}\end{minipage}\hfill
  \begin{minipage}[t]{0.235\linewidth}\poptile{popgreen}{Fiche bilan}{l'essentiel à retenir}\end{minipage}\par}

% Sommaire
\newtcolorbox{sommaire}{enhanced,colback=white,colframe=popink,boxrule=2.2pt,arc=10pt,
  drop shadow=popink,left=18pt,right=18pt,top=13pt,bottom=13pt,before skip=6pt,after skip=20pt,
  before upper={\poppill{Sommaire}{poppurple}{white}\par\vspace{9pt}\popsemi\fontsize{10.6}{14.5}\selectfont\color{popink}}}

% Signature de fin de cours — poussée tout en bas de la dernière page
\newcommand{\findecours}{%
  \par\vfill\nopagebreak
  \begin{center}\popsquiggle{poporange}\end{center}\vspace{4pt}
  \signatureonbuch
}
\AtBeginDocument{\def\llbracket{\mathopen{[\mkern-3.5mu[}}\def\rrbracket{\mathclose{]\mkern-3.5mu]}}} % intervalles d'entiers (glyphe ⟦ absent de la police)
% Glyphes absents de Fira Math : redéfinis à partir de glyphes disponibles
\AtBeginDocument{%
  \def\setminus{\mathbin{\backslash}}\let\smallsetminus\setminus
  \def\vdots{\mathord{\vbox{\baselineskip3.2pt\lineskiplimit0pt\kern4pt\hbox{.}\hbox{.}\hbox{.}}}}%
  \def\ddots{\mathinner{\mkern1mu\raise6.5pt\hbox{.}\mkern2mu\raise3.6pt\hbox{.}\mkern2mu\raise0.7pt\hbox{.}\mkern1mu}}%
  \def\triangle{\mathord{\scalebox{0.9}{$\symup{\Delta}$}}}%
  \def\nabla{\mathord{\raisebox{\depth}{\scalebox{1}[-1]{$\symup{\Delta}$}}}}%
  \def\checkmark{\popcheck{popdark}}%
  \def\star{\mathbin{\ast}}\def\bigstar{\mathord{\ast}}%
  \def\diamond{\mathbin{\scalebox{0.6}{$\Diamond$}}}%
  \def\bigcup{\mathop{\mathchoice{\raisebox{-0.3em}{\scalebox{1.7}{$\cup$}}}{\scalebox{1.25}{$\cup$}}{\cup}{\cup}}\displaylimits}%
  \def\bigcap{\mathop{\mathchoice{\raisebox{-0.3em}{\scalebox{1.7}{$\cap$}}}{\scalebox{1.25}{$\cap$}}{\cap}{\cap}}\displaylimits}%
  \def\lesseqqgtr{\mathrel{\lessgtr}}\def\bigoplus{\mathop{\oplus}\displaylimits}%
}
\providecommand{\ssi}{si et seulement si} % abréviation fréquente
\AtBeginDocument{\providecommand{\wideparen}{\overparen}} % arc de cercle

%% FIGFIX-BEGIN v4 — correctifs globaux des figures (docs/onbuch-generation/tools/figfix)
\usepackage{adjustbox}\usepackage{etoolbox}
% toute figure TikZ/pgfplots plus large que la ligne est réduite à la largeur disponible
\BeforeBeginEnvironment{tikzpicture}{\begin{adjustbox}{max width=\linewidth}}
\AfterEndEnvironment{tikzpicture}{\end{adjustbox}}
% légendes pgfplots lisibles par-dessus les courbes
\pgfplotsset{every axis/.append style={legend style={fill=white,fill opacity=0.92,text opacity=1,draw=black!15,inner sep=3pt}}}
% étiquettes d'arêtes (syntaxe edge["..."]) avec fond blanc
\tikzset{every edge quotes/.append style={fill=white,inner sep=1.2pt}}
%% FIGFIX-END
\begin{document}

\pagegarde{Fonction exponentielle népérienne}{Mathématiques}{Tle Tech}{Mathématiques – classe de Terminale technique}{N13}{4 séances (fiche de progression harmonisée)}{Cette leçon introduit la fonction exponentielle de base e (exp), ses propriétés algébriques, ses limites, sa dérivée et ses primitives. Elle permet d'étudier des fonctions composées avec exp et de modéliser des situations concrètes de croissance ou décroissance continue (démographie, finance, radioactivité, chimie).
\vspace{4pt}
\begin{multicols}{2}\small
\begin{itemize}
\item Définir la fonction exponentielle népérienne exp comme la fonction réciproque du logarithme népérien ln et en déduire ses propriétés algébriques fondamentales.
\item Calculer les limites de exp en -∞ et +∞, et résoudre des équations et inéquations du type exp(u(x)) = k ou exp(u(x)) > k.
\item Démontrer que la fonction exp est dérivable sur ℝ, qu'elle est sa propre dérivée, et appliquer la règle de dérivation des composées u'(x)exp(u(x)).
\item Déterminer les primitives de fonctions de la forme u'(x)exp(u(x)) et calculer des intégrales définies associées.
\item Réaliser une étude complète (domaine, limites, dérivée, variations, tableau de signes, équation f(x)=k, graphique) d'une fonction définie à l'aide de exp.
\item Modéliser une situation concrète (croissance démographique, intérêt composé continu, décroissance radioactive, chimie) par une fonction exponentielle et interpréter les résultats (temps de doublement, demi-vie, taux horaire).
\end{itemize}\end{multicols}}

\begin{sommaire}
\begin{multicols}{2}
\begin{enumerate}
\item 1. Définition et propriétés algébriques de la fonction exp
\item 2. Limites, équations et inéquations exponentielles
\item 3. Dérivée de la fonction exp et fonctions composées
\item 4. Primitives et intégrales de fonctions exponentielles
\item 5. Étude complète de fonctions faisant intervenir exp
\item 6. Modélisation concrète : croissance et décroissance continues
\item 7. Compléments et ouverture : Éuler, complexes et suites (Culture mathématique)
\item Activité d'intégration
\item Méthodes et exercices résolus
\item Exercices
\item Corrigés des exercices
\item Fiche bilan
\end{enumerate}
\end{multicols}
\end{sommaire}

\begin{objectifs}
\begin{itemize}
\item Définir la fonction exponentielle népérienne exp comme la fonction réciproque du logarithme népérien ln et en déduire ses propriétés algébriques fondamentales.
\item Calculer les limites de exp en -∞ et +∞, et résoudre des équations et inéquations du type exp(u(x)) = k ou exp(u(x)) > k.
\item Démontrer que la fonction exp est dérivable sur ℝ, qu'elle est sa propre dérivée, et appliquer la règle de dérivation des composées u'(x)exp(u(x)).
\item Déterminer les primitives de fonctions de la forme u'(x)exp(u(x)) et calculer des intégrales définies associées.
\item Réaliser une étude complète (domaine, limites, dérivée, variations, tableau de signes, équation f(x)=k, graphique) d'une fonction définie à l'aide de exp.
\item Modéliser une situation concrète (croissance démographique, intérêt composé continu, décroissance radioactive, chimie) par une fonction exponentielle et interpréter les résultats (temps de doublement, demi-vie, taux horaire).
\item Rédiger un raisonnement mathématique rigoureux en justifiant chaque étape (théorèmes des valeurs intermédiaires, de la bijection, de dérivation).
\item Utiliser la calculatrice ou un logiciel (GeoGebra, Python) pour conjecturer des variations, vérifier des limites ou tracer des courbes représentatives.
\end{itemize}
\end{objectifs}

\begin{prerequis}
\begin{itemize}
\item Fonction logarithme népérien ln : définition comme primitive de 1/x sur ]0;+∞[, propriétés algébriques, limites, dérivée, variations, équations ln(x)=k.
\item Notion de fonction réciproque (bijection) : conditions d'existence, propriétés graphiques (symétrie par rapport à y=x), lien dérivées réciproques.
\item Dérivation des fonctions usuelles (puissances, ln) et règles de dérivation (somme, produit, quotient, composée).
\item Recherche de primitives usuelles et calcul d'intégrales définies (Théorème fondamental de l'analyse).
\item Méthodologie d'étude complète de fonction (domaine, parité, limites, dérivée, tableau de variations, tableau de signes, graphique).
\end{itemize}
\end{prerequis}

\newpage

\coursec{Définition et propriétés algébriques de la fonction exp}

\textbf{Situation d'accroche.}\par
À Yaoundé, la population d'un quartier populaire double environ tous les 20 ans. Un élève de Terminale Tech veut prévoir le nombre d'habitants dans 30 ans pour dimensionner un forage d'eau. Comment modéliser cette croissance rapide et continue ? La fonction exponentielle népérienne est l'outil mathématique incontournable pour décrire de tels phénomènes d'évolution proportionnelle à la quantité présente.

\courssub{Rappel : la fonction logarithme népérien, une bijection}

Avant de définir l'exponentielle, rappelons le résultat fondamental établi en Première : la fonction logarithme népérien $\ln$ est une \textbf{bijection} de l'intervalle $]0;+\infty[$ sur $\mathbb{R}$. Elle est \textbf{strictement croissante} et continue. Ses propriétés algébriques sont :
\[
\forall (a,b)\in ]0;+\infty[^2,\quad \ln(ab)=\ln a+\ln b \quad;\quad \ln\left(\frac{a}{b}\right)=\ln a-\ln b \quad;\quad \ln(a^r)=r\ln a\ (r\in\mathbb{Q}).
\]
Cette bijectivité garantit l'existence d'une fonction réciproque, que nous allons étudier.

\courssub{Définition de la fonction exponentielle népérienne}

\begin{definition}[Fonction exponentielle népérienne]
La fonction \textbf{exponentielle népérienne}, notée $\exp$, est la fonction réciproque de la fonction logarithme népérien $\ln$. Elle est définie par :
\[
\exp : \mathbb{R} \longrightarrow ]0;+\infty[
\]
et vérifie les deux équivalences fondamentales, pour tout $x\in\mathbb{R}$ et tout $y>0$ :
\[
\boxed{\ln(\exp(x)) = x} \qquad \text{et} \qquad \boxed{\exp(\ln(y)) = y}.
\]
\end{definition}

\begin{propriete}[Notation alternative et nombre $e$]
On note $e = \exp(1)$. Ce nombre irrationnel, appelé \textbf{constante d'Euler}, vaut approximativement $e \approx 2,718\,281\,828\ldots$
Par extension, pour tout réel $x$, on écrit couramment :
\[
\exp(x) = e^x.
\]
Cette notation de puissance sera justifiée par les propriétés algébriques ci-dessous.
\end{propriete}

\begin{popfigure}
\begin{tikzpicture}
\begin{axis}[
    width=\linewidth, height=0.7\linewidth,
    axis lines=middle,
    xlabel=$x$, ylabel=$y$,
    xmin=-2.5, xmax=3.5,
    ymin=-2.5, ymax=3.5,
    xtick={-2,-1,1,2,3}, ytick={-2,-1,1,2,3},
    minor tick num=1,
    grid=both, grid style={popline},
    major grid style={popline},
    axis line style={popdark, thick},
    tick style={popdark},
    label style={popdark},
    clip=false,
    domain=-2:1.6, samples=200,
    restrict y to domain=-2.5:3.5,
]
% Droite y=x (axe de symétrie)
\addplot[popmuted, dashed, thick, domain=-2.5:3.5] {x};
\node[popmuted, anchor=south west] at (axis cs:3,3) {$y=x$};

% Courbe Ln (x>0)
\addplot[popblue, very thick, domain=0.02:3.5] {ln(x)};
\node[popblue, anchor=south east] at (axis cs:3.5,1.25) {$C_{\ln}\ :\ y=\ln(x)$};

% Courbe Exp
\addplot[poporange, very thick, domain=-2:1.3] {exp(x)};
\node[poporange, anchor=north west] at (axis cs:-2,3) {$C_{\exp}\ :\ y=\exp(x)$};

% Tangente en (0,1) : y = x+1
\addplot[popgreen, thick, dashed, domain=-1.5:2.5] {x+1};
\node[popgreen, anchor=south west] at (axis cs:2,3) {$T_0\ :\ y=x+1$};

% Points caractéristiques
\fill[poporange] (axis cs:0,1) circle (2.5pt) node[above right, popdark] {$(0,1)$};
\fill[popblue] (axis cs:1,0) circle (2.5pt) node[below right, popdark] {$(1,0)$};
\fill[popdark] (axis cs:0,0) circle (1.5pt);

% Flèches de symétrie (optionnel, illustration)
\draw[poppurple, <->, thick] (axis cs:0.5, -0.693) -- (axis cs:-0.693, 0.5) node[midway, above left, poppurple, font=\small] {symétrie};
\end{axis}
\end{tikzpicture}
\legende{Graphiques de $C_{\ln}$ et $C_{\exp}$ : symétrie par rapport à $y=x$, point $(0,1)$ sur $C_{\exp}$, tangente $y=x+1$ illustrant $\exp(x) \ge x+1$.}
\end{popfigure}

\courssub{Propriétés algébriques fondamentales}

Les propriétés de $\exp$ découlent directement de celles de $\ln$ par passage à la fonction réciproque.

\begin{propriete}[Propriétés algébriques de $\exp$]
Pour tous réels $a$ et $b$ :
\begin{enumerate}
    \item $\exp(0) = 1$ \quad (car $\ln(1)=0$).
    \item $\exp(a+b) = \exp(a) \times \exp(b)$.
    \item $\exp(a-b) = \dfrac{\exp(a)}{\exp(b)}$.
    \item $\exp(-x) = \dfrac{1}{\exp(x)}$.
    \item $\forall r \in \mathbb{Q},\ \exp(r) = e^r$ (cohérence avec la notation puissance).
\end{enumerate}
\end{propriete}

\begin{experience}[Démonstration guidée des propriétés algébriques]
On utilise le fait que $\ln$ est bijective : $u=v \iff \ln(u)=\ln(v)$ pour $u,v>0$.
\begin{enumerate}
    \item \textbf{$\exp(0)=1$ :} $\ln(\exp(0)) = 0 = \ln(1) \implies \exp(0)=1$.
    \item \textbf{Produit :} Soient $A=\exp(a+b)$ et $B=\exp(a)\exp(b)$. Calculons leur logarithme :
    \[
    \ln(A) = \ln(\exp(a+b)) = a+b.
    \]
    \[
    \ln(B) = \ln(\exp(a)\exp(b)) = \ln(\exp(a)) + \ln(\exp(b)) = a+b.
    \]
    Comme $\ln$ est injective, $A=B$, d'où $\exp(a+b)=\exp(a)\exp(b)$.
    \item \textbf{Quotient et opposé :} Déduits immédiatement du produit en posant $b=-b$ ou $a=0$.
    \item \textbf{Puissances rationnelles :} Pour $r=\frac{p}{q}\in\mathbb{Q}$, $\ln(\exp(r)^q) = q\ln(\exp(r)) = qr = p = \ln(e^p) = \ln((e^r)^q)$. Par injectivité, $\exp(r)^q = e^p$, donc $\exp(r) = e^r$.
\end{enumerate}
\end{experience}

\begin{exoresolu}[Calcul sans calculatrice]
Calculer $E = \exp(\ln(3)) - \ln(\exp(2)) + \exp(0)$.
\tcblower
\textbf{Solution.}
On applique directement les relations de réciprocité et la valeur en $0$ :
\begin{align*}
\exp(\ln(3)) &= 3 \quad (\text{car } 3>0), \\
\ln(\exp(2)) &= 2 \quad (\text{car } 2\in\mathbb{R}), \\
\exp(0) &= 1.
\end{align*}
Donc $E = 3 - 2 + 1 = \boxed{2}$.
\end{exoresolu}

\courssub{Inégalité fondamentale et croissance}

\begin{propriete}[Inégalité fondamentale]
Pour tout réel $x$ :
\[
\boxed{\exp(x) \ge x + 1}
\]
avec égalité \textbf{si et seulement si} $x=0$.
\end{propriete}

\begin{experience}[Démonstration guidée de l'inégalité]
On utilise l'inégalité du logarithme : $\forall u > -1,\ \ln(1+u) \le u$ (égalité si $u=0$).
\begin{enumerate}
    \item Posons $u = \exp(x) - 1$. Comme $\exp(x) > 0$, on a $u > -1$.
    \item Appliquons l'inégalité du logarithme :
    \[
    \ln(1 + u) \le u \iff \ln(1 + \exp(x) - 1) \le \exp(x) - 1.
    \]
    \item Simplifions : $\ln(\exp(x)) \le \exp(x) - 1 \iff x \le \exp(x) - 1$.
    \item En ajoutant $1$ : $x + 1 \le \exp(x)$.
    \item L'égalité $\ln(1+u)=u$ a lieu seulement pour $u=0$, c'est-à-dire $\exp(x)-1=0 \iff \exp(x)=1 \iff x=0$.
\end{enumerate}
\end{experience}

\begin{propriete}[Croissance stricte]
La fonction $\exp$ est \textbf{strictement croissante} sur $\mathbb{R}$.
\emph{Justification :} $\ln$ est strictement croissante sur $]0;+\infty[$, sa fonction réciproque $\exp$ l'est donc aussi sur $\mathbb{R}$.
\end{propriete}

\courssub{Puissances rationnelles et cohérence}

Pour tout rationnel $r = \frac{p}{q}$ ($p\in\mathbb{Z}, q\in\mathbb{N}^*$), on a prouvé que $\exp(r) = e^r$. Cela justifie pleinement l'écriture $e^x$ pour tout réel $x$ : la fonction $x \mapsto e^x$ \emph{est} la fonction $\exp$, unique prolongement continu de la fonction puissance rationnelle.

\courssub{Pièges à éviter}

\begin{attention}[Erreurs fréquentes]
\begin{itemize}
    \item \textbf{Confusion $\exp(x)$ vs $x^e$ :} $\exp(x) = e^x$ (base constante $e$, exposant variable $x$). Ne pas confondre avec la fonction puissance $x \mapsto x^e$ (base variable, exposant constant).
    \item \textbf{Signe de $\exp(x)$ :} $\exp(x) > 0$ pour \textbf{tout} réel $x$. Jamais $\exp(x) \le 0$. L'équation $\exp(x) = -2$ n'a \textbf{aucune solution}.
    \item \textbf{Notation $\exp(x)^2$ :} Écrire $\exp(x)^2$ est ambigu. Préférer $(\exp(x))^2$ ou $\exp(2x)$ (car $\exp(2x) = (\exp(x))^2$).
    \item \textbf{$\ln(\exp(x)) = x$ toujours vrai} (pour tout $x\in\mathbb{R}$), mais \textbf{$\exp(\ln(x)) = x$ seulement si $x>0$}.
\end{itemize}
\end{attention}

\courssub{Le savais-tu ?}

\begin{savaistu}[Le nombre $e$ : une constante universelle]
\begin{itemize}
    \item \textbf{Transcendance :} En 1873, Charles Hermite démontre que $e$ est \textbf{transcendant} (il n'est racine d'aucun polynôme à coefficients entiers non nul). C'était le premier nombre « naturel » prouvé transcendant (avant $\pi$ par Lindemann en 1882).
    \item \textbf{Intérêt composé continu :} Si l'on place un capital $C_0$ à un taux $t$ (ex: $5\%$ soit $t=0,05$) avec capitalisation \emph{continue}, le capital au bout d'un an est $C_0 e^t$. La formule $e = \lim_{n\to\infty} (1+\frac{1}{n})^n$ modélise une capitalisation infiniment fréquente.
    \item \textbf{Partout dans la technique :} Désintégration radioactive, charge/décharge d'un condensateur (électricité), loi de refroidissement de Newton, croissance bactérienne, amortissement d'emprunts... Partout où la variation instantanée est proportionnelle à la quantité présente, $e$ apparaît.
\end{itemize}
\end{savaistu}

\coursec{Limites, équations et inéquations exponentielles}

Dans la section précédente, nous avons défini la fonction exponentielle népérienne $\exp$ comme l'unique fonction dérivable sur $\mathbb{R}$ vérifiant $f' = f$ et $f(0) = 1$, et nous avons établi ses propriétés algébriques fondamentales : $\exp(a+b) = \exp(a)\times\exp(b)$, $\exp(-a) = \dfrac{1}{\exp(a)}$, ainsi que le fait que $\exp(x) > 0$ pour tout réel $x$. Nous savons aussi que $\exp$ est \cle{strictement croissante} sur $\mathbb{R}$.

Mais pour étudier des phénomènes concrets — la charge d'un condensateur, la croissance d'une production, la décroissance d'un produit chimique — il faut savoir \emph{ce que devient} $\exp(x)$ quand $x$ devient très grand ou très négatif, et savoir \emph{résoudre} des équations du type $\exp(u(x)) = k$. C'est l'objet de cette section : d'abord les \cle{limites}, puis les \cle{équations} et \cle{inéquations}, enfin une application directe aux temps de doublement et aux demi-vies.

\courssub{Limites de la fonction exponentielle aux bornes}

\textbf{Intuition.} Observons quelques valeurs : $\exp(1) \approx 2,72$ ; $\exp(5) \approx 148$ ; $\exp(10) \approx 22026$ ; $\exp(20) \approx 4,85\times 10^{8}$. Quand $x$ augmente, $\exp(x)$ explose littéralement. À l'inverse, $\exp(-5) \approx 0,0067$ et $\exp(-10) \approx 0,0000454$ : quand $x$ devient très négatif, $\exp(x)$ se rapproche de $0$ sans jamais l'atteindre (puisque $\exp(x) > 0$).

\begin{propriete}[Limites de la fonction exponentielle]
\[
\lim_{x \to -\infty} \exp(x) = 0 \qquad \text{et} \qquad \lim_{x \to +\infty} \exp(x) = +\infty.
\]
\end{propriete}

\begin{aretenir}[Conséquence graphique : l'asymptote]
La première limite signifie que la droite d'équation $y = 0$ (l'axe des abscisses) est une \cle{asymptote horizontale} à la courbe de $\exp$ en $-\infty$ : la courbe « colle » à l'axe des abscisses à gauche, sans jamais le toucher. En $+\infty$, il n'y a aucune asymptote : la courbe monte indéfiniment, de plus en plus vite.
\end{aretenir}

\begin{popfigure}
\begin{center}
\begin{tikzpicture}
\begin{axis}[
    width=0.9\linewidth, height=7cm,
    axis lines=middle,
    xlabel={$x$}, ylabel={$y$},
    xmin=-4.5, xmax=3.2, ymin=-0.5, ymax=8,
    xtick={-4,-3,-2,-1,0,1,2,3},
    ytick={1,2.718,5},
    yticklabels={$1$,$e$,$5$},
    grid=both, grid style={popline!40},
    legend style={at={(0.02,0.98)}, anchor=north west, draw=none, fill=none}
]
\addplot[popcurve, domain=-4.5:2.05, samples=200]{exp(x)};
\addlegendentry{$y = \exp(x)$}
\addplot[poporange, dashed, thick, domain=-4.5:3.2, samples=2]{0};
\addlegendentry{asymptote $y=0$}
\addplot[poppurple, dashed, thick, domain=-4.5:3.2, samples=2]{5};
\addlegendentry{$y=k$ (ici $k=5$)}
\addplot[only marks, mark=*, popdark, mark size=2pt] coordinates {(0,1) (1.609,5)};
\node[popdark, anchor=south west] at (axis cs:0.1,1.1) {$(0\,;\,1)$};
\node[poppurple, anchor=south] at (axis cs:1.609,5.15) {$\bigl(\ln k\,;\,k\bigr)$};
\draw[poppurple, dashed] (axis cs:1.609,0) -- (axis cs:1.609,5);
\node[poppurple, anchor=north] at (axis cs:1.609,-0.1) {$\ln k$};
\end{axis}
\end{tikzpicture}
\legende{Courbe de $y=\exp(x)$ : asymptote $y=0$ en $-\infty$, passage par $(0\,;\,1)$, croissance explosive à droite. L'équation $\exp(x)=k$ se lit graphiquement comme l'intersection avec l'horizontale $y=k$ : l'unique solution est $x=\ln k$.}
\end{center}
\end{popfigure}

\courssub{Croissance comparée : $\exp(x)$ face aux puissances de $x$}

\textbf{Intuition.} En $+\infty$, qui « gagne » entre $\exp(x)$ et une puissance $x^n$ ? Comparons : pour $x = 100$, on a $x^3 = 10^{6}$, mais $\exp(100) \approx 2,7 \times 10^{43}$. L'exponentielle écrase toutes les puissances, même $x^{100}$. C'est un résultat essentiel pour lever les formes indéterminées du type « $\dfrac{\infty}{\infty}$ » ou « $\infty - \infty$ ».

\begin{propriete}[Croissance comparée (admise)]
Pour tout entier naturel $n$ non nul :
\[
\lim_{x \to +\infty} \frac{x^n}{\exp(x)} = 0 \qquad \text{et, de manière équivalente,} \qquad \lim_{x \to +\infty} \frac{\exp(x)}{x^n} = +\infty.
\]
On dit que \emph{l'exponentielle domine toute puissance de $x$ en $+\infty$}. Ce résultat est admis en Terminale technique ; sa démonstration fait appel à des outils plus avancés.
\end{propriete}

\begin{exemplebox}[Lever une forme indéterminée]
Calculer $\displaystyle\lim_{x \to +\infty} \bigl(\exp(x) - x^2\bigr)$. C'est une forme indéterminée « $\infty - \infty$ ». On factorise par le terme dominant :
\[
\exp(x) - x^2 = \exp(x)\left(1 - \frac{x^2}{\exp(x)}\right).
\]
Or $\dfrac{x^2}{\exp(x)} \to 0$ par croissance comparée, donc la parenthèse tend vers $1$, et $\exp(x) \to +\infty$. Conclusion : $\displaystyle\lim_{x \to +\infty} \bigl(\exp(x) - x^2\bigr) = +\infty$.
\end{exemplebox}

\begin{savaistu}[Pourquoi les algorithmes exponentiels sont inutilisables]
En informatique, on mesure le coût d'un algorithme par le nombre d'opérations en fonction de la taille $n$ des données. Un algorithme en $n^2$ (tri naïf) traitant $n = 10^{6}$ données demande environ $10^{12}$ opérations : long mais faisable. Un algorithme en $\exp(n)$ demanderait $\exp(10^{6})$ opérations, un nombre astronomique dépassant tout ce qu'un ordinateur peut faire. La relation $\exp(n) \gg n^k$ explique pourquoi les problèmes dits « exponentiels » (comme le voyageur de commerce par force brute) restent hors de portée dès que $n$ dépasse quelques dizaines. C'est toute la différence entre un problème « facile » et un problème « intraitable ».
\end{savaistu}

\courssub{Résolution d'équations du type $\exp(u(x)) = k$}

\textbf{Intuition.} La fonction $\exp$ est strictement croissante et prend toutes les valeurs de $]0\,;\,+\infty[$. Donc l'horizontale $y = k$ coupe la courbe en \emph{un seul point} si $k > 0$, et \emph{jamais} si $k \leq 0$ (voir la figure ci-dessus). L'outil qui permet de « descendre » de l'exponentielle est sa fonction réciproque, le logarithme népérien $\ln$.

\begin{propriete}[Équation exponentielle]
Soit $u$ une fonction définie sur un ensemble $D$ et $k$ un réel.
\begin{itemize}
\item Si $k > 0$ : $\exp(u(x)) = k \iff u(x) = \ln(k)$.
\item Si $k \leq 0$ : l'équation $\exp(u(x)) = k$ n'a \textbf{aucune solution}.
\end{itemize}
Cas particuliers à connaître par cœur : $\exp(x) = 1 \iff x = 0$ et $\exp(x) = e \iff x = 1$.
\end{propriete}

\begin{methode}[Résoudre $\exp(u(x)) = k$]
\begin{enumerate}
\item Vérifier le signe de $k$. Si $k \leq 0$, conclure immédiatement : pas de solution.
\item Si $k > 0$, appliquer $\ln$ aux deux membres : l'équation devient $u(x) = \ln(k)$.
\item Résoudre cette équation « ordinaire » en $x$.
\item Vérifier que les solutions appartiennent à l'ensemble de définition, puis conclure par l'ensemble $S$.
\end{enumerate}
\end{methode}

\begin{attention}[Pièges à éviter absolument]
\begin{itemize}
\item \textbf{Ne jamais écrire} « $\exp(x) = 0$ donc $x = \ln(0)$ » : $\ln(0)$ n'existe pas, et l'équation $\exp(x) = 0$ est tout simplement impossible car $\exp(x) > 0$ pour tout $x$.
\item \textbf{Toujours vérifier la condition $k > 0$} avant d'appliquer $\ln$. Écrire $\ln(-3)$ est une erreur grave qui coûte des points au Baccalauréat technique.
\item Ne pas confondre : $\exp(x) = k$ donne $x = \ln k$, et non $x = \dfrac{k}{e}$ ni $x = e^{k}$.
\end{itemize}
\end{attention}

\begin{exoresolu}[Résoudre $\exp(2x-1) = 5$]
Résoudre dans $\mathbb{R}$ l'équation $\exp(2x-1) = 5$.
\tcblower
\textbf{Étape 1.} $k = 5 > 0$ : la condition est satisfaite, l'équation a une unique solution.

\textbf{Étape 2.} On applique $\ln$ : $2x - 1 = \ln 5$.

\textbf{Étape 3.} On résout : $2x = 1 + \ln 5$, donc $x = \dfrac{1 + \ln 5}{2}$.

\textbf{Étape 4.} Valeur approchée : $\ln 5 \approx 1,609$, donc $x \approx \dfrac{2,609}{2} \approx 1,305$.
\[
S = \left\{ \frac{1 + \ln 5}{2} \right\}.
\]
\end{exoresolu}

\courssub{Changement de variable $X = \exp(x)$}

Certaines équations font intervenir $\exp(2x)$ et $\exp(x)$. L'astuce consiste à remarquer que $\exp(2x) = \bigl(\exp(x)\bigr)^2$ : en posant $X = \exp(x)$, on retombe sur une équation du second degré, que l'on sait résoudre depuis la classe de Première technique.

\begin{methode}[Équation quadratique en $\exp(x)$]
\begin{enumerate}
\item Poser $X = \exp(x)$, avec la contrainte fondamentale $X > 0$.
\item Résoudre l'équation du second degré en $X$ (discriminant $\Delta$).
\item \textbf{Écarter toute racine $X \leq 0$} (impossible, car $\exp(x) > 0$).
\item Pour chaque racine $X > 0$, conclure par $x = \ln X$.
\end{enumerate}
\end{methode}

\begin{exoresolu}[Résoudre $\exp(2x) - 3\exp(x) + 2 = 0$]
Résoudre dans $\mathbb{R}$ : $\exp(2x) - 3\exp(x) + 2 = 0$.
\tcblower
\textbf{Étape 1.} Posons $X = \exp(x)$, donc $X > 0$. Comme $\exp(2x) = X^2$, l'équation devient :
\[
X^2 - 3X + 2 = 0.
\]
\textbf{Étape 2.} Discriminant : $\Delta = (-3)^2 - 4\times 1\times 2 = 9 - 8 = 1 > 0$. Racines : $X_1 = \dfrac{3-1}{2} = 1$ et $X_2 = \dfrac{3+1}{2} = 2$.

\textbf{Étape 3.} Les deux racines sont strictement positives : on les garde toutes les deux.

\textbf{Étape 4.} $\exp(x) = 1 \iff x = 0$ ; $\exp(x) = 2 \iff x = \ln 2$.
\[
S = \{0 \,;\, \ln 2\}.
\]
\textbf{Vérification} pour $x = \ln 2$ : $\exp(2\ln 2) = \bigl(\exp(\ln 2)\bigr)^2 = 4$, $3\exp(\ln 2) = 6$, et $4 - 6 + 2 = 0$. Correct.
\end{exoresolu}

\begin{exoresolu}[Résoudre $\exp(x) + 2\exp(-x) = 3$]
Résoudre dans $\mathbb{R}$ : $\exp(x) + 2\exp(-x) = 3$.
\tcblower
\textbf{Étape 1.} Comme $\exp(-x) = \dfrac{1}{\exp(x)}$, posons $X = \exp(x)$ ($X > 0$). L'équation devient :
\[
X + \frac{2}{X} = 3.
\]
\textbf{Étape 2.} On multiplie les deux membres par $X > 0$ (licite, le sens des égalités est conservé) : $X^2 + 2 = 3X$, soit $X^2 - 3X + 2 = 0$.

\textbf{Étape 3.} C'est la même équation que précédemment : $X = 1$ ou $X = 2$, toutes deux strictement positives.

\textbf{Étape 4.} $x = \ln 1 = 0$ ou $x = \ln 2$.
\[
S = \{0 \,;\, \ln 2\}.
\]
\textbf{Vérification} pour $x = 0$ : $\exp(0) + 2\exp(0) = 1 + 2 = 3$. Correct.
\end{exoresolu}

\courssub{Résolution d'inéquations exponentielles}

\textbf{Intuition.} La fonction $\exp$ est \cle{strictement croissante} : elle conserve donc l'ordre. Dire $\exp(a) < \exp(b)$ revient exactement à dire $a < b$. Comme $k = \exp(\ln k)$ pour $k > 0$, on peut comparer $\exp(u(x))$ à $k$ en comparant $u(x)$ à $\ln k$.

\begin{propriete}[Inéquations exponentielles]
Soit $u$ une fonction et $k > 0$. Alors :
\[
\exp(u(x)) < k \iff u(x) < \ln(k) \qquad \text{et} \qquad \exp(u(x)) > k \iff u(x) > \ln(k).
\]
Si $k \leq 0$ : l'inéquation $\exp(u(x)) < k$ est impossible, et l'inéquation $\exp(u(x)) > k$ est vraie pour tout $x$ de l'ensemble de définition.
\end{propriete}

\begin{attention}[Le sens de l'inégalité ne change jamais ici]
Contrairement à la multiplication par un nombre négatif, appliquer $\ln$ (fonction strictement croissante) ou $\exp$ aux deux membres d'une inégalité \textbf{conserve toujours le sens}. Le seul point de vigilance reste la condition $k > 0$ pour pouvoir écrire $\ln k$.
\end{attention}

\begin{exoresolu}[Résoudre $\exp(3x + 1) \geq 7$]
Résoudre dans $\mathbb{R}$ : $\exp(3x+1) \geq 7$.
\tcblower
$k = 7 > 0$ : on applique $\ln$, qui conserve le sens de l'inégalité :
\[
3x + 1 \geq \ln 7 \iff 3x \geq \ln 7 - 1 \iff x \geq \frac{\ln 7 - 1}{3}.
\]
Numériquement : $\ln 7 \approx 1,946$, donc $x \geq \dfrac{0,946}{3} \approx 0,315$.
\[
S = \left[ \frac{\ln 7 - 1}{3} \,;\, +\infty \right[.
\]
\end{exoresolu}

\courssub{Application : temps de doublement et demi-vie}

De nombreuses situations professionnelles suivent un modèle exponentiel : un capital placé à intérêts composés en continu, la charge d'un condensateur, l'effectif d'une population de bactéries, la décroissance d'un produit chimique ou radioactif.

\textbf{Croissance.} Une quantité évolue selon $Q(t) = Q_0 \exp(kt)$ avec $k > 0$. Le \cle{temps de doublement} $T_2$ est le temps nécessaire pour que la quantité double : $Q(T_2) = 2Q_0$. On résout :
\[
Q_0 \exp(kT_2) = 2Q_0 \iff \exp(kT_2) = 2 \iff kT_2 = \ln 2 \iff T_2 = \frac{\ln 2}{k}.
\]

\textbf{Décroissance.} Une quantité décroît selon $Q(t) = Q_0 \exp(-\lambda t)$ avec $\lambda > 0$. La \cle{demi-vie} $T_{1/2}$ est le temps au bout duquel il ne reste que la moitié : $Q(T_{1/2}) = \dfrac{Q_0}{2}$. On résout :
\[
\exp(-\lambda T_{1/2}) = \frac{1}{2} \iff -\lambda T_{1/2} = \ln\frac{1}{2} = -\ln 2 \iff T_{1/2} = \frac{\ln 2}{\lambda}.
\]

\begin{aretenir}[Une formule remarquable]
Dans les deux cas, le temps caractéristique est $\dfrac{\ln 2}{\text{taux}}$, avec $\ln 2 \approx 0,693$. Il ne dépend \textbf{pas} de la quantité initiale $Q_0$ : c'est la signature des phénomènes exponentiels.
\end{aretenir}

\begin{exoresolu}[Demi-vie d'un produit chimique]
En atelier de traitement de surface, la concentration d'un produit suit $C(t) = 8\exp(-0,05t)$, où $t$ est en heures et $C$ en g/L. Déterminer la demi-vie du produit, puis au bout de combien de temps la concentration passe sous $1$ g/L.
\tcblower
\textbf{Demi-vie.} Ici $\lambda = 0,05$, donc :
\[
T_{1/2} = \frac{\ln 2}{0,05} \approx \frac{0,693}{0,05} \approx 13,9 \text{ heures}.
\]
Toutes les $13,9$ heures environ, la concentration est divisée par deux.

\textbf{Seuil de $1$ g/L.} On résout $8\exp(-0,05t) < 1$ :
\[
\exp(-0,05t) < \frac{1}{8} \iff -0,05t < \ln\frac{1}{8} = -\ln 8 \iff t > \frac{\ln 8}{0,05}.
\]
(Attention : à la dernière étape, on divise par $-0,05 < 0$, donc le sens de l'inégalité change.) Or $\ln 8 = 3\ln 2 \approx 2,079$, donc $t > \dfrac{2,079}{0,05} \approx 41,6$ heures. La concentration passe sous $1$ g/L au bout d'environ $41,6$ heures — soit exactement trois demi-vies, puisque $\left(\tfrac12\right)^3 = \tfrac18$.
\end{exoresolu}

\begin{exercice}[Pour s'entraîner]
\begin{enumerate}
\item Résoudre dans $\mathbb{R}$ : a) $\exp(x+2) = 3$ ; b) $\exp(2x) = -1$ ; c) $\exp(2x) - 5\exp(x) + 6 = 0$.
\item Résoudre dans $\mathbb{R}$ : a) $\exp(1 - x) \leq 4$ ; b) $\exp(x^2) > e$.
\item Un capital suit $C(t) = 5000\exp(0,07t)$ (en FCFA, $t$ en années). Calculer le temps de doublement du capital, arrondi au dixième d'année.
\end{enumerate}
\end{exercice}

\begin{corrige}[Exercice « Pour s'entraîner »]
\textbf{1.a)} $x + 2 = \ln 3$, donc $S = \{\ln 3 - 2\}$. \quad \textbf{1.b)} $-1 \leq 0$ : pas de solution, $S = \emptyset$. \quad \textbf{1.c)} Avec $X = \exp(x)$ : $X^2 - 5X + 6 = 0$, $\Delta = 25 - 24 = 1$, $X = 2$ ou $X = 3$, toutes deux strictement positives. Donc $S = \{\ln 2 \,;\, \ln 3\}$.

\textbf{2.a)} $1 - x \leq \ln 4 \iff x \geq 1 - \ln 4$ : $S = [1 - \ln 4 \,;\, +\infty[$. \quad \textbf{2.b)} $\exp(x^2) > e \iff x^2 > 1 \iff x < -1$ ou $x > 1$ : $S = \,]-\infty\,;\,-1[ \cup ]1\,;\,+\infty[$.

\textbf{3.} $T_2 = \dfrac{\ln 2}{0,07} \approx \dfrac{0,693}{0,07} \approx 9,9$ années.
\end{corrige}

En résumé, nous savons désormais décrire le comportement de $\exp$ aux bornes de $\mathbb{R}$, lever les formes indéterminées grâce à la croissance comparée, et résoudre toute équation ou inéquation exponentielle, y compris par changement de variable. Il reste à exploiter l'outil le plus puissant du programme : la \emph{dérivation}, qui fera l'objet de la section suivante et permettra l'étude complète de fonctions construites avec $\exp$.

\coursec{Dérivée de la fonction exp et fonctions composées}

Dans la section précédente, nous avons étudié les limites de la fonction exponentielle népérienne et résolu des équations et inéquations de la forme $\exp(a) = \exp(b)$ ou $\exp(a) > \exp(b)$. Nous savons donc comment se comporte $\exp$ aux bornes de son domaine et comment comparer des images. Il reste une question fondamentale pour l'étude des fonctions : \emph{comment varie la fonction exponentielle, et à quelle vitesse ?} La réponse tient en une propriété remarquable, qui fait toute la richesse de cette fonction : la fonction exponentielle est \cle{égale à sa propre dérivée}. Cette section établit ce résultat, le généralise aux fonctions composées $\exp(u(x))$, et montre comment l'exploiter pour étudier des variations, chercher des extrema et déterminer des tangentes.

\courssub{Dérivabilité de exp sur $\mathbb{R}$ et calcul de $\exp'$}

\textbf{Intuition.} Tracons mentalement la courbe de $y = \exp(x)$. Elle est lisse, sans angle ni rupture : on s'attend donc à ce que la fonction soit dérivable partout. De plus, on observe que plus $x$ est grand, plus la courbe monte vite : la pente de la tangente semble croître en même temps que l'ordonnée du point. Nous allons démontrer que cette intuition est exacte : \textbf{la pente en tout point est égale à l'ordonnée du point}.

\begin{propriete}[Dérivée de la fonction exponentielle]
La fonction exponentielle népérienne est dérivable sur $\mathbb{R}$ et, pour tout réel $x$ :
\[
\exp'(x) = \exp(x).
\]
Autrement dit, $\exp$ est sa propre dérivée : $\exp' = \exp$.
\end{propriete}

\begin{experience}[Démonstration guidée : dérivée de la fonction réciproque]
La fonction $\exp$ est la bijection réciproque de $\ln$. Nous allons utiliser un résultat général : si $g$ est la réciproque d'une fonction $f$ dérivable, alors $g$ est dérivable là où la dérivée de $f$ ne s'annule pas, et
\[
g'(x) = \frac{1}{f'\bigl(g(x)\bigr)}.
\]
\textbf{Étape 1.} La fonction $\ln$ est dérivable sur $]0\,;+\infty[$ et $(\ln)'(y) = \dfrac{1}{y}$. Cette dérivée ne s'annule jamais sur $]0\,;+\infty[$.

\textbf{Étape 2.} Appliquons la formule de la dérivée de la réciproque avec $f = \ln$ et $g = \exp$. Pour tout réel $x$ :
\[
\exp'(x) = \frac{1}{(\ln)'\bigl(\exp(x)\bigr)}.
\]

\textbf{Étape 3.} Or $(\ln)'(y) = \dfrac{1}{y}$, donc $(\ln)'\bigl(\exp(x)\bigr) = \dfrac{1}{\exp(x)}$. Par suite :
\[
\exp'(x) = \frac{1}{\dfrac{1}{\exp(x)}} = \exp(x).
\]
Comme $\exp(x)$ est défini pour tout réel $x$, la fonction $\exp$ est dérivable sur $\mathbb{R}$ tout entier, et $\exp' = \exp$. \hfill$\blacksquare$
\end{experience}

\begin{attention}[Une dérivée qui ne s'annule jamais]
Comme $\exp(x) > 0$ pour tout réel $x$, on a $\exp'(x) > 0$ pour tout $x$ : la fonction $\exp$ est \textbf{strictement croissante sur $\mathbb{R}$}, ce que nous savions déjà, mais que la dérivation confirme élégamment. Retenir aussi : la dérivée de $\exp$ ne s'annule \emph{jamais}.
\end{attention}

\courssub{La seule fonction égale à sa dérivée}

\begin{propriete}[Caractérisation des fonctions égales à leur dérivée]
Une fonction $f$ dérivable sur $\mathbb{R}$ vérifie $f' = f$ si et seulement s'il existe une constante réelle $C$ telle que, pour tout réel $x$ :
\[
f(x) = C\,\exp(x).
\]
En particulier, $\exp$ est l'\textbf{unique} fonction dérivable sur $\mathbb{R}$ vérifiant $f' = f$ et $f(0) = 1$.
\end{propriete}

\begin{experience}[Démonstration guidée]
\textbf{Étape 1 (sens facile).} Si $f(x) = C\exp(x)$, alors $f'(x) = C\exp(x) = f(x)$ : la condition est bien vérifiée.

\textbf{Étape 2 (sens difficile).} Supposons $f' = f$. Posons $g(x) = f(x)\exp(-x)$. Dérivons $g$ comme un produit :
\begin{align*}
g'(x) &= f'(x)\exp(-x) + f(x)\times\bigl(-\exp(-x)\bigr) \\
&= f(x)\exp(-x) - f(x)\exp(-x) = 0.
\end{align*}
La dérivée de $g$ est nulle sur $\mathbb{R}$, donc $g$ est constante : il existe $C$ tel que $f(x)\exp(-x) = C$, c'est-à-dire $f(x) = C\exp(x)$. \hfill$\blacksquare$
\end{experience}

\begin{savaistu}[L'équation $y' = y$ : le moteur du monde vivant et industriel]
L'équation $y' = y$ est la plus simple des \emph{équations différentielles}. Elle modélise \textbf{tout processus dont le taux de variation instantané est proportionnel à la quantité présente} : croissance d'une population de bactéries, intérêts composés en continu sur un compte épargne à Douala, échauffement d'un moteur, désintégration radioactive, charge d'un condensateur en électricité. C'est pourquoi la fonction exponentielle apparaît dans presque tous les métiers techniques : dès qu'une grandeur « croît d'autant plus vite qu'elle est grande », $\exp$ n'est pas loin.
\end{savaistu}

\courssub{Dérivée d'une fonction composée : $(\exp(u))' = u'\exp(u)$}

En pratique, on dérive rarement $\exp(x)$ seul : on rencontre $\exp(3x)$, $\exp(x^2+1)$, $\exp(-x)$\ldots{} On applique alors la \cle{règle de dérivation des fonctions composées} (règle de la chaîne).

\begin{propriete}[Dérivée de $\exp \circ\, u$]
Soit $u$ une fonction dérivable sur un intervalle $I$. Alors la fonction $x \mapsto \exp\bigl(u(x)\bigr)$ est dérivable sur $I$ et :
\[
\bigl(\exp(u(x))\bigr)' = u'(x)\,\exp\bigl(u(x)\bigr).
\]
On écrit souvent, en abrégé : $(e^u)' = u'\,e^u$.
\end{propriete}

\begin{methode}[Dériver une fonction de la forme $\exp(u(x))$]
\begin{enumerate}
\item \textbf{Identifier} la fonction $u$ « à l'intérieur » de l'exponentielle.
\item \textbf{Calculer} la dérivée $u'(x)$ séparément.
\item \textbf{Écrire} le résultat : $u'(x)\times\exp\bigl(u(x)\bigr)$, sans modifier l'intérieur de l'exponentielle.
\item \textbf{Simplifier} éventuellement $u'(x)$ (factoriser, réduire).
\end{enumerate}
\end{methode}

\begin{exemplebox}[Premières applications de la règle de la chaîne]
\begin{itemize}
\item $f(x) = \exp(3x)$ : ici $u(x) = 3x$, $u'(x) = 3$, donc $f'(x) = 3\exp(3x)$.
\item $g(x) = \exp(-x)$ : $u(x) = -x$, $u'(x) = -1$, donc $g'(x) = -\exp(-x)$.
\item $h(x) = \exp(x^2)$ : $u(x) = x^2$, $u'(x) = 2x$, donc $h'(x) = 2x\,\exp(x^2)$.
\item $k(x) = \exp\bigl(\ln(x)\bigr) = x$ pour $x > 0$ : on retrouve bien $k'(x) = 1$, et la règle de la chaîne donne $\dfrac{1}{x}\exp(\ln x) = \dfrac{1}{x}\times x = 1$. Cohérent !
\end{itemize}
\end{exemplebox}

\begin{attention}[Pièges fréquents de dérivation]
\begin{itemize}
\item \textbf{Ne pas confondre} $\bigl(\exp(x)\bigr)' = \exp(x)$ avec la dérivée d'une puissance $(x^n)' = n x^{n-1}$. La fonction $\exp(x)$ \textbf{n'est pas} une puissance usuelle : la variable $x$ est \emph{en exposant}, pas en base. Il n'existe pas de formule du type « on descend l'exposant » pour $\exp$.
\item \textbf{Ne pas oublier le facteur $u'(x)$} : écrire $\bigl(\exp(3x)\bigr)' = \exp(3x)$ est faux ; la bonne réponse est $3\exp(3x)$.
\item \textbf{Ne pas dériver l'intérieur} : l'exponentielle garde $u(x)$ tel quel, c'est devant que $u'(x)$ apparaît en facteur.
\end{itemize}
\end{attention}

\courssub{Dérivation de produits et de quotients avec exp}

La fonction $\exp$ se combine souvent avec des polynômes. On utilise alors les formules classiques, valables pour toutes fonctions dérivables :
\[
(uv)' = u'v + uv' \qquad\text{et}\qquad \left(\frac{u}{v}\right)' = \frac{u'v - uv'}{v^2}.
\]

\begin{exoresolu}[Dériver $f(x) = (x^2 - 1)\exp(x)$]
\textbf{Énoncé.} Soit $f$ la fonction définie sur $\mathbb{R}$ par $f(x) = (x^2 - 1)\exp(x)$. Calculer $f'(x)$ et le factoriser.
\tcblower
\textbf{Solution.} $f$ est un produit $u \times v$ avec $u(x) = x^2 - 1$ et $v(x) = \exp(x)$. On a $u'(x) = 2x$ et $v'(x) = \exp(x)$. D'après la formule du produit :
\begin{align*}
f'(x) &= u'(x)v(x) + u(x)v'(x) \\
&= 2x\,\exp(x) + (x^2 - 1)\exp(x) \\
&= \bigl(2x + x^2 - 1\bigr)\exp(x) \\
&= (x^2 + 2x - 1)\exp(x).
\end{align*}
On a mis $\exp(x)$ en facteur : c'est presque toujours la bonne stratégie, car $\exp(x) > 0$ et n'influence donc pas le signe de $f'$.
\end{exoresolu}

\begin{exoresolu}[Dériver $f(x) = \exp(3x^2 + 2x)$]
\textbf{Énoncé.} Soit $f$ définie sur $\mathbb{R}$ par $f(x) = \exp(3x^2 + 2x)$. Calculer $f'(x)$.
\tcblower
\textbf{Solution.} C'est une composée $\exp \circ\, u$ avec $u(x) = 3x^2 + 2x$, donc $u'(x) = 6x + 2$. Par la règle de la chaîne :
\[
f'(x) = (6x + 2)\,\exp(3x^2 + 2x).
\]
L'intérieur de l'exponentielle reste inchangé ; seul le facteur $(6x+2)$ apparaît devant.
\end{exoresolu}

\begin{exemplebox}[{Un quotient : $f(x) = \dfrac{\exp(x)}{x}$ sur $]0\,;+\infty[$}]
Avec $u(x) = \exp(x)$ et $v(x) = x$ :
\[
f'(x) = \frac{\exp(x)\times x - \exp(x)\times 1}{x^2} = \frac{(x-1)\exp(x)}{x^2}.
\]
Comme $\exp(x) > 0$ et $x^2 > 0$, le signe de $f'$ est celui de $(x-1)$ : $f$ décroît sur $]0\,;1]$ puis croît sur $[1\,;+\infty[$, avec un minimum en $x = 1$ valant $f(1) = e$.
\end{exemplebox}

\courssub{Variations et extrema : le signe de $f'$ se lit sur $u'$}

Voici le point le plus utile en pratique. Si $f(x) = \exp\bigl(u(x)\bigr)$, alors $f'(x) = u'(x)\exp\bigl(u(x)\bigr)$. Comme $\exp\bigl(u(x)\bigr) > 0$ pour tout $x$, on obtient le principe fondamental :

\begin{aretenir}[Règle d'or du signe]
Pour $f = \exp \circ\, u$ :
\[
\text{signe de } f' \;=\; \text{signe de } u'.
\]
En particulier, $f'(x) = 0 \iff u'(x) = 0$. Les variations de $\exp(u)$ sont \textbf{exactement celles de $u$} : $\exp$ « transmet » les variations sans les déformer.
\end{aretenir}

\begin{methode}[Étudier les variations d'une fonction faisant intervenir exp]
\begin{enumerate}
\item Dériver $f$ en appliquant la règle de la chaîne et/ou la formule du produit.
\item Factoriser $f'(x)$ en mettant l'exponentielle (toujours strictement positive) en facteur.
\item Conclure que le signe de $f'$ est celui du \textbf{facteur restant} (souvent un polynôme).
\item Dresser le tableau de signes de ce facteur, puis le tableau de variations de $f$.
\item Identifier les extrema : points où $f'$ s'annule en changeant de signe.
\end{enumerate}
\end{methode}

\begin{exoresolu}[Extremum de $f(x) = x\exp(-x)$]
\textbf{Énoncé.} Soit $f$ définie sur $\mathbb{R}$ par $f(x) = x\exp(-x)$. Étudier les variations de $f$ et déterminer ses extrema éventuels.
\tcblower
\textbf{Solution.} $f$ est un produit avec $u(x) = x$ et $v(x) = \exp(-x)$. On a $u'(x) = 1$ et, par la règle de la chaîne, $v'(x) = -\exp(-x)$. Donc :
\begin{align*}
f'(x) &= 1\times\exp(-x) + x\times\bigl(-\exp(-x)\bigr) \\
&= (1 - x)\exp(-x).
\end{align*}
Comme $\exp(-x) > 0$ pour tout $x$, le signe de $f'(x)$ est celui de $(1 - x)$ :
\begin{itemize}
\item $f'(x) > 0$ pour $x < 1$ : $f$ est strictement croissante sur $]-\infty\,;1]$ ;
\item $f'(1) = 0$ ;
\item $f'(x) < 0$ pour $x > 1$ : $f$ est strictement décroissante sur $[1\,;+\infty[$.
\end{itemize}
\begin{center}
\begin{tabular}{|c|ccccc|}\hline
$x$ & $-\infty$ & & $1$ & & $+\infty$ \\ \hline
$f'(x)$ & & $+$ & $0$ & $-$ & \\ \hline
$f$ & $-\infty$ & $\nearrow$ & $\dfrac{1}{e}$ & $\searrow$ & $0$ \\ \hline
\end{tabular}
\end{center}
$f$ admet donc un \textbf{maximum} en $x = 1$, valant $f(1) = 1\times\exp(-1) = \dfrac{1}{e} \approx 0{,}37$.
\end{exoresolu}

\courssub{Tangente à la courbe de exp}

\begin{propriete}[Équation de la tangente à $\mathcal{C}_{\exp}$]
La tangente à la courbe représentative de $\exp$ au point d'abscisse $a$ a pour équation :
\[
y = \exp(a)\,(x - a) + \exp(a).
\]
En particulier, au point $(0\,;1)$, la tangente a pour équation $y = x + 1$.
\end{propriete}

En effet, l'équation générale de la tangente est $y = f'(a)(x-a) + f(a)$, et ici $f'(a) = f(a) = \exp(a)$ : le coefficient directeur et l'ordonnée du point de contact sont \emph{égaux}, ce qui est la signature géométrique de la propriété $\exp' = \exp$.

\begin{popfigure}
\begin{center}
\begin{tikzpicture}
\begin{axis}[
  width=0.85\linewidth, height=6.5cm,
  axis lines=middle, xlabel=$x$, ylabel=$y$,
  xmin=-2.5, xmax=2.5, ymin=-1, ymax=5,
  xtick={-2,-1,0,1,2}, ytick={1,2.718},
  yticklabels={$1$,$e$},
  grid=both, grid style={popline!40},
  legend style={at={(0.02,0.98)}, anchor=north west, draw=none, fill=none}
]
\addplot[popcurve,domain=-2.5:1.6,samples=200]{exp(x)};
\addlegendentry{$y=\exp(x)$}
\addplot[poporange,thick,domain=-1.5:2.2,samples=2]{x+1};
\addlegendentry{Tangente : $y=x+1$}
\addplot[only marks,mark=*,popdark] coordinates {(0,1)};
\node[popdark,anchor=south west] at (axis cs:0.05,1.05) {$(0\,;1)$};
\addplot[only marks,mark=*,poppurple] coordinates {(1,2.718)};
\node[poppurple,anchor=south] at (axis cs:1,2.718) {pente $=$ ordonnée $=e$};
\end{axis}
\end{tikzpicture}
\end{center}
\legende{Courbe de $y=\exp(x)$ et sa tangente en $(0\,;1)$, d'équation $y = x+1$. En tout point, la pente de la tangente est égale à l'ordonnée du point : c'est la traduction graphique de $\exp' = \exp$.}
\end{popfigure}

\begin{popfigure}
\begin{center}
\begin{tikzpicture}[scale=1.0, >=Stealth]
% Schéma des pentes inverses
\draw[->] (-0.5,0) -- (6.5,0) node[right] {$x$};
\draw[->] (0,-0.5) -- (0,4.5) node[above] {$y$};
\draw[popblue, very thick, domain=0.15:2.4, samples=100] plot (\x,{exp(\x/1.6)}) node[right] {$y=\exp(x)$};
\draw[poppurple, very thick, domain=0.4:5.8, samples=100] plot (\x,{1.6*ln(\x)}) node[above right] {$y=\ln(x)$};
\draw[dashed, popmuted] (0.3,0.3) -- (4.2,4.2) node[pos=0.9, below right] {$y=x$};
% point sur exp et symétrique sur ln
\fill[popdark] (1.6,{exp(1)}) circle (2pt);
\fill[popdark] ({exp(1)},1.6) circle (2pt);
\node[popdark, anchor=south east] at (1.6,{exp(1)}) {$M$};
\node[popdark, anchor=north west] at ({exp(1)},1.6) {$M'$};
\node[align=center, popink] at (4.9,0.9) {pentes\\ inverses};
\end{tikzpicture}
\end{center}
\legende{Les courbes de $\exp$ et de $\ln$ sont symétriques par rapport à la droite $y=x$. Aux points symétriques $M$ et $M'$, les pentes des tangentes sont inverses l'une de l'autre : c'est le contenu géométrique de la formule de dérivation d'une fonction réciproque.}
\end{popfigure}

\begin{exercice}[À vous de jouer]
\begin{enumerate}
\item Dériver les fonctions suivantes : $f(x) = \exp(2x-5)$ ; $g(x) = x^2\exp(x)$ ; $h(x) = \dfrac{\exp(x)}{x^2+1}$.
\item Déterminer l'équation de la tangente à $\mathcal{C}_{\exp}$ au point d'abscisse $a = 1$.
\item Étudier les variations de $f(x) = \exp(x^2 - 4x)$ sur $\mathbb{R}$ et préciser son extremum.
\end{enumerate}
\end{exercice}

\begin{corrige}[Exercice 3]
\begin{enumerate}
\item $f'(x) = 2\exp(2x-5)$ (chaîne, $u'=2$). $g'(x) = 2x\exp(x) + x^2\exp(x) = (x^2+2x)\exp(x) = x(x+2)\exp(x)$ (produit). $h'(x) = \dfrac{\exp(x)(x^2+1) - 2x\exp(x)}{(x^2+1)^2} = \dfrac{(x-1)^2\,\exp(x)}{(x^2+1)^2}$ (quotient).
\item $y = \exp(1)(x-1) + \exp(1)$, soit $y = e\,x$.
\item $f'(x) = (2x-4)\exp(x^2-4x)$. Comme $\exp > 0$, le signe est celui de $2x-4$ : $f$ décroît sur $]-\infty\,;2]$ et croît sur $[2\,;+\infty[$. Minimum en $x=2$ : $f(2) = \exp(-4) = \dfrac{1}{e^4}$.
\end{enumerate}
\end{corrige}

\begin{aretenir}[L'essentiel de la section]
\begin{itemize}
\item $\exp$ est dérivable sur $\mathbb{R}$ et $\exp' = \exp$ (démonstration par la dérivée de la réciproque).
\item $f' = f \iff f(x) = C\exp(x)$ : l'exponentielle est la seule fonction (à constante près) égale à sa dérivée.
\item Règle de la chaîne : $\bigl(\exp(u)\bigr)' = u'\exp(u)$.
\item Comme $\exp > 0$, le signe de $f'$ se lit sur le facteur qui accompagne l'exponentielle ; pour $\exp(u)$, les variations sont celles de $u$.
\item Tangente en $a$ : $y = \exp(a)(x-a) + \exp(a)$ ; en $0$ : $y = x+1$.
\end{itemize}
\end{aretenir}

La section suivante exploitera la propriété $\exp' = \exp$ « dans l'autre sens » : puisque dériver $\exp$ redonne $\exp$, une primitive de $\exp$ est... $\exp$ elle-même. Nous étudierons les primitives et les intégrales de fonctions exponentielles.

\coursec{Primitives et intégrales de fonctions exponentielles}

\courssub{Rappel et transition depuis la dérivation}

Dans la section précédente, nous avons établi que la fonction exponentielle népérienne $\exp$ (notée aussi $e^x$) est l'unique fonction dérivable sur $\mathbb{R}$ qui est égale à sa propre dérivée et qui vaut $1$ en $0$ :
\[
\forall x\in\mathbb{R},\quad \exp'(x)=\exp(x).
\]
Cette propriété fondamentale fait de $\exp$ sa propre primitive (à une constante près). L'intégration, opération inverse de la dérivation, devient donc particulièrement simple pour les fonctions exponentielles et leurs composées. Nous allons voir comment exploiter cette simplicité pour calculer des primitives, des intégrales définies, des aires et même modéliser des phénomènes de fiabilité.

\begin{propriete}[Primitives de $\exp$ et de $u'\exp(u)$]
Soit $u$ une fonction dérivable sur un intervalle $I$.
\begin{enumerate}
\item $\displaystyle \int \exp(x)\,dx = \exp(x) + C,\quad \forall x\in\mathbb{R}$.
\item $\displaystyle \int u'(x)\exp(u(x))\,dx = \exp(u(x)) + C,\quad \forall x\in I$.
\end{enumerate}
\end{propriete}

\begin{experience}[Justification par dérivation]
Vérifions la deuxième formule en dérivant $F(x)=\exp(u(x))$ :
\[
F'(x) = u'(x)\exp(u(x)) \quad\text{(règle de la chaîne)}.
\]
Donc $F$ est bien une primitive de $u'\exp(u)$. La première formule est le cas particulier $u(x)=x$ ($u'(x)=1$).
\end{experience}

\courssub{Primitives de $\exp(u(x))$ : la méthode du « $u'$ devant »}

La formule $\int u'(x)\exp(u(x))\,dx = \exp(u(x))+C$ ne s'applique \textbf{que si le facteur $u'(x)$ est présent exactement} devant l'exponentielle. Si ce n'est pas le cas, on ne peut pas « inventer » $u'$ sans compenser.

\begin{exemplebox}[Cas direct : $u'$ est présent]
Calculer $\displaystyle \int (2x+3)\exp(x^2+3x)\,dx$.
\par\medskip
Ici $u(x)=x^2+3x$, donc $u'(x)=2x+3$. Le facteur $(2x+3)$ est bien devant $\exp$.
\[
\int (2x+3)\exp(x^2+3x)\,dx = \exp(x^2+3x) + C.
\]
\end{exemplebox}

\begin{attention}[Piège fréquent : « multiplier et diviser par $u'$ »]
On ne peut \textbf{pas} écrire :
\[
\int \exp(x^2)\,dx \neq \frac{1}{2x}\exp(x^2)+C \quad\text{(FAUX !)}
\]
car $\frac{1}{2x}$ n'est pas une constante et sa dérivée n'est pas nulle. La primitive de $\exp(x^2)$ n'est pas exprimable avec les fonctions usuelles (voir \emph{Le savais-tu ?} plus loin).
\end{attention}

\begin{propriete}[Primitive de $\exp(ax+b)$ — Cas très fréquent au Probatoire]
Pour $a\neq 0$ et $b\in\mathbb{R}$ :
\[
\int \exp(ax+b)\,dx = \frac{1}{a}\exp(ax+b) + C.
\]
\end{propriete}

\begin{exemplebox}[Application : $\int \exp(3x-2)\,dx$]
$u(x)=3x-2$, $u'(x)=3$. On a $\exp(3x-2) = \frac{1}{3}\cdot 3\exp(3x-2)$.
\[
\int \exp(3x-2)\,dx = \frac{1}{3}\exp(3x-2) + C.
\]
Vérification : $\left(\frac{1}{3}\exp(3x-2)\right)' = \frac{1}{3}\cdot 3\exp(3x-2) = \exp(3x-2)$.
\end{exemplebox}

\courssub{Intégrales définies et calcul d'aires}

\begin{propriete}[Intégrale définie de $u'\exp(u)$]
Si $u$ est dérivable et continue sur $[a,b]$ :
\[
\int_a^b u'(x)\exp(u(x))\,dx = \Bigl[\exp(u(x))\Bigr]_a^b = \exp(u(b)) - \exp(u(a)).
\]
\end{propriete}

\begin{exoresolu}[Calcul d'intégrales définies]
\textbf{Exercice 1.} Calculer $\displaystyle I_1 = \int_0^1 (2x+3)\exp(x^2+3x)\,dx$.\par
\textbf{Exercice 2.} Calculer $\displaystyle I_2 = \int_0^1 x\exp(x)\,dx$ (nécessite une intégration par parties).
\tcblower
\textbf{Solution 1.} $u(x)=x^2+3x$, $u(0)=0$, $u(1)=4$.
\[
I_1 = \Bigl[\exp(x^2+3x)\Bigr]_0^1 = \exp(4) - \exp(0) = e^4 - 1.
\]

\textbf{Solution 2.} Intégration par parties : $\int f g' = [fg] - \int f' g$.
Posons $f(x)=x \Rightarrow f'(x)=1$ et $g'(x)=\exp(x) \Rightarrow g(x)=\exp(x)$.
\[
\begin{aligned}
I_2 &= \Bigl[x\exp(x)\Bigr]_0^1 - \int_0^1 1\cdot\exp(x)\,dx \\
&= (1\cdot e^1 - 0\cdot e^0) - \Bigl[\exp(x)\Bigr]_0^1 \\
&= e - (e - 1) = 1.
\end{aligned}
\]
\end{exoresolu}

\begin{popfigure}
\begin{tikzpicture}[scale=0.9]
\begin{axis}[
    axis lines=middle,
    xlabel=$x$, ylabel=$y$,
    domain=0:1.5, samples=200,
    xmin=0, xmax=1.5, ymin=0, ymax=5,
    xtick={0,1}, xticklabels={$0$,$1$},
    ytick={2.718}, yticklabels={$e$},
    width=\linewidth, height=7cm,
    clip=false
]
\addplot[popcurve, domain=0:1.5] {exp(x)};
\addplot[poporange, fill=poporange!30, domain=0:1] {exp(x)} \closedcycle;
\addplot[dashed, popmuted] coordinates {(1,0) (1,2.718)};
\node[popdark, anchor=south] at (axis cs:0.5,3.5) {Aire $= \int_0^1 e^x\,dx = e-1$};
\node[popdark, anchor=north east] at (axis cs:1.5,4.48) {$y=e^x$};
\end{axis}
\end{tikzpicture}
\legende{Aire sous la courbe $y=\exp(x)$ entre $0$ et $1$.}
\end{popfigure}

\courssub{Intégration par parties : $\int P(x)\exp(x)\,dx$}

Lorsque l'intégrande est un produit d'un polynôme $P(x)$ par $\exp(x)$, la méthode systématique est l'intégration par parties répétée : on dérive le polynôme (qui finit par s'annuler) et on intègre $\exp(x)$ (qui reste inchangé).

\begin{methode}[Intégration par parties pour $\int P(x)\exp(x)\,dx$]
\begin{enumerate}
\item Poser $f(x)=P(x)$ (à dériver) et $g'(x)=\exp(x)$ (à intégrer).
\item Appliquer $\int f g' = fg - \int f' g$.
\item Répéter l'opération sur $\int f' g$ tant que le polynôme n'est pas réduit à une constante.
\item Factoriser $\exp(x)$ à la fin pour simplifier l'expression.
\end{enumerate}
\end{methode}

\begin{exemplebox}[Exemple complet : $\int (x^2-2x+3)\exp(x)\,dx$]
\begin{align*}
\int (x^2-2x+3)e^x\,dx 
&= (x^2-2x+3)e^x - \int (2x-2)e^x\,dx \\
&= (x^2-2x+3)e^x - \Bigl[(2x-2)e^x - \int 2e^x\,dx\Bigr] \\
&= (x^2-2x+3)e^x - (2x-2)e^x + 2e^x + C \\
&= e^x\bigl(x^2-2x+3 -2x+2 + 2\bigr) + C \\
&= e^x(x^2-4x+7) + C.
\end{align*}
\end{exemplebox}

\begin{attention}[Erreurs classiques en intégration par parties]
\begin{itemize}
\item Oublier le signe $-$ devant l'intégrale restante : $\int fg' = fg - \int f'g$.
\item Ne pas dériver correctement le polynôme (signe, coefficient).
\item Oublier la constante $C$ pour une primitive indéfinie.
\item Ne pas vérifier en dérivant le résultat final.
\end{itemize}
\end{attention}

\begin{popfigure}
\begin{tikzpicture}[scale=0.9]
\begin{axis}[
    axis lines=middle,
    xlabel=$x$, ylabel=$y$,
    domain=0:1.2, samples=200,
    xmin=0, xmax=1.2, ymin=0, ymax=3.5,
    xtick={0,1}, xticklabels={$0$,$1$},
    ytick={1,2.718}, yticklabels={$1$,$e$},
    width=\linewidth, height=7cm,
    clip=false
]
% Courbe y = x e^x
\addplot[poppurple, thick, domain=0:1.2] {x*exp(x)};
% Rectangle pour l'intégration par parties : [x e^x]_0^1
\addplot[poppurple!50, fill=poppurple!20, domain=0:1] {x*exp(x)} \closedcycle;
\addplot[dashed, popmuted] coordinates {(1,0) (1,2.718)};
% Courbe y = e^x (dérivée de x e^x moins e^x)
\addplot[popgreen, thick, domain=0:1.2] {exp(x)};
\addplot[popgreen!50, fill=popgreen!20, domain=0:1] {exp(x)} \closedcycle;
\node[poppurple, anchor=south west] at (axis cs:0.2,2.5) {$\int_0^1 x e^x\,dx$};
\node[popgreen, anchor=north east] at (axis cs:1.1,1.5) {$\int_0^1 e^x\,dx$};
\node[popdark, anchor=north] at (axis cs:0.6,3.2) {$[x e^x]_0^1 = e$};
\end{axis}
\end{tikzpicture}
\legende{Illustration géométrique de $\int_0^1 x e^x\,dx = [x e^x]_0^1 - \int_0^1 e^x\,dx = e - (e-1) = 1$. L'aire violette (sous $x e^x$) égale l'aire du rectangle $e$ moins l'aire verte (sous $e^x$).}
\end{popfigure}

\courssub{Applications : aire, valeur moyenne, fiabilité}

\begin{propriete}[{Valeur moyenne d'une fonction sur $[a,b]$}]
La valeur moyenne de $f$ continue sur $[a,b]$ est
\[
m = \frac{1}{b-a}\int_a^b f(x)\,dx.
\]
\end{propriete}

\begin{exemplebox}[{Valeur moyenne de $\exp(x)$ sur $[0,1]$}]
\[
m = \frac{1}{1-0}\int_0^1 e^x\,dx = \Bigl[e^x\Bigr]_0^1 = e - 1 \approx 1,718.
\]
\end{exemplebox}

\begin{savaistu}[Loi exponentielle et fiabilité — Culture probabiliste]
En ingénierie de fiabilité, la durée de vie $T$ d'un composant suit souvent une \emph{loi exponentielle} de paramètre $\lambda>0$ (taux de défaillance). Sa densité de probabilité est
\[
f(t) = \lambda e^{-\lambda t},\quad t\ge 0.
\]
La probabilité totale vaut $1$ :
\[
\int_0^{+\infty} \lambda e^{-\lambda t}\,dt = \Bigl[-e^{-\lambda t}\Bigr]_0^{+\infty} = 0 - (-1) = 1.
\]
L'espérance de vie (moyenne) est $\mathbb{E}[T] = \frac{1}{\lambda}$. C'est un modèle fondamental pour la maintenance préventive, la garantie des pièces, la gestion des stocks de rechange.
\end{savaistu}

\begin{exoresolu}[Application concrète : fiabilité d'un composant électronique]
Un condensateur a une durée de vie modélisée par une loi exponentielle de paramètre $\lambda = 0,0002$ h$^{-1}$ (soit une espérance de vie de $5000$ h).
\begin{enumerate}
\item Calculer la probabilité qu'il fonctionne encore après $10\,000$ h.
\item Calculer la probabilité qu'il tombe en panne entre $2000$ h et $8000$ h.
\end{enumerate}
\tcblower
\textbf{1.} $P(T>10000) = \int_{10000}^{+\infty} 0,0002\,e^{-0,0002 t}\,dt = \Bigl[-e^{-0,0002 t}\Bigr]_{10000}^{+\infty} = e^{-2} \approx 0,135$ (13,5\%).
\par\textbf{2.} $P(2000<T<8000) = \int_{2000}^{8000} 0,0002\,e^{-0,0002 t}\,dt = e^{-0,4} - e^{-1,6} \approx 0,670 - 0,202 = 0,468$ (46,8\%).
\end{exoresolu}

\courssub{Complément : l'intégrale de Gauss et la fonction Erf}

\begin{savaistu}[La fonction $\exp(-x^2)$ n'a pas de primitive « élémentaire »]
L'intégrale $\int e^{-x^2}\,dx$ ne peut pas s'exprimer avec les fonctions usuelles (polynômes, rationnelles, trigonométriques, exponentielles, logarithmes). C'est un résultat majeur (théorème de Liouville, 1835).
\par En statistiques, la loi normale (gaussienne) utilise la \emph{fonction d'erreur} :
\[
\operatorname{Erf}(x) = \frac{2}{\sqrt{\pi}}\int_0^x e^{-t^2}\,dt.
\]
Les tables de la loi normale ou les calculatrices donnent les valeurs de $\operatorname{Erf}$. L'intégrale de Gauss $\int_{-\infty}^{+\infty} e^{-x^2}\,dx = \sqrt{\pi}$ est un classique des concours.
\end{savaistu}

\begin{exercice}[Entraînement — Primitives et intégrales exponentielles]
\begin{enumerate}
\item Calculer les primitives sur $\mathbb{R}$ :
      \begin{enumerate}
      \item $\exp(5x-3)$
      \item $(3x^2-4x+2)\exp(x^3-2x^2+2x)$
      \item $x^2\exp(x)$
      \item $(2x+1)\exp(x^2+x)$
      \end{enumerate}
\item Calculer les intégrales définies :
      \begin{enumerate}
      \item $\int_0^{\ln 2} e^{2x}\,dx$
      \item $\int_1^2 (x+1)\exp(x^2+2x)\,dx$
      \item $\int_0^1 (x^2+1)\exp(x)\,dx$
      \end{enumerate}
\item Un système de refroidissement voit sa température $T(t)$ (en °C) évoluer selon $T(t)=20+80e^{-0,05t}$ ($t$ en minutes). Calculer la température moyenne sur les $30$ premières minutes.
\end{enumerate}
\end{exercice}

\begin{corrige}[Exercice n]
\begin{enumerate}
\item \begin{enumerate}
      \item $\frac{1}{5}\exp(5x-3)+C$
      \item $\exp(x^3-2x^2+2x)+C$
      \item Intégration par parties deux fois : $e^x(x^2-2x+2)+C$
      \item $\exp(x^2+x)+C$
      \end{enumerate}
\item \begin{enumerate}
      \item $\frac{1}{2}(e^{2\ln 2}-1)=\frac{1}{2}(4-1)=1,5$
      \item $u(x)=x^2+2x$, $u'(x)=2x+2=2(x+1)$. Donc $(x+1)dx = \frac{1}{2}du$.
            \[
            \int_1^2 (x+1)\exp(x^2+2x)\,dx = \frac{1}{2}\int_{u(1)}^{u(2)} e^u\,du = \frac{1}{2}\Bigl[e^u\Bigr]_3^8 = \frac{1}{2}(e^8-e^3).
            \]
      \item $\int_0^1 (x^2+1)\exp(x)\,dx = \int_0^1 x^2 e^x\,dx + \int_0^1 e^x\,dx = [e^x(x^2-2x+2)]_0^1 + [e^x]_0^1 = e(1-2+2)-2 + e-1 = 2e-3 \approx 2,436$
      \end{enumerate}
\item $m = \frac{1}{30}\int_0^{30}(20+80e^{-0,05t})\,dt = \frac{1}{30}\Bigl[20t - 1600e^{-0,05t}\Bigr]_0^{30} = \frac{1}{30}(600 - 1600e^{-1,5} + 1600) = \frac{2200 - 1600e^{-1,5}}{30} \approx 61,4$ °C.
\end{enumerate}
\end{corrige}

\begin{aretenir}[Résumé des primitives exponentielles — Fiche de révision Probatoire]
\begin{center}
\begin{tabularx}{\linewidth}{|l|X|}\hline
\textbf{Formule} & \textbf{Condition / Remarque} \\ \hline
$\int \exp(x)\,dx = \exp(x) + C$ & Base \\ \hline
$\int \exp(ax+b)\,dx = \frac{1}{a}\exp(ax+b)+C$ & $a\neq 0$ \\ \hline
$\int u'(x)\exp(u(x))\,dx = \exp(u(x))+C$ & Repérer $u$ et $u'$ \\ \hline
$\int_a^b u'(x)\exp(u(x))\,dx = \exp(u(b))-\exp(u(a))$ & Intégrale définie \\ \hline
$\int P(x)\exp(x)\,dx$ & Intégration par parties : dériver $P$, intégrer $\exp(x)$ \\ \hline
\end{tabularx}
\end{center}
\end{aretenir}

\coursec{Étude complète de fonctions faisant intervenir exp}

Dans la section précédente, nous avons appris à déterminer les primitives et les intégrales de fonctions faisant intervenir l'exponentielle. Nous disposons maintenant de tous les outils : propriétés algébriques de $\exp$, limites et croissances comparées, dérivation, primitives. Il est temps de rassembler ces savoirs pour mener à bien l'exercice le plus complet du Probatoire et du Baccalauréat technique : \cle{l'étude complète d'une fonction} définie à l'aide de $\exp$. C'est le type de problème qui mobilise le plus de points à l'examen : domaine, limites, dérivée, variations, asymptotes, convexité, tracé. Suivons une démarche systématique, étape par étape.

\courssub{La méthode générale : une checklist en 8 étapes}

L'erreur la plus fréquente des élèves n'est pas de ne pas savoir calculer, mais de ne pas savoir \emph{dans quel ordre} travailler. Voici la démarche professionnelle, à suivre dans l'ordre.

\begin{methode}[Étude complète d'une fonction faisant intervenir exp]
Pour étudier une fonction $f$ définie à l'aide de $\exp$, on suit systématiquement les 8 étapes :
\begin{enumerate}
\item \textbf{Domaine de définition} $\mathcal{D}_f$. Pour les fonctions de la forme $P(x)\mathrm{e}^{ax}$ avec $P$ polynôme, c'est presque toujours $\mathbb{R}$. Attention aux dénominateurs et aux $\ln$ éventuels.
\item \textbf{Limites aux bornes} de $\mathcal{D}_f$, en utilisant les \cle{croissances comparées} : $\lim\limits_{x\to+\infty}\dfrac{\mathrm{e}^{x}}{x}=+\infty$ et $\lim\limits_{x\to+\infty}x^{n}\mathrm{e}^{-x}=0$.
\item \textbf{Calcul de la dérivée} $f'$, puis \textbf{factorisation} : mettre $\mathrm{e}^{u(x)}$ en facteur. Comme $\mathrm{e}^{u(x)}>0$, le signe de $f'$ est le signe du facteur restant.
\item \textbf{Tableau de variations} complet (signe de $f'$, flèches, valeurs des extrema, limites aux bornes).
\item \textbf{Équation $f(x)=k$} : nombre de solutions déduit du tableau de variations (position de $k$ par rapport aux extrema et aux limites).
\item \textbf{Tableau de signes} de $f$ si nécessaire (position de la courbe par rapport à l'axe des abscisses).
\item \textbf{Points particuliers} : intersections avec les axes, tangentes remarquables (en particulier horizontales aux extrema).
\item \textbf{Tracé du graphique} : asymptotes, extrema, points d'intersection, allure générale. On valide avec la calculatrice ou GeoGebra.
\end{enumerate}
\end{methode}

\begin{attention}[Le réflexe qui simplifie tout : $\mathrm{e}^{x}>0$]
Pour tout réel $x$, $\mathrm{e}^{x}>0$. Conséquence capitale : après avoir mis $\mathrm{e}^{u(x)}$ en facteur dans $f'(x)$, \textbf{le signe de $f'(x)$ est exactement le signe du facteur polynomial ou rationnel restant}. Inutile d'étudier le signe de l'exponentielle ! Par exemple, si $f'(x)=(x-1)\mathrm{e}^{x}$, alors $f'(x)$ a le signe de $x-1$. Oublier ce réflexe conduit à des tableaux de signes inutilement compliqués.
\end{attention}

\courssub{Limites et asymptotes : que dire selon la forme de $f$ ?}

\textbf{Le cas $f(x)=P(x)\mathrm{e}^{ax}$}\par

C'est le cas le plus fréquent au baccalauréat. Tout dépend du signe de $a$.

\begin{propriete}[Limites de $P(x)\mathrm{e}^{ax}$]
Soit $P$ un polynôme non nul et $a>0$.
\begin{itemize}
\item Pour $f(x)=P(x)\mathrm{e}^{ax}$ : $\lim\limits_{x\to+\infty}f(x)=\pm\infty$ (selon le signe du coefficient dominant de $P$) ; la croissance est trop rapide pour une asymptote : \textbf{pas d'asymptote oblique en $+\infty$}. En $-\infty$, $\lim\limits_{x\to-\infty}P(x)\mathrm{e}^{ax}=0$ : la droite $y=0$ est \cle{asymptote horizontale}.
\item Pour $f(x)=P(x)\mathrm{e}^{-ax}$ : par croissances comparées, $\lim\limits_{x\to+\infty}P(x)\mathrm{e}^{-ax}=0$ : la droite $y=0$ (l'axe des abscisses) est \cle{asymptote horizontale} en $+\infty$.
\end{itemize}
\end{propriete}

\begin{attention}[Asymptote oblique : mythe et réalité]
Une asymptote oblique $y=ax+b$ exige $\lim\limits_{x\to\pm\infty}\dfrac{f(x)}{x}=a$ (fini, non nul) puis $\lim\limits_{x\to\pm\infty}\bigl(f(x)-ax\bigr)=b$ (fini). Pour $f(x)=P(x)\mathrm{e}^{x}$, $\dfrac{f(x)}{x}\to\pm\infty$ : \textbf{jamais d'asymptote oblique}, l'exponentielle « écrase » tout polynôme. Pour $f(x)=P(x)\mathrm{e}^{-x}$, la limite en $+\infty$ est $0$ : l'asymptote est \textbf{horizontale} ($y=0$), pas oblique. Retenir : avec l'exponentielle seule, on rencontre des asymptotes \emph{horizontales} ; les asymptotes obliques apparaissent plutôt pour des sommes du type $f(x)=ax+b+\mathrm{e}^{-x}$.
\end{attention}

\textbf{Autres formes classiques}\par

\begin{itemize}
\item $f(x)=\dfrac{\mathrm{e}^{u(x)}}{v(x)}$ : le domaine exclut les zéros de $v$ ; on cherche des \cle{asymptotes verticales} en ces valeurs (limite infinie).
\item $f(x)=\ln(u(x))+\mathrm{e}^{v(x)}$ : le domaine exige $u(x)>0$ ; les limites se traitent par somme, en identifiant le terme dominant.
\item $f(x)=ax+b+\mathrm{e}^{-x}$ : en $+\infty$, $\mathrm{e}^{-x}\to 0$, donc la droite $y=ax+b$ est asymptote oblique, et la courbe est \emph{au-dessus} de son asymptote car $\mathrm{e}^{-x}>0$.
\end{itemize}

\courssub{Exemple chiffré complet : étude de $f(x)=(x^{2}-4)\mathrm{e}^{x}$}

\begin{exemplebox}[Étude complète de $f(x)=(x^{2}-4)\mathrm{e}^{x}$ sur $\mathbb{R}$]
Appliquons la checklist des 8 étapes.

\textbf{Étape 1 — Domaine.} $x^{2}-4$ et $\mathrm{e}^{x}$ sont définis sur $\mathbb{R}$ : $\mathcal{D}_f=\mathbb{R}$.

\textbf{Étape 2 — Limites.} En $-\infty$ : $\mathrm{e}^{x}\to 0$ et, par croissances comparées, $(x^{2}-4)\mathrm{e}^{x}\to 0$. Donc $\lim\limits_{x\to-\infty}f(x)=0$ : asymptote horizontale $y=0$. En $+\infty$ : $x^{2}-4\to+\infty$ et $\mathrm{e}^{x}\to+\infty$, donc $\lim\limits_{x\to+\infty}f(x)=+\infty$.

\textbf{Étape 3 — Dérivée.} Avec $(uv)'=u'v+uv'$ :
\[ f'(x)=2x\,\mathrm{e}^{x}+(x^{2}-4)\mathrm{e}^{x}=(x^{2}+2x-4)\mathrm{e}^{x}. \]
Comme $\mathrm{e}^{x}>0$, $f'(x)$ a le signe de $x^{2}+2x-4$. Discriminant : $\Delta=4+16=20$, racines $x_{1}=-1-\sqrt{5}\approx -3{,}24$ et $x_{2}=-1+\sqrt{5}\approx 1{,}24$. Le trinôme est positif à l'extérieur des racines, négatif entre les racines.

\textbf{Étape 4 — Tableau de variations.}
\begin{center}
\begin{tabular}{|c|ccccccc|}\hline
$x$ & $-\infty$ & & $-1-\sqrt{5}$ & & $-1+\sqrt{5}$ & & $+\infty$ \\ \hline
$f'(x)$ & & $+$ & $0$ & $-$ & $0$ & $+$ & \\ \hline
$f$ & $0$ & $\nearrow$ & $f(x_{1})\approx 0{,}25$ & $\searrow$ & $f(x_{2})\approx -8{,}5$ & $\nearrow$ & $+\infty$ \\ \hline
\end{tabular}
\end{center}
Valeurs approchées : $f(x_{1})=(x_{1}^{2}-4)\mathrm{e}^{x_{1}}\approx 0{,}25$ (maximum local) et $f(x_{2})\approx -8{,}5$ (minimum local).

\textbf{Étape 5 — Équation $f(x)=0$.} $(x^{2}-4)\mathrm{e}^{x}=0 \iff x^{2}-4=0 \iff x=-2$ ou $x=2$ (car $\mathrm{e}^{x}\neq 0$).

\textbf{Étape 6 — Signe de $f$.} $f(x)$ a le signe de $x^{2}-4$ : positif sur $]-\infty;-2[\cup]2;+\infty[$, négatif sur $]-2;2[$.

\textbf{Étape 7 — Points particuliers.} Intersection avec l'axe des ordonnées : $f(0)=-4$. Intersections avec l'axe des abscisses : $(-2;0)$ et $(2;0)$. Tangentes horizontales en $x_{1}$ et $x_{2}$.

\textbf{Étape 8 — Tracé.} Voir la figure ci-dessous, validée sur GeoGebra.
\end{exemplebox}

\begin{popfigure}
\begin{center}
\begin{tikzpicture}
\begin{axis}[width=0.85\linewidth,height=7cm,axis lines=middle,xlabel={$x$},ylabel={$y$},xmin=-5.5,xmax=2.8,ymin=-9,ymax=1.5,xtick={-4,-2,2},ytick={-8,-4,0.25},grid=both,grid style={popline!40},samples=200,clip=true]
\addplot[popcurve,domain=-5.5:2.05]{(x^2-4)*exp(x)};
\addplot[poporange,dashed,domain=-5.5:2.8]{0};
\addplot[mark=*,popdark,only marks] coordinates {(-2,0) (2,0) (0,-4) (-3.236,0.254) (1.236,-8.51)};
\node[popblue,anchor=south] at (axis cs:-3.2,0.3) {\footnotesize max};
\node[popblue,anchor=north] at (axis cs:1.3,-8.8) {\footnotesize min};
\node[poporange,anchor=south west] at (axis cs:-5.4,0.3) {\footnotesize asymptote $y=0$};
\end{axis}
\end{tikzpicture}
\end{center}
\legende{Courbe de $f(x)=(x^{2}-4)\mathrm{e}^{x}$ : asymptote horizontale en $-\infty$, maximum local en $x\approx-3{,}24$ ($y\approx0{,}25$), minimum local en $x\approx1{,}24$ ($y\approx-8{,}5$), intersections avec les axes.}
\end{popfigure}

\courssub{Convexité et points d'inflexion}

L'allure d'une courbe (creuse vers le haut ou vers le bas) se lit sur la dérivée seconde.

\begin{definition}[Convexité]
Soit $f$ deux fois dérivable sur un intervalle $I$.
\begin{itemize}
\item Si $f''(x)\geq 0$ sur $I$, $f$ est \cle{convexe} sur $I$ : sa courbe est « creuse vers le haut » (elle est au-dessus de ses tangentes).
\item Si $f''(x)\leq 0$ sur $I$, $f$ est \cle{concave} : sa courbe est « creuse vers le bas ».
\item Un \cle{point d'inflexion} est un point où $f''$ s'annule \emph{en changeant de signe} : la courbe y traverse sa tangente.
\end{itemize}
\end{definition}

\begin{propriete}[Dérivée seconde d'une fonction exponentielle composée]
Si $f(x)=\mathrm{e}^{u(x)}$ avec $u$ deux fois dérivable, alors
\[ f'(x)=u'(x)\mathrm{e}^{u(x)} \qquad\text{et}\qquad f''(x)=\bigl(u''(x)+[u'(x)]^{2}\bigr)\mathrm{e}^{u(x)}. \]
Comme $\mathrm{e}^{u(x)}>0$, le signe de $f''$ est le signe de $u''+(u')^{2}$.
\end{propriete}

\begin{exemplebox}[Convexité de $g(x)=x\mathrm{e}^{-x}$]
$g'(x)=(1-x)\mathrm{e}^{-x}$, puis $g''(x)=-\mathrm{e}^{-x}+(1-x)(-\mathrm{e}^{-x})=(x-2)\mathrm{e}^{-x}$. Comme $\mathrm{e}^{-x}>0$, $g''(x)$ a le signe de $x-2$ : $g$ est concave sur $]-\infty;2]$, convexe sur $[2;+\infty[$, et le point d'abscisse $x=2$, de coordonnées $\left(2\,;\,\dfrac{2}{\mathrm{e}^{2}}\right)$, est un \cle{point d'inflexion}.
\end{exemplebox}

\begin{popfigure}
\begin{center}
\begin{tikzpicture}
\begin{axis}[width=0.85\linewidth,height=6.5cm,axis lines=middle,xlabel={$x$},ylabel={$y$},xmin=-1.5,xmax=7,ymin=-0.35,ymax=0.5,xtick={1,2},ytick={0.368},yticklabels={$1/\mathrm{e}$},grid=both,grid style={popline!40},samples=200]
\addplot[popcurve,domain=-1.4:7]{x*exp(-x)};
\addplot[poporange,dashed,domain=-1.5:7]{0};
\addplot[mark=*,popdark,only marks] coordinates {(1,0.3679) (2,0.2707) (0,0)};
\node[popblue,anchor=south] at (axis cs:1,0.38) {\footnotesize max $(1\,;\,1/\mathrm{e})$};
\node[poppurple,anchor=west] at (axis cs:2.1,0.27) {\footnotesize inflexion};
\addplot[poppurple,dashed,domain=0.5:3.5]{0.5412 - 0.1353*x};
\node[popdark,anchor=south west] at (axis cs:-1.4,0.02) {\footnotesize asymptote $y=0$};
\end{axis}
\end{tikzpicture}
\end{center}
\legende{Courbe de $g(x)=x\mathrm{e}^{-x}$ : maximum en $x=1$, point d'inflexion en $x=2$ (avec sa tangente), asymptote horizontale $y=0$ en $+\infty$.}
\end{popfigure}

\courssub{Optimisation : maximiser un profit, minimiser un coût}

De nombreuses situations techniques se modélisent par $f(x)=P(x)\mathrm{e}^{-kx}$ : rendement d'une machine en fonction de sa vitesse, efficacité d'un médicament en fonction du temps, profit en fonction du prix. Chercher l'optimum revient à chercher l'extremum de $f$.

\begin{exoresolu}[Optimisation d'un profit en atelier]
\textbf{Énoncé.} Un atelier de Douala produit des pièces mécaniques. Le bénéfice journalier, en milliers de FCFA, est modélisé par $B(x)=(10x)\mathrm{e}^{-0{,}5x}$, où $x$ est le nombre de dizaines de pièces produites par jour ($x\geq 0$). Déterminer la production qui maximise le bénéfice et la valeur de ce bénéfice maximal.
\tcblower
\textbf{Solution détaillée.}
\begin{enumerate}
\item \textbf{Dérivée} : $B'(x)=10\,\mathrm{e}^{-0{,}5x}+10x(-0{,}5)\mathrm{e}^{-0{,}5x}=(10-5x)\mathrm{e}^{-0{,}5x}=5(2-x)\mathrm{e}^{-0{,}5x}$.
\item \textbf{Signe} : comme $\mathrm{e}^{-0{,}5x}>0$, $B'(x)$ a le signe de $2-x$. Donc $B$ est croissante sur $[0;2]$ et décroissante sur $[2;+\infty[$.
\item \textbf{Maximum} : atteint en $x=2$, c'est-à-dire pour $20$ pièces par jour. Valeur : $B(2)=20\,\mathrm{e}^{-1}=\dfrac{20}{\mathrm{e}}\approx 7{,}36$ milliers de FCFA, soit environ \num{7358} FCFA.
\item \textbf{Conclusion} : l'atelier maximise son bénéfice en produisant $20$ pièces par jour ; au-delà, les coûts de surproduction dégradent le profit.
\end{enumerate}
\end{exoresolu}

\begin{methode}[Résoudre un problème d'optimisation avec exp]
\begin{enumerate}
\item Identifier la grandeur à optimiser et la variable ; écrire la fonction $f$ (souvent donnée dans l'énoncé).
\item Préciser l'intervalle d'étude (contraintes concrètes : $x\geq 0$, capacité maximale...).
\item Calculer $f'$, factoriser par l'exponentielle, étudier le signe du facteur restant.
\item Conclure par le tableau de variations, puis \textbf{rédiger la réponse dans le langage de la situation} (unités, quantités concrètes).
\end{enumerate}
\end{methode}

\courssub{Validation numérique et erreurs fréquentes}

Avant de rendre sa copie, on \textbf{valide} son étude : sur calculatrice (mode TABLE, on vérifie quelques valeurs de $f$ et le sens de variation) ou sur GeoGebra (on trace la courbe et on compare extrema, asymptotes, intersections). Un écart entre le tableau de variations et la courbe tracée signale une erreur de dérivation ou de limite.

\begin{attention}[Les quatre pièges classiques]
\begin{itemize}
\item \textbf{Dérivation} : $\bigl(\mathrm{e}^{u}\bigr)'=u'\mathrm{e}^{u}$, et non $\mathrm{e}^{u}$ ! Pour $f(x)=\mathrm{e}^{-2x}$, $f'(x)=-2\mathrm{e}^{-2x}$.
\item \textbf{Signe} : ne jamais écrire « $\mathrm{e}^{x}$ est négatif pour $x<0$ » : $\mathrm{e}^{x}>0$ pour \emph{tout} réel $x$.
\item \textbf{Limites} : en $-\infty$, $P(x)\mathrm{e}^{x}\to 0$ (l'exponentielle l'emporte) ; c'est une forme indéterminée « $\infty\times 0$ » qu'on lève par croissances comparées, jamais par « $0\times\infty=0$ » sans justification.
\item \textbf{Équation} $f(x)=k$ : le nombre de solutions se justifie par le tableau de variations (théorème des valeurs intermédiaires sur chaque intervalle de monotonie), pas par un simple dessin.
\end{itemize}
\end{attention}

\begin{exercice}[Pour s'entraîner]
On considère $h(x)=(x-2)\mathrm{e}^{x}$ sur $\mathbb{R}$.
\begin{enumerate}
\item Calculer les limites de $h$ en $-\infty$ et $+\infty$. Préciser l'asymptote éventuelle.
\item Calculer $h'(x)$, dresser le tableau de variations de $h$.
\item Déterminer les intersections de la courbe avec les axes et le signe de $h$.
\item Calculer $h''(x)$ et en déduire la convexité et le point d'inflexion éventuel.
\item Tracer la courbe sur GeoGebra et vérifier les résultats.
\end{enumerate}
\end{exercice}

\begin{corrige}[Exercice 5.6]
\begin{enumerate}
\item En $-\infty$ : par croissances comparées, $(x-2)\mathrm{e}^{x}\to 0$ : asymptote horizontale $y=0$. En $+\infty$ : $h(x)\to+\infty$.
\item $h'(x)=\mathrm{e}^{x}+(x-2)\mathrm{e}^{x}=(x-1)\mathrm{e}^{x}$. Signe de $x-1$ : $h$ décroît sur $]-\infty;1]$, croît sur $[1;+\infty[$ ; minimum $h(1)=-\mathrm{e}\approx -2{,}72$.
\item $h(x)=0\iff x=2$ ; $h(0)=-2$. $h$ est négative sur $]-\infty;2[$, positive sur $]2;+\infty[$.
\item $h''(x)=\mathrm{e}^{x}+(x-1)\mathrm{e}^{x}=x\,\mathrm{e}^{x}$ : signe de $x$. Concave sur $\mathbb{R}^{-}$, convexe sur $\mathbb{R}^{+}$, point d'inflexion en $(0;-2)$.
\item La courbe GeoGebra confirme : minimum en $(1;-\mathrm{e})$, passage par $(0;-2)$ et $(2;0)$, asymptote $y=0$ en $-\infty$.
\end{enumerate}
\end{corrige}

\begin{savaistu}[La courbe $y=x\mathrm{e}^{-x}$ et la loi Gamma]
La courbe de $g(x)=x\mathrm{e}^{-x}$ que nous avons étudiée n'est pas qu'un exercice scolaire : à un coefficient près, c'est la \cle{densité de la loi Gamma} de paramètre $2$. En théorie des probabilités, elle modélise le \textbf{temps d'attente avant la survenue du 2\textsuperscript{ème} événement} dans un processus de Poisson : par exemple le temps avant la deuxième panne d'une machine, le deuxième appel à un standard, le deuxième client à un guichet. Le maximum en $x=1$ correspond au temps d'attente le plus « probable ». Les ingénieurs télécoms et les automaticiens l'utilisent quotidiennement !
\end{savaistu}

En résumé, l'étude complète d'une fonction avec exponentielle n'est pas une suite de calculs isolés mais une \textbf{enquête méthodique} : domaine, limites, dérivée factorisée, variations, convexité, points remarquables, tracé validé numériquement. Cette démarche, rédigée avec soin, est exactement ce que les correcteurs du Baccalauréat technique attendent. La section suivante montrera comment ces mêmes fonctions modélisent les phénomènes de croissance et de décroissance continues de la vie courante.

\coursec{Modélisation concrète : croissance et décroissance continues}

Après avoir maîtrisé l'étude complète de fonctions faisant intervenir l'exponentielle (section 5), nous disposons désormais de tous les outils — limites, dérivées, primitives, variations — pour comprendre \emph{ pourquoi } la fonction $\exp$ s'impose comme le modèle universel des phénomènes où la vitesse de variation d'une grandeur est proportionnelle à cette grandeur elle-même. C'est le cœur de la modélisation en croissance et décroissance continues, omniprésente en biologie, finance, physique, chimie et médecine.

\courssub{Le modèle universel : l'équation différentielle $y' = ky$}

Toute la puissance de la modélisation exponentielle repose sur une idée simple : \emph{ le taux de variation relative est constant }.

\begin{definition}[Taux de variation relative instantanée]
Soit $y(t)$ une fonction dérivable et strictement positive sur un intervalle $I$. Le \cle{taux de variation relative instantanée} de $y$ en $t$ est le quotient $\dfrac{y'(t)}{y(t)}$.
\end{definition}

\begin{propriete}[Caractérisation de l'exponentielle]
Une fonction $y$ dérivable sur $\mathbb{R}$ (ou un intervalle) satisfait $y'(t) = k \, y(t)$ avec $k \in \mathbb{R}$ constant \textbf{si et seulement si} il existe $y_0 \in \mathbb{R}$ tel que $y(t) = y_0 \, e^{kt}$ pour tout $t$.
\end{propriete}

\begin{experience}[Démonstration guidée : de l'équation différentielle à la solution]
\textbf{Étape 1 :} On cherche $y$ telle que $y' = ky$. Si $y$ ne s'annule pas, on peut diviser : $\dfrac{y'}{y} = k$.
\textbf{Étape 2 :} On reconnaît la dérivée de $\ln|y|$ : $(\ln|y|)' = k$.
\textbf{Étape 3 :} On intègre : $\ln|y(t)| = kt + C$.
\textbf{Étape 4 :} On exponentie : $|y(t)| = e^{kt+C} = e^C e^{kt}$. En posant $y_0 = \pm e^C$, on obtient $y(t) = y_0 e^{kt}$.
\textbf{Étape 5 :} Réciproquement, si $y(t) = y_0 e^{kt}$, alors $y'(t) = y_0 k e^{kt} = k y(t)$. \hfill $\square$
\end{experience}

\begin{propriete}[Modèle exponentiel général]
\begin{itemize}
    \item \textbf{Équation différentielle :} $y'(t) = k \, y(t)$.
    \item \textbf{Solution :} $y(t) = y_0 \, e^{kt}$ où $y_0 = y(0)$ est la valeur initiale.
    \item \textbf{Paramètre $k$ (taux instantané) :}
    \begin{itemize}
        \item $k > 0$ : \cle{croissance exponentielle} (la quantité augmente d'autant plus vite qu'elle est grande).
        \item $k < 0$ : \cle{décroissance exponentielle} (la quantité diminue d'autant plus vite qu'elle est grande).
        \item $k = 0$ : quantité constante.
    \end{itemize}
\end{itemize}
\end{propriete}

\begin{popfigure}
\begin{tikzpicture}[scale=0.9, every node/.style={font=\small}]
\begin{axis}[
    width=\linewidth, height=7cm,
    axis lines=middle,
    xlabel=$t$, ylabel=$y(t)$,
    xmin=0, xmax=5, ymin=0, ymax=5,
    xtick={0,1,2,3,4,5}, ytick={0,1,2,3,4,5},
    domain=0:5, samples=100,
    legend style={at={(0.02,0.98)}, anchor=north west, font=\tiny},
    grid=major, grid style={popline!50}
]
\addplot[popblue, thick, domain=0:5] {exp(0.5*x)}; \addlegendentry{$k=0,5 > 0$ (Croissance)}
\addplot[popgreen, thick, domain=0:5] {exp(-0.5*x)}; \addlegendentry{$k=-0,5 < 0$ (Décroissance)}
\addplot[poporange, thick, dashed, domain=0:5] {1}; \addlegendentry{$k=0$ (Constant)}
\addplot[only marks, mark=*, popdark] coordinates {(0,1)};
\node[above right, popdark] at (axis cs:0,1) {$y_0=1$};
\end{axis}
\end{tikzpicture}
\legende{Évolution de $y(t)=y_0 e^{kt}$ pour $y_0=1$ selon le signe de $k$. La courbe passe toujours par $(0,y_0)$.}
\end{popfigure}

\courssub{Croissance démographique et intérêts composés continus}

C'est l'application historique et financière par excellence.

\begin{propriete}[Capitalisation continue]
Si un capital $C_0$ est placé à un \cle{taux instantané} $k$ (force d'intérêt, exprimé en décimal) avec capitalisation \textbf{continue}, le capital au bout de $t$ années est :
\[
C(t) = C_0 \, e^{kt}.
\]
Le taux $k$ est ici le \textbf{taux de variation relative instantanée} : $\dfrac{C'(t)}{C(t)} = k$.
Si l'on connaît le \textbf{taux annuel effectif} $r$ (capitalisation annuelle), le taux instantané équivalent est $k = \ln(1+r)$.
\end{propriete}

\begin{attention}[Continu vs Discret : ne pas confondre $r$ et $k$ !]
\begin{itemize}
    \item \textbf{Capitalisation discrète (annuelle) :} $C_{\text{disc}}(t) = C_0 (1+r)^t$. Le taux $r$ est un \emph{taux périodique effectif}.
    \item \textbf{Capitalisation continue :} $C_{\text{cont}}(t) = C_0 e^{kt}$. Le taux $k$ est un \emph{taux instantané} (force d'intérêt).
    \item \textbf{Lien exact :} Pour produire le même résultat après 1 an, il faut $e^k = 1+r \iff k = \ln(1+r)$.
    \item \textbf{Approximation :} Pour $r$ petit, $\ln(1+r) \approx r - \frac{r^2}{2}$, donc $k \approx r$. L'erreur relative est de l'ordre de $r/2$.
\end{itemize}
\end{attention}

\begin{exoresolu}[Placement à 5\% : continu vs annuel]
\textbf{Énoncé :} On place $10\,000$ FCFA à un taux annuel effectif de $5\%$ ($r=0,05$).
\begin{enumerate}
    \item Calculer le capital après 10 ans en capitalisation annuelle.
    \item Calculer le capital après 10 ans en capitalisation continue (taux instantané $k = \ln(1,05)$).
    \item Comparer et déterminer le temps de doublement dans les deux cas.
\end{enumerate}
\tcblower
\textbf{Solution détaillée :}
\begin{enumerate}
    \item \textbf{Annuel :} $C_{10} = 10000 \times (1,05)^{10} \approx 10000 \times 1,62889 = \textbf{16 289 FCFA}$.
    \item \textbf{Continu :} $k = \ln(1,05) \approx 0,04879$.
    $C(10) = 10000 \times e^{0,04879 \times 10} = 10000 \times e^{0,4879} \approx 10000 \times 1,62889 = \textbf{16 289 FCFA}$.
    \emph{Remarque :} Les résultats sont identiques car $k$ a été choisi pour équivaloir à $r=5\%$ annuel.
    \item \textbf{Temps de doublement $T_2$ :} On cherche $t$ tel que $C(t) = 2C_0$.
    \begin{itemize}
        \item Annuel : $(1,05)^{T_2} = 2 \Rightarrow T_2 \ln(1,05) = \ln 2 \Rightarrow T_2 = \dfrac{\ln 2}{\ln 1,05} \approx \dfrac{0,6931}{0,04879} \approx \textbf{14,21 ans}$.
        \item Continu : $e^{k T_2} = 2 \Rightarrow k T_2 = \ln 2 \Rightarrow T_2 = \dfrac{\ln 2}{k} = \dfrac{\ln 2}{\ln 1,05} \approx \textbf{14,21 ans}$.
    \end{itemize}
    \textbf{Conclusion :} Si on définit correctement $k = \ln(1+r)$, les deux modèles sont \emph{mathématiquement équivalents}. La capitalisation continue est juste l'écriture « différentielle » du même phénomène.
\end{enumerate}
\end{exoresolu}

\begin{popfigure}
\begin{tikzpicture}[scale=0.85]
\begin{axis}[
    width=\linewidth, height=7cm,
    axis lines=middle,
    xlabel=$t$ (années), ylabel=$C(t)$ (FCFA),
    xmin=0, xmax=20, ymin=0, ymax=35000,
    xtick={0,5,10,14.21,15,20}, ytick={10000,20000,30000},
    domain=0:20, samples=100,
    legend style={at={(0.02,0.98)}, anchor=north west, font=\tiny},
    grid=major, grid style={popline!50}
]
% Continu (lisse)
\addplot[popblue, very thick] {10000*exp(ln(1.05)*x)}; \addlegendentry{Continu $C_0 e^{kt}$}
% Discret (points + traits d'escalier) - simulation via const plot
\addplot[poporange, thick, mark=*, only marks, samples at={0,1,...,20}] {10000*1.05^x}; \addlegendentry{Annuel $C_0(1+r)^t$ (points)}
\addplot[poporange, thick, const plot, samples at={0,1,...,20}] {10000*1.05^x};
% Ligne doublement
\addplot[popdark, dashed, domain=0:20] {20000};
\draw[popdark, dashed] (axis cs:14.21,0) -- (axis cs:14.21,20000);
\node[popdark, anchor=north, font=\tiny] at (axis cs:14.21,0) {$T_2 \approx 14,2$};
\node[popdark, anchor=east, font=\tiny] at (axis cs:0,20000) {$2C_0$};
\end{axis}
\end{tikzpicture}
\legende{Comparaison capitalisation continue (courbe lisse) vs annuelle (escalier). Le temps de doublement $T_2$ est identique si $k=\ln(1+r)$.}
\end{popfigure}

\courssub{Décroissance radioactive, chimie d'ordre 1 et datation au Carbone 14}

La loi de décroissance radioactive est l'archétype de la décroissance exponentielle : le nombre de noyaux qui se désintègrent par unité de temps est proportionnel au nombre de noyaux présents.

\begin{propriete}[Loi de décroissance radioactive / Réaction d'ordre 1]
Soit $N(t)$ le nombre de noyaux (ou la concentration d'un réactif) à l'instant $t$. La loi s'écrit :
\[
N(t) = N_0 \, e^{-\lambda t} \quad (\lambda > 0 \text{ constante de décroissance}).
\]
La \cle{demi-vie} (ou période radioactive) $T_{1/2}$ est le temps nécessaire pour que la quantité soit divisée par 2 :
\[
N(T_{1/2}) = \frac{N_0}{2} \iff e^{-\lambda T_{1/2}} = \frac{1}{2} \iff \lambda T_{1/2} = \ln 2 \iff \boxed{T_{1/2} = \frac{\ln 2}{\lambda}}.
\]
\end{propriete}

\begin{exemplebox}[Datation au Carbone 14 ($^{14}\text{C}$)]
Le $^{14}\text{C}$ est un isotope radioactif du carbone présent dans l'atmosphère (formation par rayons cosmiques). Un organisme vivant maintient un rapport $^{14}\text{C}/^{12}\text{C}$ constant. À sa mort, les échanges s'arrêtent et le $^{14}\text{C}$ décroît.
\begin{itemize}
    \item Constante de décroissance : $\lambda \approx 1,21 \times 10^{-4} \text{ an}^{-1}$.
    \item Demi-vie : $T_{1/2} = \dfrac{\ln 2}{1,21 \times 10^{-4}} \approx \textbf{5 730 ans}$.
    \item Mesure : On trouve l'activité actuelle $A$ (ou le rapport $R$) et on compare à l'activité de référence $A_0$ (organisme vivant).
    \item Âge $t$ : $\dfrac{A}{A_0} = e^{-\lambda t} \Rightarrow t = -\dfrac{1}{\lambda} \ln\left(\dfrac{A}{A_0}\right)$.
\end{itemize}
\end{exemplebox}

\begin{exoresolu}[Datation d'un échantillon]
\textbf{Énoncé :} Un fragment de bois trouvé sur un site archéologique présente une activité $^{14}\text{C}$ égale à $12,5\%$ de celle d'un bois actuel. Quelle est son âge approximatif ?
\tcblower
\textbf{Solution :}
$\dfrac{A}{A_0} = 0,125 = \dfrac{1}{8} = \left(\dfrac{1}{2}\right)^3$.
Cela correspond à \textbf{3 demi-vies}.
$t = 3 \times T_{1/2} = 3 \times 5730 \approx \textbf{17 190 ans}$.
\textit{Vérification par la formule :} $t = -\dfrac{1}{1,21\times 10^{-4}} \ln(0,125) \approx 8264 \times 2,079 \approx 17 180 \text{ ans}$.
\end{exoresolu}

\begin{popfigure}
\begin{tikzpicture}[scale=0.9]
\begin{axis}[
    width=\linewidth, height=7cm,
    axis lines=middle,
    xlabel=$t$ (années $\times 10^3$), ylabel=$N(t)/N_0$,
    xmin=0, xmax=20, ymin=0, ymax=1.1,
    xtick={0,5.73,11.46,17.19}, xticklabels={$0$, $T_{1/2}$, $2T_{1/2}$, $3T_{1/2}$},
    ytick={0,0.5,0.25,0.125,1}, yticklabels={$0$, $1/2$, $1/4$, $1/8$, $1$},
    domain=0:20, samples=200,
    grid=major, grid style={popline!50},
    legend style={at={(0.98,0.98)}, anchor=north east, font=\tiny}
]
\addplot[poppurple, very thick] {exp(-ln(2)/5.73*x)};
\addlegendentry{$N(t)=N_0 e^{-\lambda t}$}
% Marquer les demi-vies
\draw[poppurple, dashed] (axis cs:5.73,0) -- (axis cs:5.73,0.5) -- (axis cs:0,0.5);
\fill[poppurple] (axis cs:5.73,0.5) circle (2pt);
\draw[poppurple, dashed] (axis cs:11.46,0) -- (axis cs:11.46,0.25) -- (axis cs:0,0.25);
\fill[poppurple] (axis cs:11.46,0.25) circle (2pt);
\draw[poppurple, dashed] (axis cs:17.19,0) -- (axis cs:17.19,0.125) -- (axis cs:0,0.125);
\fill[poppurple] (axis cs:17.19,0.125) circle (2pt);
\node[poppurple, anchor=north west, font=\tiny] at (axis cs:5.73,0.5) {$T_{1/2}$};
\node[poppurple, anchor=north west, font=\tiny] at (axis cs:11.46,0.25) {$2T_{1/2}$};
\node[poppurple, anchor=north west, font=\tiny] at (axis cs:17.19,0.125) {$3T_{1/2} \approx 17\,190$ ans};
\end{axis}
\end{tikzpicture}
\legende{Courbe de décroissance radioactive. Les demi-vies successives divisent la quantité par 2 à intervalles réguliers.}
\end{popfigure}

\courssub{Pharmacocinétique : élimination d'un médicament (modèle à un compartiment)}

Après une injection intraveineuse (bolus IV), le médicament se distribue quasi-instantanément dans un « volume de distribution » $V$ et est éliminé (rein, foie) avec une cinétique d'ordre 1.

\begin{propriete}[Modèle à un compartiment - Bolus IV]
\begin{itemize}
    \item $D$ : dose administrée (mg).
    \item $V$ : volume de distribution (L). Concentration initiale $C_0 = D/V$.
    \item $k$ : constante d'élimination (h$^{-1}$). Taux de variation relative de la concentration.
    \item Concentration plasmatique : $\boxed{C(t) = \frac{D}{V} \, e^{-kt}}$.
    \item Demi-vie d'élimination : $t_{1/2} = \dfrac{\ln 2}{k}$.
    \item Clairance (Clearance) : $Cl = k \cdot V$ (volume de sang épuré par unité de temps).
\end{itemize}
\end{propriete}

\begin{popfigure}
\begin{tikzpicture}[scale=0.9]
\begin{axis}[
    width=\linewidth, height=6cm,
    axis lines=middle,
    xlabel=$t$ (h), ylabel=$C(t)$ (mg/L),
    xmin=0, xmax=8, ymin=0, ymax=28,
    xtick={0,1,2,4,6.9}, ytick={0,2,6,18,25},
    domain=0:8, samples=200,
    grid=major, grid style={popline!50},
    legend style={at={(0.98,0.98)}, anchor=north east, font=\tiny}
]
\addplot[poppurple, very thick] {25*exp(-0.366*x)};
\addlegendentry{$C(t)=25 e^{-0,366 t}$}
% Demi-vie
\draw[poppurple, dashed] (axis cs:1.89,0) -- (axis cs:1.89,12.5) -- (axis cs:0,12.5);
\node[poppurple, anchor=south east, font=\tiny] at (axis cs:1.89,12.5) {$t_{1/2}$};
% Seuil
\draw[poporange, dashed] (axis cs:0,2) -- (axis cs:6.9,2);
\draw[poporange, dashed] (axis cs:6.9,0) -- (axis cs:6.9,2);
\node[poporange, anchor=north west, font=\tiny] at (axis cs:6.9,2) {$C_{\min}=2$};
\fill[poppurple] (axis cs:0,25) circle (2pt) node[above left, font=\tiny] {$C_0=25$};
\end{axis}
\end{tikzpicture}
\legende{Élimination du médicament : concentration plasmatique, demi-vie et seuil thérapeutique.}
\end{popfigure}

\begin{methode}[Résoudre un problème de pharmacocinétique]
\begin{enumerate}
    \item Identifier les données : $D$, $V$ (ou $C_0$), $k$ (ou $t_{1/2}$, ou deux mesures $C(t_1), C(t_2)$).
    \item Écrire la loi : $C(t) = C_0 e^{-kt}$.
    \item \textbf{Si $k$ inconnu :} utiliser deux points $\Rightarrow k = \dfrac{1}{t_2-t_1} \ln\left(\dfrac{C(t_1)}{C(t_2)}\right)$.
    \item Calculer la quantité demandée : $C(t)$, $t$ pour un seuil $C_{\text{seuil}}$, $t_{1/2}$, $Cl$, AUC (Aire Under Curve = $C_0/k$).
\end{enumerate}
\end{methode}

\begin{exoresolu}[Détermination de $k$ et temps sous seuil thérapeutique]
\textbf{Énoncé :} Injection IV de $500$ mg d'un antibiotique ($V = 20$ L). Mesures : $C(1\text{h}) = 18$ mg/L, $C(4\text{h}) = 6$ mg/L.
\begin{enumerate}
    \item Déterminer $k$ et $t_{1/2}$.
    \item Vérifier $C_0$. Combien de temps la concentration reste-t-elle au-dessus du seuil thérapeutique $C_{\text{min}} = 2$ mg/L ?
\end{enumerate}
\tcblower
\textbf{Solution :}
\begin{enumerate}
    \item $k = \dfrac{1}{4-1} \ln\left(\dfrac{18}{6}\right) = \dfrac{\ln 3}{3} \approx \dfrac{1,0986}{3} \approx \textbf{0,366 h}\textsuperscript{$-1$}$.
    $t_{1/2} = \dfrac{\ln 2}{k} \approx \dfrac{0,693}{0,366} \approx \textbf{1,89 h}$.
    \item $C_0 = \dfrac{D}{V} = \dfrac{500}{20} = 25$ mg/L.
    Vérif : $C(1) = 25 e^{-0,366} \approx 25 \times 0,693 \approx 17,3 \approx 18$ (cohérence).
    Seuil : $25 e^{-kt} = 2 \Rightarrow e^{-kt} = 0,08 \Rightarrow -kt = \ln 0,08 \approx -2,526$.
    $t = \dfrac{2,526}{0,366} \approx \textbf{6,9 h}$.
\end{enumerate}
\end{exoresolu}

\courssub{Refroidissement de Newton (Loi de refroidissement)}

La température d'un corps chaud placé dans un milieu à température constante évolue exponentiellement vers la température du milieu.

\begin{propriete}[Loi de Newton]
Soit $T(t)$ la température d'un corps, $T_a$ la température ambiante (constante), $T_0 = T(0)$ la température initiale. La vitesse de refroidissement est proportionnelle à l'écart de température :
\[
\frac{dT}{dt} = -k (T - T_a), \quad k > 0.
\]
En posant $y(t) = T(t) - T_a$ (l'écart), on a $y' = -ky \Rightarrow y(t) = y_0 e^{-kt}$.
\[
\boxed{T(t) = T_a + (T_0 - T_a) e^{-kt}}
\]
\end{propriete}

\begin{attention}[Piège classique : la variable est l'écart $T-T_a$, pas $T$]
Ne JAMAIS écrire $T(t) = T_0 e^{-kt}$ pour le refroidissement (sauf si $T_a = 0^\circ\text{C}$). La température ne tend pas vers 0 mais vers $T_a$. C'est l'écart $T-T_a$ qui décroît exponentiellement vers 0.
\end{attention}

\begin{popfigure}
\begin{tikzpicture}[scale=0.9]
\begin{axis}[
    width=\linewidth, height=6cm,
    axis lines=middle,
    xlabel=$t$ (h), ylabel=$T(t)$ ($^\circ$C),
    xmin=0, xmax=5, ymin=15, ymax=40,
    xtick={0,1,2,2.62,3,4,5}, ytick={20,26,28,37},
    xticklabels={$0$, $1$, $2$, $t_1$, $3$, $4$, $5$},
    domain=0:5, samples=200,
    grid=major, grid style={popline!50},
    legend style={at={(0.98,0.98)}, anchor=north east, font=\tiny}
]
\addplot[popblue, very thick] {20 + 17*exp(-0.288*x)};
\addlegendentry{$T(t)=20+17 e^{-0,288 t}$}
\addplot[poporange, dashed, domain=0:5] {20};
\addlegendentry{$T_a=20^\circ$C}
% Points de mesure
\addplot[only marks, mark=*, popdark] coordinates {(2.62,28) (3.62,26)};
\node[above left, popdark, font=\tiny] at (axis cs:2.62,28) {22h00 : $28^\circ$C};
\node[above left, popdark, font=\tiny] at (axis cs:3.62,26) {23h00 : $26^\circ$C};
% Heure décès
\draw[popdark, dashed] (axis cs:0,37) -- (axis cs:0,0);
\node[above right, popdark, font=\tiny] at (axis cs:0,37) {Décès $\approx$ 19h23 ($37^\circ$C)};
\end{axis}
\end{tikzpicture}
\legende{Refroidissement du corps : estimation de l'heure du décès par la loi de Newton.}
\end{popfigure}

\begin{exemplebox}[Enquête judiciaire : estimation de l'heure du décès]
Un corps est découvert à 22h00 à $28^\circ\text{C}$ dans une pièce à $20^\circ\text{C}$. Une heure plus tard (23h00), sa température est $26^\circ\text{C}$. On suppose $T_0 = 37^\circ\text{C}$ (vivant) et la loi de Newton.
\begin{enumerate}
    \item Trouver $k$.
    \item Estimer l'heure du décès.
\end{enumerate}
\textbf{Résolution :}
Écart initial $y_0 = 37-20 = 17$. Écart à 22h : $y_1 = 28-20 = 8$. Écart à 23h : $y_2 = 26-20 = 6$.
$y(t) = 17 e^{-kt}$ (origine $t=0$ = heure du décès).
À 22h ($t=t_1$) : $8 = 17 e^{-k t_1}$.
À 23h ($t=t_1+1$) : $6 = 17 e^{-k(t_1+1)} = 8 e^{-k}$.
$\Rightarrow e^{-k} = 6/8 = 0,75 \Rightarrow k = -\ln 0,75 \approx \textbf{0,288 h}\textsuperscript{$-1$}$.
Maintenant $8 = 17 e^{-0,288 t_1} \Rightarrow e^{-0,288 t_1} = 8/17 \approx 0,4706$.
$-0,288 t_1 = \ln 0,4706 \approx -0,754 \Rightarrow t_1 \approx \textbf{2,62 h} \approx 2\text{h }37\text{min}$.
Le décès a eu lieu environ \textbf{2h37 avant 22h00}, soit vers \textbf{19h23}.
\end{exemplebox}

\courssub{Problèmes inverses : retrouver $t$, $k$ ou $y_0$}

La modélisation conduit souvent à des « problèmes inverses » où l'on connaît le résultat et on cherche la cause ou le paramètre.

\begin{methode}[Boîte à outils pour les problèmes inverses]
\begin{enumerate}
    \item \textbf{Chercher $t$ (temps pour atteindre un seuil $y_c$) :} $y_c = y_0 e^{kt} \Rightarrow t = \dfrac{1}{k} \ln\left(\dfrac{y_c}{y_0}\right)$.
    \item \textbf{Chercher $k$ (taux) à partir de deux mesures $(t_1, y_1), (t_2, y_2)$ :} $\dfrac{y_2}{y_1} = e^{k(t_2-t_1)} \Rightarrow k = \dfrac{1}{t_2-t_1} \ln\left(\dfrac{y_2}{y_1}\right)$.
    \item \textbf{Chercher $y_0$ (valeur initiale) :} $y_0 = y(t) e^{-kt}$.
    \item \textbf{Temps de doublement (croissance, $k>0$) :} $T_2 = \dfrac{\ln 2}{k}$.
    \item \textbf{Demi-vie (décroissance, $k<0$, écrire $k=-\lambda$) :} $T_{1/2} = \dfrac{\ln 2}{\lambda}$.
\end{enumerate}
\end{methode}

\begin{exoresolu}[Croissance bactérienne : détermination du temps de doublement]
\textbf{Énoncé :} Une culture bactérienne passe de $10^3$ à $10^6$ bactéries en 5 heures. Croissance exponentielle supposée.
\begin{enumerate}
    \item Déterminer $k$ et le temps de doublement $T_2$.
    \item Combien de bactéries après 24h ?
\end{enumerate}
\tcblower
\textbf{Solution :}
\begin{enumerate}
    \item $k = \dfrac{1}{5} \ln\left(\dfrac{10^6}{10^3}\right) = \dfrac{\ln 1000}{5} = \dfrac{3\ln 10}{5} \approx \dfrac{6,9078}{5} \approx \textbf{1,38 h}\textsuperscript{$-1$}$.
    $T_2 = \dfrac{\ln 2}{k} \approx \dfrac{0,693}{1,38} = \textbf{0,5 h} = \textbf{30 min}$.
    \item $N(24) = 10^3 e^{1,38 \times 24} = 10^3 e^{33,12} \approx 10^3 \times 2,7 \times 10^{14} \approx \textbf{2,7 $\times$ 10}^{17}$ bactéries.
    \emph{Note :} Ce nombre astronomique montre les limites du modèle exponentiel pur (épuisement des ressources, place). Voir complément Gompertz.
\end{enumerate}
\end{exoresolu}

\courssub{Schéma synthèse : du réel au modèle mathématique}

\begin{popfigure}
\begin{tikzpicture}[
    node distance=1.5cm and 2cm,
    block/.style={draw, rounded corners, fill=popblueL, thick, minimum width=3cm, minimum height=1.2cm, align=center, font=\small},
    arrow/.style={-Stealth, thick, popdark},
    math/.style={draw, ellipse, fill=popgreenL, thick, minimum width=3.5cm, minimum height=1.2cm, align=center, font=\small},
    real/.style={draw, rectangle, fill=poporangeL, thick, minimum width=3.5cm, minimum height=1.2cm, align=center, font=\small},
]
\node[real, align=center] (real){Phénomène réel\\ (Bactéries, Capital, Radioactivité,\\ Médicament, Température...)};
\node[math, right=of real, align=center] (diff){Modélisation :\\ $\frac{dy}{dt} = k y$\\ (Taux relatif constant)};
\node[math, right=of diff, align=center] (sol){Résolution maths :\\ $y(t) = y_0 e^{kt}$\\ (Étude : limites, dérivées,\\ variations, primitives)};
\node[real, right=of sol, align=center] (interp){Interprétation \& Prédiction :\\ Temps de doublement / Demi-vie,\\ Date de décès, Âge fossile,\\ Concentration efficace, Capital futur};
\draw[arrow] (real) -- node[above, font=\tiny, align=center]{Observation \\ Hypothèse $y' \propto y$} (diff);
\draw[arrow] (diff) -- node[above, font=\tiny, align=center]{Intégration \\ $\int dy/y = \int k dt$} (sol);
\draw[arrow] (sol) -- node[above, font=\tiny, align=center]{Calculs \\ $t, k, y_0, T_2$} (interp);
\draw[arrow, bend right=40] (interp) to node[below, font=\tiny, align=center]{Validation / Ajustement \\ (Régression exponentielle)} (real);
\end{tikzpicture}
\legende{Cycle de la modélisation exponentielle : du phénomène observé à l'équation différentielle, sa résolution analytique, et l'exploitation concrète des paramètres.}
\end{popfigure}

\courssub{Compléments et ouverture}

\begin{savaistu}[Le Carbone 14 et Toumaï]
La datation au $^{14}\text{C}$ (limite $\approx 50\,000$ ans) ne suffit pas pour dater \emph{Sahelanthropus tchadensis} (Toumaï, $\sim 7$ millions d'années), découvert au Tchad en 2001. Pour ces âges, on utilise d'autres couples isotopiques à demi-vie beaucoup plus longue : \textbf{Potassium-Argon} ($^{40}\text{K} \to ^{40}\text{Ar}$, $T_{1/2} \approx 1,25$ milliard d'années) ou \textbf{Uranium-Plomb} (zircons, $T_{1/2} \sim 4,5$ milliards d'années). Le principe mathématique reste \textbf{exactement le même} : $N = N_0 e^{-\lambda t}$.
\end{savaistu}

\begin{aretenir}[Limites du modèle exponentiel pur]
Le modèle $y' = ky$ (croissance/décroissance illimitée) est un \textbf{modèle de référence} (premier ordre). Il échoue quand :
\begin{itemize}
    \item Ressources limitées (nourriture, place) $\Rightarrow$ \textbf{Logistique (Verhulst)} : $y' = ky(1 - y/M)$.
    \item Taux de variation relatif non constant (vieillissement, saturation) $\Rightarrow$ \textbf{Gompertz} : $y' = ky \ln(M/y)$.
    \item Phénomènes discrets (générations non chevauchantes) $\Rightarrow$ Suites géométriques / Modèles de Leslie.
\end{itemize}
\end{aretenir}

\begin{propriete}[Loi de Gompertz (Croissance tumorale / Épidémies)]
Équation différentielle : $\dfrac{dy}{dt} = k y \ln\left(\dfrac{M}{y}\right)$ avec $M$ capacité maximale (asymptote).
Solution : $y(t) = M \exp\left( \ln(y_0/M) e^{-kt} \right)$.
\begin{itemize}
    \item Croissance rapide au début (quasi exponentielle), puis ralentissement progressif vers $M$.
    \item Plus réaliste que l'exponentiel pur pour les tumeurs solides ou la diffusion d'innovations.
\end{itemize}
\end{propriete}

\begin{exercice}[Entraînement : Modélisation variée]
\begin{enumerate}
    \item \textbf{Finance :} Un capital de $2\,000\,000$ FCFA est placé à un \textbf{taux annuel effectif} de $6\%$ (capitalisation continue). Calculer le capital après 8 ans et le temps nécessaire pour tripler le capital.
    \item \textbf{Chimie :} La concentration d'un réactif suit $C(t) = 0,8 e^{-0,02 t}$ (mol/L, $t$ en minutes). Calculer la demi-vie. Au bout de combien de temps la concentration est-elle divisée par 10 ?
    \item \textbf{Refroidissement :} Une pièce de forge à $800^\circ\text{C}$ est plongée dans l'huile à $30^\circ\text{C}$. Après 2 min, elle est à $500^\circ\text{C}$. Déterminer $k$ et la température après 5 min.
    \item \textbf{Pharmacocinétique :} Dose $D=100$ mg, $V=10$ L. Mesure $C(2\text{h})=4$ mg/L. Trouver $k$, $t_{1/2}$, et la clairance $Cl$.
\end{enumerate}
\end{exercice}

\begin{corrige}[Exercice n°6.8]
\begin{enumerate}
    \item $k = \ln(1,06) \approx 0,05827$. $C(8) = 2\cdot 10^6 e^{0,05827 \times 8} \approx 2\cdot 10^6 \times 1,5938 \approx \textbf{3,19 $\times$ 10}^6$ FCFA. Tripler : $3 = e^{kt} \Rightarrow t = \ln 3 / k \approx 1,0986 / 0,05827 \approx \textbf{18,85 ans}$.
    \item $\lambda = 0,02$. $T_{1/2} = \ln 2 / 0,02 \approx \textbf{34,66 min}$. Divisé par 10 : $0,08 = 0,8 e^{-0,02t} \Rightarrow e^{-0,02t}=0,1 \Rightarrow t = \ln 10 / 0,02 \approx \textbf{115,1 min}$.
    \item $T(t) = 30 + 770 e^{-kt}$. $500 = 30 + 770 e^{-2k} \Rightarrow 470 = 770 e^{-2k} \Rightarrow e^{-2k} = 47/77 \approx 0,6104 \Rightarrow k \approx 0,247$. $T(5) = 30 + 770 e^{-0,247 \times 5} \approx 30 + 770 \times 0,291 \approx \textbf{254}^\circ\text{C}$.
    \item $C_0 = 10$ mg/L. $4 = 10 e^{-2k} \Rightarrow e^{-2k} = 0,4 \Rightarrow k = -\frac{1}{2}\ln 0,4 \approx \textbf{0,458 h}\textsuperscript{$-1$}$. $t_{1/2} = \ln 2 / 0,458 \approx \textbf{1,51 h}$. $Cl = k V = 0,458 \times 10 = \textbf{4,58 L/h}$.
\end{enumerate}
\end{corrige}

\coursec{Compléments et ouverture : Éuler, complexes et suites (Culture mathématique)}

Dans la section précédente, nous avons vu comment la fonction exponentielle modélise des phénomènes concrets : charge d'un condensateur, refroidissement d'une pièce mécanique, décroissance radioactive, placement à intérêts composés. Dans tous ces modèles, $\exp$ apparaît comme la fonction \emph{naturelle} de l'évolution continue. Cette dernière section ouvre des portes : elle montre que la fonction exponentielle est bien plus profonde qu'une simple courbe de croissance. Elle relie entre eux les nombres les plus célèbres des mathématiques, elle est au c\oe ur des calculs de votre calculatrice, et elle gouverne des modèles de populations utilisés en écologie.

\begin{attention}[Statut de cette section]
Le contenu de cette section dépasse le programme officiel de Terminale technique du MINESEC : il ne sera pas exigé au Probatoire ni au Baccalauréat technique. Il s'agit de \cle{culture mathématique} indispensable pour les études post-bac : concours d'entrée (ENS, écoles d'ingénieurs, Polytechnique de Yaoundé, Médecine), classes préparatoires, licences scientifiques. Tout ce qui suit repose uniquement sur des savoirs de Terminale technique : exponentielle, logarithme, suites, dérivation.
\end{attention}

\courssub{La formule d'Euler : l'exponentielle rencontre la trigonométrie}

\textbf{Intuition.} Nous savons calculer $\exp(x)$ pour tout réel $x$. Peut-on donner un sens à $\exp(i\theta)$, où $i$ est le nombre imaginaire vérifiant $i^2=-1$ ? Leonhard Euler (1707--1783) a répondu par l'affirmative, et sa réponse est l'une des plus belles surprises des mathématiques : l'exponentielle d'un imaginaire pur \emph{tourne en rond} sur le cercle trigonométrique.

\begin{definition}[Formule d'Euler (admise)]
Pour tout réel $\theta$, on pose :
\[ e^{i\theta} = \cos\theta + i\sin\theta. \]
Ainsi $e^{i\theta}$ est le nombre complexe de module $1$ et d'argument $\theta$ : c'est le point du \cle{cercle unité} (cercle de centre $O$ et de rayon $1$) repéré par l'angle $\theta$.
\end{definition}

\begin{propriete}[L'identité d'Euler]
En prenant $\theta = \pi$, comme $\cos\pi = -1$ et $\sin\pi = 0$, on obtient :
\[ e^{i\pi} + 1 = 0. \]
Cette unique égalité relie les \textbf{cinq constantes fondamentales} des mathématiques : $0$ (le neutre de l'addition), $1$ (le neutre de la multiplication), $e$ (base de l'exponentielle), $\pi$ (le cercle) et $i$ (l'imaginaire).
\end{propriete}

\begin{popfigure}
\begin{center}
\begin{tikzpicture}[scale=2.4]
\draw[popline] (-1.5,0) -- (1.5,0) node[below left, popmuted] {\small Réel};
\draw[popline] (0,-1.3) -- (0,1.4) node[below left, popmuted] {\small Imaginaire};
\draw[popblue, thick] (0,0) circle (1);
\draw[->, poporange, very thick] (0.55,0) arc (0:60:0.55) node[midway, right] {\small $\theta$};
\draw[->, popdark, thick] (0,0) -- (60:1);
\fill[poporange] (60:1) circle (0.03) node[above right, popdark] {\small $e^{i\theta}=\cos\theta+i\sin\theta$};
\draw[popgreen, thick] (0,0) -- (0.5,0) node[midway, below] {\small $\cos\theta$};
\draw[popgreen, thick] (0.5,0) -- (0.5,0.866) node[midway, right] {\small $\sin\theta$};
\draw[popline] (0.5,0) rectangle (0.42,0.08);
\fill[popblue] (-1,0) circle (0.03) node[below left, popdark] {\small $e^{i\pi}=-1$};
\fill[popblue] (1,0) circle (0.03) node[below right, popdark] {\small $e^{i0}=1$};
\fill[popblue] (0,1) circle (0.03) node[above right, popdark] {\small $e^{i\pi/2}=i$};
\end{tikzpicture}
\end{center}
\legende{Le point $e^{i\theta}$ parcourt le cercle unité quand $\theta$ varie.}
\end{popfigure}

\begin{exemplebox}[Quelques valeurs remarquables]
\begin{itemize}
\item $e^{i\pi/2} = \cos\dfrac{\pi}{2} + i\sin\dfrac{\pi}{2} = i$ ;
\item $e^{i\pi} = -1$ (identité d'Euler) ;
\item $e^{2i\pi} = 1$ : après un tour complet, on revient au point de départ ;
\item $e^{i\pi/3} = \dfrac{1}{2} + i\dfrac{\sqrt{3}}{2}$.
\end{itemize}
Les règles de calcul habituelles restent valables : $e^{i\theta}\times e^{i\theta'} = e^{i(\theta+\theta')}$, ce qui redémontre élégamment les formules d'addition de $\cos$ et $\sin$.
\end{exemplebox}

\begin{savaistu}[La plus belle formule des mathématiques]
Le physicien Richard Feynman (prix Nobel 1965) qualifiait la formule d'Euler de \og joyau \fg{} et de \og plus remarquable formule des mathématiques \fg{}. L'identité $e^{i\pi}+1=0$ arrive régulièrement en tête des sondages sur la \og plus belle équation \fg{}. Euler, mathématicien suisse presque aveugle en fin de vie, a publié plus de 800 articles : c'est lui qui a imposé les notations $e$, $i$, $\pi$ et $f(x)$ que vous utilisez chaque jour.
\end{savaistu}

\courssub{Définition par série : comment la calculatrice calcule $\exp(x)$}

\textbf{Intuition.} Votre calculatrice affiche $e^{0,5} \approx 1,6487$. Comment fait-elle ? Elle ne connaît ni courbe ni table : elle \emph{additionne des puissances}. L'idée est de chercher une fonction égale à sa dérivée sous forme de somme infinie de monômes.

\begin{propriete}[Développement en série de l'exponentielle (admis)]
Pour tout réel $x$ :
\[ \exp(x) = \sum_{n=0}^{+\infty} \frac{x^n}{n!} = 1 + x + \frac{x^2}{2} + \frac{x^3}{6} + \frac{x^4}{24} + \cdots \]
où $n! = 1\times 2\times \cdots \times n$ (factorielle de $n$), avec $0! = 1$.
\end{propriete}

\begin{experience}[Pourquoi cette série est-elle égale à sa dérivée ?]
Dérivons terme à terme (opération admise ici) :
\begin{align*}
\frac{d}{dx}\left(1 + x + \frac{x^2}{2} + \frac{x^3}{6} + \frac{x^4}{24} + \cdots\right) &= 0 + 1 + \frac{2x}{2} + \frac{3x^2}{6} + \frac{4x^3}{24} + \cdots\\
&= 1 + x + \frac{x^2}{2} + \frac{x^3}{6} + \cdots
\end{align*}
On retombe sur la série de départ : elle vérifie $f'=f$, et $f(0)=1$. C'est donc bien $\exp$, par unicité de la solution de $y'=y$ avec $y(0)=1$.
\end{experience}

\begin{exemplebox}[Calcul de $e$ à la main]
En prenant $x=1$ et en sommant les premiers termes :
\begin{align*}
e &\approx 1 + 1 + \frac{1}{2} + \frac{1}{6} + \frac{1}{24} + \frac{1}{120} + \frac{1}{720} + \frac{1}{5040}\\
&\approx 1 + 1 + 0,5 + 0,16667 + 0,04167 + 0,00833 + 0,00139 + 0,00020\\
&\approx 2,71825
\end{align*}
Avec une dizaine de termes on obtient $e \approx 2,718281828459045\ldots$ : c'est exactement le procédé utilisé par les ordinateurs, qui ajoutent des termes jusqu'à ce que la précision voulue soit atteinte (les termes deviennent vite minuscules car $n!$ explose).
\end{exemplebox}

\courssub{La limite caractéristique : $\exp(x)$ comme limite d'intérêts composés}

\textbf{Intuition.} Placez 1 FCFA à un taux de $100\,\%$ par an. Intérêts versés une fois : vous avez $2$ FCFA. Versés deux fois (à $50\,\%$ chaque fois) : $(1+\tfrac12)^2 = 2,25$ FCFA. Douze fois : $(1+\tfrac1{12})^{12} \approx 2,613$ FCFA. Et si la banque verse les intérêts \emph{à chaque instant} ? La fortune tend vers $e \approx 2,718$ FCFA.

\begin{propriete}[Limite caractéristique de l'exponentielle]
Pour tout réel $x$ :
\[ \exp(x) = \lim_{n\to+\infty}\left(1+\frac{x}{n}\right)^n. \]
En particulier, pour $x=1$ : $e = \displaystyle\lim_{n\to+\infty}\left(1+\frac{1}{n}\right)^n$.
\end{propriete}

\begin{experience}[Démonstration via le logarithme]
Posons $u_n = \left(1+\dfrac{x}{n}\right)^n$ et passons au logarithme :
\[ \ln(u_n) = n\,\ln\left(1+\frac{x}{n}\right) = x\times\frac{\ln\left(1+\frac{x}{n}\right)}{\frac{x}{n}}. \]
Posons $h = \dfrac{x}{n}$. Quand $n\to+\infty$, $h\to 0$. Or le taux d'accroissement de $\ln$ en $1$ donne :
\[ \lim_{h\to 0}\frac{\ln(1+h)}{h} = \ln'(1) = \frac{1}{1} = 1. \]
Donc $\ln(u_n) \to x\times 1 = x$, et par continuité de $\exp$ : $u_n = e^{\ln(u_n)} \to e^x$. \hfill$\square$
\end{experience}

\begin{exemplebox}[Intérêts composés infiniment fractionnés]
Un capital $C_0$ placé au taux annuel $t$, avec intérêts composés $n$ fois par an pendant 1 an, devient $C_0\left(1+\dfrac{t}{n}\right)^n$. Quand $n\to+\infty$ (capitalisation \cle{continue}), il tend vers $C_0\,e^{t}$. Pour $t = 5\,\%$ : le facteur continu est $e^{0,05} \approx 1,0513$, contre $(1+\tfrac{0,05}{12})^{12} \approx 1,0512$ pour une capitalisation mensuelle : la différence est faible, mais le modèle continu est bien plus simple à étudier.
\end{exemplebox}

\courssub{Suites récurrentes définies avec $\exp$ : le modèle de Ricker}

\textbf{Intuition.} En écologie, on modélise la taille d'une population de poissons d'une année sur l'autre par une relation de récurrence. Le modèle de \cle{Ricker} (1954) s'écrit $u_{n+1} = r\,u_n\,e^{-u_n}$ : le facteur $e^{-u_n}$ traduit la \emph{compétition} (plus la population est dense, plus la reproduction est freinée). La question centrale : la population se stabilise-t-elle ?

\begin{methode}[Étudier la convergence d'une suite $u_{n+1}=f(u_n)$]
\begin{enumerate}
\item Chercher les \cle{points fixes} : résoudre $f(x)=x$.
\item Étudier $f$ (dérivée, variations) pour encadrer la suite.
\item Montrer que la suite est monotone et bornée, donc convergente.
\item La limite $\ell$ vérifie nécessairement $f(\ell)=\ell$ (par continuité de $f$).
\end{enumerate}
\end{methode}

\begin{exoresolu}[Convergence de la suite $u_{n+1}=2e^{-u_n}$]
\textbf{Énoncé.} Soit $(u_n)$ définie par $u_0 = 0$ et $u_{n+1} = 2e^{-u_n}$. On admet que $(u_n)$ converge. Déterminer sa limite et donner une valeur approchée.
\tcblower
\textbf{Solution.} La fonction $f(x)=2e^{-x}$ est continue sur $\mathbb{R}$. Si $u_n\to\ell$, alors en passant à la limite dans $u_{n+1}=2e^{-u_n}$ :
\[ \ell = 2e^{-\ell}. \]
Étudions $g(x)=2e^{-x}-x$ : $g'(x)=-2e^{-x}-1<0$, donc $g$ est strictement décroissante. De plus $g(0)=2>0$ et $g(1)=2e^{-1}-1\approx -0,264<0$. Par le théorème des valeurs intermédiaires, l'équation $\ell=2e^{-\ell}$ admet une \textbf{unique} solution, située dans $]0\,;1[$.
Calculons les premiers termes :
\begin{align*}
u_0 &= 0 & u_1 &= 2 & u_2 &= 2e^{-2}\approx 0,271\\
u_3 &= 2e^{-0,271}\approx 1,525 & u_4 &\approx 0,435 & u_5 &\approx 1,294\\
u_6 &\approx 0,548 & u_7 &\approx 1,156 & u_8 &\approx 0,629
\end{align*}
Les valeurs encadrent de mieux en mieux la limite. En itérant davantage : $u_{20}\approx 0,853$. La limite est $\ell \approx 0,853$, unique solution de $x=2e^{-x}$ (on vérifie : $2e^{-0,853}\approx 0,853$).
\end{exoresolu}

\begin{popfigure}
\begin{center}
\begin{tikzpicture}[scale=4.2]
\draw[->, popline] (-0.1,0) -- (2.3,0) node[below left, popmuted] {\small $u_n$};
\draw[->, popline] (0,-0.1) -- (0,2.3) node[below left, popmuted] {\small $u_{n+1}$};
\draw[popmuted, dashed] (0,0) -- (2.2,2.2) node[above left, popmuted] {\small $y=x$};
\draw[popblue, very thick, domain=0:2.2, samples=100] plot (\x,{2*exp(-\x)}) node[above left, popblue] {\small $y=2e^{-x}$};
\fill[poporange] (0.8526,0.8526) circle (0.025) node[below right, popdark] {\small point fixe $\ell\approx 0,85$};
\draw[poporange, thick, ->] (0,0) -- (0,2);
\draw[poporange, thick, ->] (0,2) -- (2,2);
\draw[poporange, thick, ->] (2,2) -- (2,0.2707);
\draw[poporange, thick, ->] (2,0.2707) -- (0.2707,0.2707);
\draw[poporange, thick, ->] (0.2707,0.2707) -- (0.2707,1.5254);
\draw[poporange, thick, ->] (0.2707,1.5254) -- (1.5254,1.5254);
\draw[poporange, thick, ->] (1.5254,1.5254) -- (1.5254,0.4354);
\draw[poporange, thick, ->] (1.5254,0.4354) -- (0.4354,0.4354);
\draw[poporange, thick, ->] (0.4354,0.4354) -- (0.4354,1.2938);
\draw[poporange, thick, ->] (0.4354,1.2938) -- (1.2938,1.2938);
\draw[poporange, thick, ->] (1.2938,1.2938) -- (1.2938,0.5479);
\end{tikzpicture}
\end{center}
\legende{Construction en \og toile d'araignée \fg{} de $u_{n+1}=2e^{-u_n}$ : les itérés spiralent autour du point fixe $\ell\approx 0,85$.}
\end{popfigure}

\begin{attention}[Piège : une suite récurrente ne converge pas toujours]
Pour le modèle de Ricker $u_{n+1}=r\,u_n\,e^{-u_n}$, tout dépend du paramètre $r$ : pour $r$ petit, la population se stabilise ; quand $r$ augmente, elle oscille entre deux valeurs, puis quatre (phénomène de \cle{bifurcation}), puis devient chaotique. Une suite de la forme $u_{n+1}=\exp(u_n)$, elle, diverge toujours vers $+\infty$ (car $e^x > x$ pour tout $x$, donc la suite est strictement croissante et non majorée). Ne jamais conclure à la convergence sans l'avoir justifiée !
\end{attention}

\courssub{Équation différentielle $y''-y=0$ et fonctions hyperboliques}

\textbf{Intuition.} Au programme, vous avez résolu $y'=ay$ (solutions $Ce^{ax}$). En mécanique (oscillations, cordes suspendues, poutres), on rencontre des équations du \emph{second ordre} faisant intervenir $y''$. La plus simple est $y''=y$.

\begin{propriete}[Équation $y''-y=0$ (admise)]
Les fonctions $x\mapsto e^{x}$ et $x\mapsto e^{-x}$ sont solutions de $y''-y=0$ sur $\mathbb{R}$. Toutes les solutions s'écrivent :
\[ y(x) = A\,e^{x} + B\,e^{-x}, \qquad (A,B)\in\mathbb{R}^2. \]
\end{propriete}

En effet, si $y=e^x$ alors $y''=e^x=y$ ; de même pour $e^{-x}$. Par linéarité, toute combinaison $Ae^x+Be^{-x}$ convient.

\begin{definition}[Fonctions hyperboliques]
On définit le \cle{cosinus hyperbolique} et le \cle{sinus hyperbolique} :
\[ \operatorname{ch}(x)=\frac{e^{x}+e^{-x}}{2}, \qquad \operatorname{sh}(x)=\frac{e^{x}-e^{-x}}{2}. \]
\end{definition}

\begin{propriete}[Propriétés des fonctions hyperboliques]
Pour tout réel $x$ :
\[ \operatorname{ch}'(x)=\operatorname{sh}(x), \qquad \operatorname{sh}'(x)=\operatorname{ch}(x), \qquad \operatorname{ch}^2(x)-\operatorname{sh}^2(x)=1. \]
\end{propriete}

\begin{experience}[Vérification de $\operatorname{ch}^2-\operatorname{sh}^2=1$]
\begin{align*}
\operatorname{ch}^2(x)-\operatorname{sh}^2(x) &= \frac{(e^x+e^{-x})^2}{4}-\frac{(e^x-e^{-x})^2}{4}\\
&= \frac{(e^{2x}+2+e^{-2x})-(e^{2x}-2+e^{-2x})}{4} = \frac{4}{4}=1.
\end{align*}
Cette identité ressemble à $\cos^2+\sin^2=1$, avec un signe $-$ : d'où le nom \og hyperbolique \fg{} (le point $(\operatorname{ch} t\,;\operatorname{sh} t)$ parcourt une hyperbole, comme $(\cos t\,;\sin t)$ parcourt le cercle).
\end{experience}

\begin{savaistu}[La chaînette : une courbe du quotidien]
Un câble électrique suspendu entre deux poteaux, ou une chaîne tenue par ses deux extrémités, épouse la courbe $y=a\,\operatorname{ch}\!\left(\dfrac{x}{a}\right)$, appelée \cle{chaînette}. Leibniz, Huygens et Jean Bernoulli l'ont découverte en 1691. On la retrouve partout : lignes à haute tension du Cameroun, toiles d'araignée, arches de ponts. C'est une application directe de $e^x+e^{-x}$ !
\end{savaistu}

\courssub{L'exponentielle en algorithmique (Python)}

En informatique, la fonction exponentielle est disponible dans toutes les bibliothèques scientifiques. Deux usages typiques :

\begin{itemize}
\item \textbf{Calcul ponctuel} : \texttt{math.exp(x)} renvoie $e^x$ pour un nombre $x$ (calculé par la série de la sous-section 7.2).
\item \textbf{Calcul sur tableaux} : \texttt{numpy.exp(array)} applique $\exp$ à des millions de valeurs d'un coup, ce qui permet les simulations dites de \cle{Monte-Carlo} (tirages aléatoires massifs pour estimer des probabilités ou des intégrales) et la résolution numérique d'équations différentielles.
\end{itemize}

\begin{exemplebox}[Itérer la suite de Ricker en Python]
Pour calculer $u_{20}$ avec $u_{n+1}=2e^{-u_n}$, $u_0=0$ :
\begin{itemize}
\item[] \texttt{import math}
\item[] \texttt{u = 0.0}
\item[] \texttt{for n in range(20):}
\item[] \texttt{\ \ \ \ u = 2.0 * math.exp(-u)}
\item[] \texttt{print(u)} \quad affiche \texttt{0.8526...}, la valeur du point fixe.
\end{itemize}
Cinq lignes suffisent pour retrouver numériquement le résultat de l'exercice résolu plus haut : l'algorithmique et l'analyse se répondent.
\end{exemplebox}

\begin{exercice}[Pour aller plus loin]
\begin{enumerate}
\item Calculer $e^{i\pi/6}$ sous forme algébrique, puis vérifier que $\left(e^{i\pi/6}\right)^6=-1$.
\item En utilisant la série de $\exp$, donner une valeur approchée de $e^{0,1}$ à $10^{-4}$ près (trois termes suffisent-ils ?).
\item Montrer que $\operatorname{ch}(x)+\operatorname{sh}(x)=e^x$ pour tout réel $x$, puis que $\operatorname{ch}$ est paire et $\operatorname{sh}$ impaire.
\item On considère $u_{n+1}=e^{-u_n}$ avec $u_0=0$. Montrer que l'équation $x=e^{-x}$ admet une unique solution $\alpha$ dans $[0\,;1]$ et donner un encadrement de $\alpha$ à $0,1$ près.
\end{enumerate}
\end{exercice}

\begin{corrige}[Exercice : Pour aller plus loin]
\textbf{1.} $e^{i\pi/6}=\cos\dfrac{\pi}{6}+i\sin\dfrac{\pi}{6}=\dfrac{\sqrt{3}}{2}+\dfrac{1}{2}i$. Alors $\left(e^{i\pi/6}\right)^6=e^{i\pi}=-1$.

\textbf{2.} $e^{0,1}\approx 1+0,1+\dfrac{0,01}{2}+\dfrac{0,001}{6}=1+0,1+0,005+0,000167\approx 1,1052$. Le terme suivant vaut $\dfrac{0,0001}{24}\approx 4\times 10^{-6}<10^{-4}$ : trois termes après le $1$ suffisent.

\textbf{3.} $\operatorname{ch}(x)+\operatorname{sh}(x)=\dfrac{e^x+e^{-x}+e^x-e^{-x}}{2}=e^x$. Parité : $\operatorname{ch}(-x)=\dfrac{e^{-x}+e^{x}}{2}=\operatorname{ch}(x)$ (paire) ; $\operatorname{sh}(-x)=\dfrac{e^{-x}-e^{x}}{2}=-\operatorname{sh}(x)$ (impaire).

\textbf{4.} Soit $g(x)=e^{-x}-x$ : $g'(x)=-e^{-x}-1<0$, donc $g$ strictement décroissante, continue, avec $g(0)=1>0$ et $g(1)=e^{-1}-1<0$. Par le TVI, unique $\alpha\in[0\,;1]$ avec $e^{-\alpha}=\alpha$. Tests : $g(0,5)=e^{-0,5}-0,5\approx 0,107>0$ et $g(0,6)=e^{-0,6}-0,6\approx -0,051<0$, donc $0,5<\alpha<0,6$ (en fait $\alpha\approx 0,567$).
\end{corrige}

\begin{aretenir}[L'essentiel de l'ouverture]
\begin{itemize}
\item \textbf{Euler} : $e^{i\theta}=\cos\theta+i\sin\theta$, et le joyau $e^{i\pi}+1=0$.
\item \textbf{Série} : $e^x=\displaystyle\sum_{n=0}^{+\infty}\frac{x^n}{n!}$ : c'est ainsi que les machines calculent.
\item \textbf{Limite} : $e^x=\displaystyle\lim_{n\to+\infty}\left(1+\frac{x}{n}\right)^n$ (intérêts composés continus).
\item \textbf{Suites} : la limite de $u_{n+1}=f(u_n)$ est un point fixe de $f$ ; le modèle de Ricker $ru_ne^{-u_n}$ illustre stabilité et chaos.
\item \textbf{Second ordre} : $y''=y$ a pour solutions $Ae^x+Be^{-x}$, d'où $\operatorname{ch}$ et $\operatorname{sh}$ avec $\operatorname{ch}^2-\operatorname{sh}^2=1$.
\item \textbf{Python} : \texttt{math.exp} et \texttt{numpy.exp} pour simuler et vérifier numériquement.
\end{itemize}
\end{aretenir}

\newpage

\coursec{Activité d'intégration}

\coursec{Leçon N13 : Fonction exponentielle népérienne}

\courssub{Objectifs de l'activité}
\begin{itemize}
    \item Mobiliser l'ensemble des savoirs de la leçon : définition et propriétés de \texttt{exp}, limites, dérivation, étude de fonction, primitives.
    \item Modéliser un phénomène réel de propagation (rumeur, virus, information) en comparant un modèle exponentiel pur (non borné) et un modèle logistique (borné par la population totale).
    \item Interpréter mathématiquement les paramètres : temps initial, population totale, seuil de 50\,\%, vitesse maximale, point d'inflexion.
    \item Rédiger un compte rendu synthétique à destination d'un décideur (le proviseur), en langage clair et argumenté.
\end{itemize}

\courssub{Situation}
\begin{experience}[Contexte : Propagation d'une rumeur au Lycée Technique de Douala]
Le proviseur du Lycée Technique de Douala (effectif total : \textbf{2000 élèves}) est alerté : une rumeur concernant une prétendue annulation des épreuves du Probatoire circule depuis ce matin 8h00. Il souhaite anticiper les mouvements de panique.
Deux modèles mathématiques sont proposés par l'équipe de surveillance pour prédire le nombre d'élèves $N(t)$ ayant entendu la rumeur $t$ heures après 8h00.

\medskip
\noindent \textbf{Modèle 1 : Modèle exponentiel pur (Malthusien)}
On suppose que chaque élève informé en informe d'autres à un rythme constant, sans saturation.
\[
N_1(t) = e^{0,5 t} \qquad (t \ge 0)
\]
Ce modèle prévoit une croissance illimitée.

\medskip
\noindent \textbf{Modèle 2 : Modèle logistique (Verhulst) — \emph{Modèle de référence pour l'activité}}
On tient compte de la capacité maximale du lycée (2000 élèves). La propagation ralentit quand beaucoup sont déjà au courant.
\[
N_2(t) = \frac{2000}{1 + 1999 \, e^{-0,5 t}} \qquad (t \ge 0)
\]
On pose la fonction réduite $f(x) = \dfrac{1}{1 + 1999 e^{-0,5 x}}$, de sorte que $N_2(t) = 2000 \, f(t)$.
\end{experience}

\courssub{Consignes}
\begin{enumerate}
    \item \textbf{Conditions initiales.} Vérifier que $N_1(0)=1$ et $N_2(0)=1$. Interpréter concrètement ce résultat à 8h00.
    \item \textbf{Étude de la fonction réduite $f$.} Étudier $f(x) = \dfrac{1}{1 + 1999 e^{-0,5 x}}$ sur $[0; +\infty[$ :
    \begin{itemize}
        \item Limites en $0$ et $+\infty$.
        \item Dérivée $f'(x)$, signe, variations.
        \item Seconde dérivée $f''(x)$, convexité, point d'inflexion.
    \end{itemize}
    \item \textbf{Représentation graphique.} Tracer les courbes $C_1$ (modèle 1) et $C_2$ (modèle 2) sur un même repère (GeoGebra / calculatrice) pour $t \in [0; 20]$. Observer les différences majeures.
    \item \textbf{Seuil des 50\,\% (médiane).} Déterminer l'instant $t_{1/2}$ (en heures et minutes) où 1000 élèves connaissent la rumeur avec le modèle logistique $N_2$.
    \item \textbf{Vitesse maximale de propagation.} Pour le modèle logistique, calculer la valeur maximale de la vitesse $N_2'(t)$ et l'instant $t_{\text{infl}}$ où elle est atteinte. Relier cet instant au point d'inflexion de $f$.
    \item \textbf{Critique des modèles.} Argumenter pourquoi le modèle logistique $N_2$ est plus réaliste que le modèle exponentiel pur $N_1$ pour cette situation (notion de plafond, vitesse relative).
    \item \textbf{Synthèse pour le proviseur.} Rédiger un paragraphe court (5-6 lignes) répondant : « À quelle heure la rumeur aura-t-elle touché la majorité ? Quelle est la vitesse de propagation maximale ? »
\end{enumerate}

\courssub{Ressources à mobiliser}
\begin{propriete}[Formules utiles — Rappels du cours]
\begin{itemize}
    \item $\dfrac{d}{dx} e^{u(x)} = u'(x) e^{u(x)}$ \quad ; \quad $\displaystyle \int e^{u(x)} u'(x) dx = e^{u(x)} + C$.
    \item $\dfrac{d}{dx} \left( \dfrac{1}{u(x)} \right) = -\dfrac{u'(x)}{u(x)^2}$ \quad ; \quad $\dfrac{d}{dx} \ln u(x) = \dfrac{u'(x)}{u(x)}$.
    \item $\lim_{x \to +\infty} e^{-kx} = 0$ ($k>0$) \quad ; \quad $\lim_{x \to 0} e^{-kx} = 1$.
    \item Point d'inflexion : $f''(x_0)=0$ et changement de signe de $f''$ $\iff$ extremum de $f'$.
    \item Équation $e^X = a \ (a>0) \iff X = \ln a$.
\end{itemize}
\end{propriete}

\courssub{Démarche et corrigé commenté}
\begin{exoresolu}[Corrigé détaillé de l'activité d'intégration]
\noindent \textbf{1. Conditions initiales}
\[
N_1(0) = e^{0,5 \times 0} = e^0 = 1.
\]
\[
N_2(0) = \frac{2000}{1 + 1999 e^{0}} = \frac{2000}{1 + 1999} = \frac{2000}{2000} = 1.
\]
\emph{Interprétation :} À 8h00 ($t=0$), un seul élève (le « patient zéro » ou l'initiateur de la rumeur) est au courant. Les deux modèles démarrent correctement.

\medskip
\noindent \textbf{2. Étude de $f(x) = \dfrac{1}{1 + 1999 e^{-0,5 x}}$ sur $[0; +\infty[$}

\noindent \textit{a) Limites}
\begin{itemize}
    \item En $0$ : $f(0) = \frac{1}{1+1999} = \frac{1}{2000} = 0,0005$.
    \item En $+\infty$ : $e^{-0,5 x} \to 0$, donc $1 + 1999 e^{-0,5 x} \to 1$, d'où $f(x) \to 1$.
\end{itemize}
\emph{Conséquence pour $N_2$ :} $N_2(0)=1$ et $N_2(t) \to 2000$ quand $t \to +\infty$. La population totale est l'asymptote horizontale.

\medskip
\noindent \textit{b) Dérivée première et variations}
Posons $u(x) = 1 + 1999 e^{-0,5 x}$. Alors $f(x) = u(x)^{-1}$.
$u'(x) = 1999 \times (-0,5) e^{-0,5 x} = -999,5 \, e^{-0,5 x}$.
\[
f'(x) = -\frac{u'(x)}{u(x)^2} = -\frac{-999,5 e^{-0,5 x}}{(1 + 1999 e^{-0,5 x})^2} = \frac{999,5 \, e^{-0,5 x}}{(1 + 1999 e^{-0,5 x})^2}.
\]
Le numérateur $>0$, le dénominateur $>0$, donc \textbf{$f'(x) > 0$ sur $[0; +\infty[$}.
\emph{Variations :} $f$ est \textbf{strictement croissante} sur $[0; +\infty[$. La rumeur progresse sans jamais reculer.

\medskip
\noindent \textit{c) Dérivée seconde, convexité et point d'inflexion}
Pour simplifier, notons $v(x) = e^{-0,5 x}$. Alors $f'(x) = \frac{999,5 \, v}{(1+1999 v)^2}$.
On dérive $f'$ (quotient) ou on utilise la forme logistique standard : $f' = \frac{k}{2} f (1-f)$ avec $k=0,5$ (vérification laissée en exercice).
On dérive $f'$ :
\[
f''(x) = \frac{d}{dx} \left[ \frac{999,5 e^{-0,5 x}}{(1+1999 e^{-0,5 x})^2} \right].
\]
Factorisons $999,5 e^{-0,5 x}$ au numérateur. Le signe de $f''$ est le signe de :
\[
(1+1999 v)' (1+1999 v) - 2 v (1+1999 v)' \quad \text{(schéma quotient simplifié)}
\]
Calculons rigoureusement :
$f'(x) = A \frac{v}{(1+B v)^2}$ avec $A=999,5, B=1999, v' = -0,5 v$.
\[
f'' = A \frac{v'(1+Bv)^2 - v \cdot 2(1+Bv) B v'}{(1+Bv)^4} = A \frac{v'(1+Bv) [1+Bv - 2Bv]}{(1+Bv)^3} = A \frac{v'(1+Bv)(1 - Bv)}{(1+Bv)^3}.
\]
Comme $A>0, v' = -0,5v < 0$, $1+Bv > 0$ :
Le signe de $f''$ est le signe de $v'(1 - Bv) = (-0,5v)(1 - 1999v)$.
$v>0$, donc signe opposé à $(1 - 1999v)$.
\begin{itemize}
    \item $f''(x) > 0 \iff 1 - 1999v < 0 \iff v > \frac{1}{1999} \iff e^{-0,5x} > \frac{1}{1999} \iff -0,5x > \ln(1/1999) \iff x < \frac{\ln 1999}{0,5}$.
    \item $f''(x) < 0 \iff x > \frac{\ln 1999}{0,5}$.
\end{itemize}
Calcul numérique : $\ln 1999 \approx 7,600$. $x_{\text{infl}} = \frac{7,600}{0,5} = 15,20$ heures.
\emph{Convexité :} $f$ est convexe sur $[0; 15,2]$, concave sur $[15,2; +\infty[$.
\emph{Point d'inflexion :} $x_0 = 15,2$ h. $f(x_0) = \frac{1}{1+1999 \times \frac{1}{1999}} = \frac{1}{2}$.
\emph{Pour $N_2$ :} Point d'inflexion à $t_0 = 15,2$ h, $N_2(t_0) = 1000$ élèves. C'est le moment où la propagation est la plus rapide.

\medskip
\noindent \textbf{Tableau de variations de $f$}
\begin{center}
\begin{tabular}{|c|ccccc|}\hline
$x$ & $0$ & & $15,2$ & & $+\infty$ \\ \hline
$f'(x)$ & & $+$ & & $+$ & \\ \hline
$f(x)$ & $\frac{1}{2000}$ & $\nearrow$ & $\frac{1}{2}$ & $\nearrow$ & $1$ \\ \hline
\end{tabular}
\end{center}
\emph{Remarque :} $f'$ est strictement positive ; le maximum de $f'$ (vitesse max) est atteint en $x=15,2$ (point d'inflexion).

\medskip
\noindent \textbf{3. Représentation graphique}
\begin{popfigure}
\begin{tikzpicture}
\begin{axis}[
    width=\linewidth, height=8cm,
    axis lines=middle,
    xlabel=$t$ (heures), ylabel=$N(t)$ (élèves),
    xmin=0, xmax=20, ymin=0, ymax=2500,
    xtick={0,5,10,15.2,20}, ytick={0,1000,2000},
    xticklabels={0,5,10,$t_0\approx15{,}2$,20},
    yticklabels={0,1000,2000},
    legend pos=north west,
    grid=major, grid style={popline},
]
\addplot[popcurve, domain=0:20, samples=200] {exp(0.5*x)};
\addlegendentry{$N_1(t)=e^{0,5t}$ (Modèle pur)}
\addplot[popcurve, poporange, domain=0:20, samples=200] {2000/(1+1999*exp(-0.5*x))};
\addlegendentry{$N_2(t)=\frac{2000}{1+1999e^{-0,5t}}$ (Logistique)}
\addplot[dashed, popmuted, domain=0:20] {2000};
\addlegendentry{Asymptote $y=2000$}
\addplot[dashed, popmuted] coordinates {(15.2,0) (15.2,1000)};
\node[popdark, anchor=south west] at (axis cs:15.2,1000) {$I(15,2;1000)$};
\end{axis}
\end{tikzpicture}
\legende{Comparaison des deux modèles de propagation sur $[0;20]$ h. Le modèle exponentiel (bleu) diverge ; le modèle logistique (orange) sature à 2000. Point d'inflexion $I$ à $t\approx15,2$ h.}
\end{popfigure}
\emph{Observation :} $C_1$ n'a pas de limite supérieure (irréaliste pour une population close). $C_2$ ralentit naturellement en approchant 2000.

\medskip
\noindent \textbf{4. Temps pour 50\,\% (médiane) — Modèle logistique}
On cherche $t$ tel que $N_2(t) = 1000$.
\[
\frac{2000}{1 + 1999 e^{-0,5 t}} = 1000 \iff 1 + 1999 e^{-0,5 t} = 2 \iff 1999 e^{-0,5 t} = 1.
\]
\[
e^{-0,5 t} = \frac{1}{1999} \iff -0,5 t = \ln(1/1999) = -\ln 1999 \iff t = \frac{\ln 1999}{0,5}.
\]
$\ln 1999 \approx 7,6004$. $t_{1/2} \approx 15,2008$ h.
\emph{Conversion :} $0,2008 \times 60 \approx 12$ minutes.
\emph{Réponse :} \textbf{Vers 23h12 (11h12 du soir)}. La majorité est atteinte 15h12 après le début.

\medskip
\noindent \textbf{5. Vitesse maximale de propagation}
La vitesse est $N_2'(t) = 2000 f'(t)$.
Le maximum de $N_2'$ coïncide avec le maximum de $f'$, donc au point d'inflexion $t_0 = 15,2$ h.
Calculons $f'(t_0)$ : à l'inflexion, $f=1/2$, $v=1/1999$.
\[
f'(t_0) = \frac{999,5 \times \frac{1}{1999}}{(1+1)^2} = \frac{0,5}{4} = 0,125.
\]
(Vérification formule logistique : $f' = \frac{k}{4} = \frac{0,5}{4} = 0,125$).
Donc $N_2'(t_0) = 2000 \times 0,125 = \textbf{250 élèves par heure}$.
\emph{Réponse :} Vitesse max \textbf{250 élèves/heure} atteinte à \textbf{23h12} (même instant que la médiane).

\medskip
\noindent \textbf{6. Pourquoi le modèle logistique est plus réaliste}
\begin{itemize}
    \item \textbf{Plafond populationnel :} $N_1(t) \to +\infty$. Au bout de 24h, $N_1(24) = e^{12} \approx 162\,754$ élèves $>$ 2000. Impossible physiquement. $N_2(t) \to 2000$, respectant la contrainte de l'effectif total.
    \item \textbf{Vitesse relative :} Modèle pur : $\frac{N_1'}{N_1} = 0,5$ constant. Chaque élève « infecte » autant au début qu'à la fin. Modèle logistique : $\frac{N_2'}{N_2} = 0,5 \left(1 - \frac{N_2}{2000}\right)$. La vitesse relative diminue quand la rumeur est connue de tous (moins de « cibles » disponibles).
    \item \textbf{Forme en S :} Correspond à la réalité sociologique : démarrage lent (peu d'informateurs), emballement, puis saturation.
\end{itemize}

\medskip
\noindent \textbf{7. Synthèse pour le proviseur}
\begin{aretenir}[Note pour le proviseur — 23h12]
Monsieur le Proviseur, selon le modèle logistique validé (plafond 2000 élèves), la rumeur aura touché la majorité (1000 élèves) \textbf{ce soir vers 23h12}. La propagation sera la plus rapide exactement à ce moment, avec une vitesse de \textbf{250 élèves par heure}. Après minuit, la progression ralentira naturellement car la quasi-totalité des élèves sera informée. Il est recommandé de communiquer officiellement \textbf{avant 22h00} pour couper court à la rumeur.
\end{aretenir}

\medskip
\noindent \textbf{8. Complément : Primitives et intégrales (Programme officiel)}
\begin{methode}[Calcul de primitives de fonctions exponentielles composées]
Pour intégrer une fonction de la forme $u'(x) e^{u(x)}$ ou $\frac{u'(x)}{u(x)}$, on utilise la substitution $U = u(x)$.
\begin{enumerate}
    \item Identifier $u(x)$ et sa dérivée $u'(x)$ dans l'expression.
    \item Écrire l'intégrale sous la forme $\int u'(x) F(u(x)) dx$.
    \item Conclure : $\int u'(x) e^{u(x)} dx = e^{u(x)} + C$ ; $\int \frac{u'(x)}{u(x)} dx = \ln|u(x)| + C$.
\end{enumerate}
\end{methode}
\begin{exemplebox}[Primitive liée au modèle logistique]
Calculer une primitive de $g(t) = \frac{999,5 e^{-0,5 t}}{(1+1999 e^{-0,5 t})^2}$ sur $\mathbb{R}$.
\emph{Solution :} On reconnaît $f'(t)$ avec $u(t) = 1+1999 e^{-0,5 t}$, $u'(t) = -999,5 e^{-0,5 t}$.
$g(t) = -\frac{u'(t)}{u(t)^2}$. Une primitive est $G(t) = \frac{1}{u(t)} = \frac{1}{1+1999 e^{-0,5 t}} = f(t)$.
Vérification : $G'(t) = -\frac{u'(t)}{u(t)^2} = g(t)$.
\emph{Application :} L'intégrale $\int_0^T N_2'(t) dt = N_2(T) - N_2(0)$ donne le nombre total d'élèves informés entre 8h00 et $8\text{h}+T$.
\end{exemplebox}
\end{exoresolu}

\courssub{Bilan de l'activité}
Cette activité a permis de :
\begin{itemize}
    \item \textbf{Ancrer les notions abstraites} (limites, dérivée, convexité, point d'inflexion) dans une situation concrète de gestion de crise scolaire.
    \item \textbf{Différencier modélisation exponentielle et logistique :} comprendre qu'un modèle exponentiel pur n'est valable qu'en phase initiale (ressources illimitées) et qu'un modèle à saturation (logistique) est indispensable pour une population close.
    \item \textbf{Maîtriser la chaîne de calcul complète} : de l'expression algébrique à l'interprétation horaire (conversion décimale $\to$ heures/minutes), en passant par la résolution d'équations exponentielles ($\ln$).
    \item \textbf{Développer la compétence « Communiquer » :} produire un avis décisionnel court, chiffré et argumenté à destination d'un non-mathématicien (le proviseur).
    \item \textbf{Valider l'usage des outils numériques} (GeoGebra/calculatrice) pour conjecturer (asymptote, inflexion) avant de démontrer analytiquement.
    \item \textbf{Mobiliser les primitives :} relier la vitesse instantanée $N_2'$ à la variation du nombre total d'informés via l'intégrale.
\end{itemize}
\begin{savaistu}[Le modèle logistique au Cameroun et dans le monde]
Le modèle logistique $\frac{N'}{N} = r(1-\frac{N}{K})$ a été introduit par \textbf{Pierre-François Verhulst} (1838) pour modéliser la population belge. Au Cameroun, il est utilisé par le \textbf{BUCREP} (Bureau Central des Recensements et Études de Population) pour projeter la démographie (population jeune, $K$ estimé par les ressources). En épidémiologie (choléra, COVID-19, rougeole au Nord), le modèle SIR (Susceptibles-Infectés-Rétablis) généralise la logistique. Le « $R_0$ » (taux de reproduction de base) correspond à notre $k=0,5$ ici. Si $R_0 > 1$, l'épidémie démarre ; la vaccination vise à faire baisser le $K$ effectif (immunité collective).
\end{savaistu}

\newpage

\coursec{Méthodes et exercices résolus}

\coursec{Fonction exponentielle népérienne}

\courssub{Méthode 1 : Appliquer les propriétés algébriques de la fonction exponentielle}

\begin{methode}[Simplifier et transformer une expression avec $\mathrm{e}^x$]
Pour simplifier une expression contenant des exponentielles :
\begin{enumerate}
\item Repérer les formes : produit $\mathrm{e}^a \times \mathrm{e}^b$, quotient $\dfrac{\mathrm{e}^a}{\mathrm{e}^b}$, puissance $\left(\mathrm{e}^a\right)^n$, inverse $\mathrm{e}^{-a}$.
\item Appliquer les formules fondamentales :
\[ \mathrm{e}^a \times \mathrm{e}^b = \mathrm{e}^{a+b} \qquad \dfrac{\mathrm{e}^a}{\mathrm{e}^b} = \mathrm{e}^{a-b} \qquad \left(\mathrm{e}^a\right)^n = \mathrm{e}^{na} \qquad \mathrm{e}^{-a} = \dfrac{1}{\mathrm{e}^a} \]
\item Regrouper en une seule exponentielle $\mathrm{e}^{A}$ où $A$ est une expression simplifiée.
\item Utiliser si besoin les valeurs remarquables : $\mathrm{e}^0 = 1$ et $\mathrm{e}^1 = \mathrm{e} \approx 2{,}718$.
\end{enumerate}
\textbf{Réflexe :} tout se ramène à \emph{une seule} exponentielle. Ne jamais écrire $\mathrm{e}^a + \mathrm{e}^b = \mathrm{e}^{a+b}$ : c'est faux !
\end{methode}

\begin{exoresolu}[Simplification d'expressions exponentielles]
\textbf{Énoncé.} Simplifier les expressions suivantes et les écrire sous la forme $\mathrm{e}^{A}$ où $A$ est une expression la plus simple possible :
\[ A = \mathrm{e}^{2x+1} \times \mathrm{e}^{3-x} \qquad ; \qquad B = \dfrac{\mathrm{e}^{5x}}{\mathrm{e}^{2x-3}} \qquad ; \qquad C = \dfrac{\left(\mathrm{e}^{x}\right)^3 \times \mathrm{e}^{-2x}}{\mathrm{e}^{x-4}} \]
\tcblower
\textbf{Solution détaillée.}

\textbf{Calcul de $A$.} On applique la formule du produit : $\mathrm{e}^a \times \mathrm{e}^b = \mathrm{e}^{a+b}$.
\[ A = \mathrm{e}^{2x+1} \times \mathrm{e}^{3-x} = \mathrm{e}^{(2x+1)+(3-x)} = \mathrm{e}^{2x+1+3-x} = \mathrm{e}^{x+4} \]

\textbf{Calcul de $B$.} On applique la formule du quotient : $\dfrac{\mathrm{e}^a}{\mathrm{e}^b} = \mathrm{e}^{a-b}$.
\[ B = \dfrac{\mathrm{e}^{5x}}{\mathrm{e}^{2x-3}} = \mathrm{e}^{5x-(2x-3)} = \mathrm{e}^{5x-2x+3} = \mathrm{e}^{3x+3} \]
\emph{Attention au signe :} soustraire $2x-3$ donne $-2x+3$.

\textbf{Calcul de $C$.} On traite d'abord le numérateur.
\begin{itemize}
\item Puissance : $\left(\mathrm{e}^{x}\right)^3 = \mathrm{e}^{3x}$.
\item Produit au numérateur : $\mathrm{e}^{3x} \times \mathrm{e}^{-2x} = \mathrm{e}^{3x-2x} = \mathrm{e}^{x}$.
\item Quotient : $\dfrac{\mathrm{e}^{x}}{\mathrm{e}^{x-4}} = \mathrm{e}^{x-(x-4)} = \mathrm{e}^{x-x+4} = \mathrm{e}^{4}$.
\end{itemize}

\textbf{Conclusion.} $A = \mathrm{e}^{x+4}$, $B = \mathrm{e}^{3x+3}$ et $C = \mathrm{e}^{4}$ (constante indépendante de $x$).
\end{exoresolu}

\courssub{Méthode 2 : Résoudre une équation ou une inéquation exponentielle}

\begin{methode}[Résoudre $\mathrm{e}^{A} = \mathrm{e}^{B}$ ou $\mathrm{e}^{A} \leq \mathrm{e}^{B}$]
\begin{enumerate}
\item Mettre l'équation ou l'inéquation sous la forme $\mathrm{e}^{A} = \mathrm{e}^{B}$ (ou $\mathrm{e}^{A} \leq \mathrm{e}^{B}$, $\mathrm{e}^{A} \geq \mathrm{e}^{B}$) en utilisant les propriétés algébriques.
\item Utiliser l'injectivité de la fonction exponentielle : $\mathrm{e}^{A} = \mathrm{e}^{B} \iff A = B$.
\item Pour une inéquation : la fonction exp est \cle{strictement croissante} sur $\mathbb{R}$, donc
\[ \mathrm{e}^{A} \leq \mathrm{e}^{B} \iff A \leq B \qquad \text{(le sens de l'inégalité est conservé)} \]
\item Résoudre l'équation ou l'inéquation obtenue (du premier ou second degré en $x$).
\item Conclure par l'ensemble des solutions, sous forme d'intervalle pour une inéquation.
\end{enumerate}
\textbf{Réflexe :} si un membre est un nombre positif $k$, écrire $k = \mathrm{e}^{\ln k}$. Rappel : $\mathrm{e}^{A} > 0$ pour tout réel $A$, donc $\mathrm{e}^{A} = 0$ ou $\mathrm{e}^{A} < 0$ n'ont \textbf{aucune solution}.
\end{methode}

\begin{exoresolu}[Équation et inéquation exponentielles]
\textbf{Énoncé.}
\begin{enumerate}
\item Résoudre dans $\mathbb{R}$ l'équation $\mathrm{e}^{2x-1} = \mathrm{e}^{x+3}$.
\item Résoudre dans $\mathbb{R}$ l'inéquation $\mathrm{e}^{3x+2} \leq \mathrm{e}^{5}$.
\item Résoudre dans $\mathbb{R}$ l'équation $\mathrm{e}^{x^2 - 4} = 1$.
\end{enumerate}
\tcblower
\textbf{Solution détaillée.}

\textbf{1.} Par injectivité de la fonction exponentielle :
\[ \mathrm{e}^{2x-1} = \mathrm{e}^{x+3} \iff 2x - 1 = x + 3 \iff 2x - x = 3 + 1 \iff x = 4 \]
L'ensemble des solutions est $\mathcal{S} = \{4\}$.

\textbf{2.} La fonction exponentielle est strictement croissante sur $\mathbb{R}$, donc elle conserve l'ordre :
\[ \mathrm{e}^{3x+2} \leq \mathrm{e}^{5} \iff 3x + 2 \leq 5 \iff 3x \leq 3 \iff x \leq 1 \]
L'ensemble des solutions est $\mathcal{S} = \left]-\infty \, ; 1\right]$.

\textbf{3.} On écrit $1 = \mathrm{e}^{0}$. L'équation devient :
\[ \mathrm{e}^{x^2-4} = \mathrm{e}^{0} \iff x^2 - 4 = 0 \iff x^2 = 4 \iff x = 2 \ \text{ou} \ x = -2 \]
L'ensemble des solutions est $\mathcal{S} = \{-2 \, ; 2\}$.

\textbf{Conclusion.} La clé est toujours de « descendre » l'exposant grâce à l'injectivité, puis de résoudre une équation classique.
\end{exoresolu}

\courssub{Méthode 3 : Calculer les limites de référence et les limites de fonctions avec exp}

\begin{methode}[Lever une indétermination avec l'exponentielle]
\begin{enumerate}
\item Connaître par cœur les limites de référence :
\[ \lim_{x \to +\infty} \mathrm{e}^x = +\infty \qquad ; \qquad \lim_{x \to -\infty} \mathrm{e}^x = 0 \qquad ; \qquad \lim_{x \to +\infty} \dfrac{\mathrm{e}^x}{x} = +\infty \qquad ; \qquad \lim_{x \to -\infty} x\,\mathrm{e}^x = 0 \]
\item Pour $\mathrm{e}^{u(x)}$ : calculer d'abord la limite $\ell$ de $u(x)$, puis appliquer la composition :
\[ \text{si } \lim_{x \to a} u(x) = \ell \ \text{ alors } \ \lim_{x \to a} \mathrm{e}^{u(x)} = \lim_{X \to \ell} \mathrm{e}^{X} \]
\item Cas des formes indéterminées : l'\cle{exponentielle l'emporte} sur toute puissance de $x$. Par exemple $\lim\limits_{x \to +\infty} \left(\mathrm{e}^x - x^2\right) = +\infty$ car $\mathrm{e}^x - x^2 = \mathrm{e}^x\left(1 - \dfrac{x^2}{\mathrm{e}^x}\right)$ et $\dfrac{x^2}{\mathrm{e}^x} \to 0$.
\item Interpréter graphiquement : si $\lim\limits_{x \to -\infty} f(x) = k$, la droite $y = k$ est \cle{asymptote horizontale}.
\end{enumerate}
\end{methode}

\begin{exoresolu}[Calculs de limites]
\textbf{Énoncé.} Calculer les limites suivantes, en justifiant :
\[ \text{a) } \lim_{x \to +\infty} \mathrm{e}^{-2x+1} \qquad ; \qquad \text{b) } \lim_{x \to +\infty} \left(3\mathrm{e}^x - 5x\right) \qquad ; \qquad \text{c) } \lim_{x \to -\infty} \left(2 + \mathrm{e}^{x}\right) \]
\tcblower
\textbf{Solution détaillée.}

\textbf{a)} Posons $u(x) = -2x + 1$. On a $\lim\limits_{x \to +\infty} (-2x+1) = -\infty$ (fonction affine de coefficient directeur négatif). Par composition, avec $X = -2x+1 \to -\infty$ :
\[ \lim_{x \to +\infty} \mathrm{e}^{-2x+1} = \lim_{X \to -\infty} \mathrm{e}^{X} = 0 \]

\textbf{b)} Forme indéterminée « $+\infty - \infty$ ». On factorise par le terme prépondérant $\mathrm{e}^x$ :
\[ 3\mathrm{e}^x - 5x = \mathrm{e}^x\left(3 - \dfrac{5x}{\mathrm{e}^x}\right) \]
Or $\lim\limits_{x \to +\infty} \dfrac{x}{\mathrm{e}^x} = 0$ (croissance comparée : l'exponentielle l'emporte), donc $\lim\limits_{x \to +\infty} \dfrac{5x}{\mathrm{e}^x} = 0$ et la parenthèse tend vers $3$. Par produit :
\[ \lim_{x \to +\infty} \mathrm{e}^x \times \lim_{x \to +\infty}\left(3 - \dfrac{5x}{\mathrm{e}^x}\right) = (+\infty) \times 3 = +\infty \]

\textbf{c)} $\lim\limits_{x \to -\infty} \mathrm{e}^{x} = 0$, donc par somme : $\lim\limits_{x \to -\infty} \left(2 + \mathrm{e}^{x}\right) = 2 + 0 = 2$.

\textbf{Conclusion.} a) $0$ ; b) $+\infty$ ; c) $2$ (la droite $y = 2$ est asymptote horizontale en $-\infty$ pour la fonction $x \mapsto 2 + \mathrm{e}^x$).
\end{exoresolu}

\courssub{Méthode 4 : Dériver des fonctions construites avec exp}

\begin{methode}[Dériver $\mathrm{e}^u$ et les produits avec exp]
\begin{enumerate}
\item Formules de base : $(\mathrm{e}^x)' = \mathrm{e}^x$ et, pour $u$ dérivable :
\[ \left(\mathrm{e}^{u}\right)' = u' \times \mathrm{e}^{u} \]
\item Cas particulier très fréquent : $\left(\mathrm{e}^{ax+b}\right)' = a\,\mathrm{e}^{ax+b}$.
\item Si $f = u \times \mathrm{e}^{v}$, appliquer la formule du produit : $f' = u'\,\mathrm{e}^{v} + u \times v'\,\mathrm{e}^{v}$, puis \textbf{factoriser par $\mathrm{e}^{v}$} (qui est toujours strictement positif).
\item Pour étudier le signe de $f'$ : comme $\mathrm{e}^{v} > 0$, le signe de $f'$ est celui du facteur restant (souvent une fonction affine).
\end{enumerate}
\textbf{Réflexe :} après chaque dérivation avec exp, factoriser systématiquement par l'exponentielle.
\end{methode}

\begin{exoresolu}[Dérivation et signe de la dérivée]
\textbf{Énoncé.} Soit $f$ la fonction définie sur $\mathbb{R}$ par $f(x) = (2x - 1)\,\mathrm{e}^{x}$.
\begin{enumerate}
\item Montrer que pour tout réel $x$, $f'(x) = (2x+1)\,\mathrm{e}^{x}$.
\item En déduire le signe de $f'(x)$ et le sens de variation de $f$.
\end{enumerate}
\tcblower
\textbf{Solution détaillée.}

\textbf{1.} $f$ est un produit $f = u \times v$ avec $u(x) = 2x - 1$ et $v(x) = \mathrm{e}^{x}$, toutes deux dérivables sur $\mathbb{R}$. On a $u'(x) = 2$ et $v'(x) = \mathrm{e}^{x}$. Formule du produit : $(uv)' = u'v + uv'$ :
\begin{align*}
f'(x) &= 2 \times \mathrm{e}^{x} + (2x-1) \times \mathrm{e}^{x} \\
&= \mathrm{e}^{x}\left[2 + (2x - 1)\right] \quad \text{(factorisation par } \mathrm{e}^x\text{)} \\
&= (2x + 1)\,\mathrm{e}^{x}
\end{align*}

\textbf{2.} Pour tout réel $x$, $\mathrm{e}^{x} > 0$, donc $f'(x)$ a le même signe que $2x + 1$.
\[ 2x + 1 > 0 \iff x > -\dfrac{1}{2} \qquad ; \qquad 2x + 1 < 0 \iff x < -\dfrac{1}{2} \qquad ; \qquad f'\left(-\tfrac{1}{2}\right) = 0 \]
Donc $f$ est \textbf{strictement décroissante} sur $\left]-\infty \, ; -\dfrac{1}{2}\right]$ et \textbf{strictement croissante} sur $\left[-\dfrac{1}{2} \, ; +\infty\right[$. Elle admet un minimum en $x = -\dfrac{1}{2}$, valant :
\[ f\left(-\dfrac{1}{2}\right) = \left(2 \times \left(-\tfrac{1}{2}\right) - 1\right)\mathrm{e}^{-1/2} = -2\,\mathrm{e}^{-1/2} = -\dfrac{2}{\sqrt{\mathrm{e}}} \]

\textbf{Conclusion.} $f'(x) = (2x+1)\mathrm{e}^x$ ; $f$ décroît puis croît, avec un minimum égal à $-2\mathrm{e}^{-1/2}$ atteint en $x = -\dfrac{1}{2}$.
\end{exoresolu}

\courssub{Méthode 5 : Déterminer des primitives et calculer une intégrale avec exp}

\begin{methode}[Primitives et intégrales de fonctions exponentielles]
\begin{enumerate}
\item Primitives de référence : une primitive de $\mathrm{e}^x$ est $\mathrm{e}^x$ ; plus généralement :
\[ \text{une primitive de } \mathrm{e}^{ax+b} \text{ est } \dfrac{1}{a}\,\mathrm{e}^{ax+b} \quad (a \neq 0) \]
\item Forme $u'\,\mathrm{e}^{u}$ : une primitive est $\mathrm{e}^{u}$. \textbf{Réflexe :} vérifier que le facteur devant l'exponentielle est exactement la dérivée de l'exposant (à une constante multiplicative près, que l'on sort).
\item Linéarité : une primitive de $\alpha f + \beta g$ est $\alpha F + \beta G$.
\item Intégrale : $\displaystyle\int_{a}^{b} f(x)\,\mathrm{d}x = F(b) - F(a)$, noté $[F(x)]_a^b$. L'intégrale d'une fonction positive sur $[a\,;b]$ ($a<b$) s'interprète comme une \cle{aire} en unités d'aire.
\end{enumerate}
\textbf{Vérification :} toujours redériver mentalement la primitive trouvée pour contrôler le résultat.
\end{methode}

\begin{exoresolu}[Primitives et calcul d'aire]
\textbf{Énoncé.} Dans un atelier, la vitesse d'usinage (en pièces par heure) d'une machine en chauffe est modélisée par $v(t) = 40\,\mathrm{e}^{0{,}5t}$ pour $t \in [0\,;2]$ ($t$ en heures).
\begin{enumerate}
\item Déterminer une primitive $V$ de $v$ sur $[0\,;2]$.
\item Calculer $\displaystyle\int_{0}^{2} v(t)\,\mathrm{d}t$. Donner la valeur exacte puis une valeur arrondie à l'unité, et interpréter le résultat.
\end{enumerate}
\tcblower
\textbf{Solution détaillée.}

\textbf{1.} $v$ est de la forme $40\,\mathrm{e}^{0{,}5t}$. Une primitive de $\mathrm{e}^{0{,}5t}$ est $\dfrac{1}{0{,}5}\,\mathrm{e}^{0{,}5t} = 2\,\mathrm{e}^{0{,}5t}$. Par linéarité (multiplication par la constante $40$) :
\[ V(t) = 40 \times 2\,\mathrm{e}^{0{,}5t} = 80\,\mathrm{e}^{0{,}5t} \]
\emph{Vérification :} $V'(t) = 80 \times 0{,}5\,\mathrm{e}^{0{,}5t} = 40\,\mathrm{e}^{0{,}5t} = v(t)$. Correct.

\textbf{2.} On applique la formule fondamentale :
\begin{align*}
\int_{0}^{2} v(t)\,\mathrm{d}t &= \left[80\,\mathrm{e}^{0{,}5t}\right]_{0}^{2} = V(2) - V(0) \\
&= 80\,\mathrm{e}^{0{,}5 \times 2} - 80\,\mathrm{e}^{0} = 80\,\mathrm{e}^{1} - 80 \times 1 = 80(\mathrm{e} - 1)
\end{align*}
Valeur approchée : $80(\mathrm{e}-1) \approx 80 \times (2{,}718 - 1) = 80 \times 1{,}718 \approx 137{,}5$, soit environ $137$.

\textbf{Conclusion.} $\displaystyle\int_{0}^{2} v(t)\,\mathrm{d}t = 80(\mathrm{e}-1) \approx 137$ : la machine a usiné environ $137$ pièces pendant ses deux premières heures de fonctionnement.
\end{exoresolu}

\courssub{Méthode 6 : Conduire l'étude complète d'une fonction définie avec exp}

\begin{methode}[Plan d'étude complet de $f(x) = (ax+b)\mathrm{e}^{cx}$ ou similaire]
\begin{enumerate}
\item \textbf{Ensemble de définition} : l'exponentielle est définie sur $\mathbb{R}$ ; vérifier s'il y a un dénominateur.
\item \textbf{Limites aux bornes} : utiliser les limites de référence et la croissance comparée ; en déduire les asymptotes éventuelles.
\item \textbf{Dérivée} : dériver (produit ou composée), factoriser par l'exponentielle.
\item \textbf{Signe de $f'$ et variations} : dresser le tableau de variations complet (limites et valeurs aux extrema).
\item \textbf{Points remarquables} : image de $0$, intersection avec les axes, extremum.
\item \textbf{Courbe} : tracer en plaçant les points remarquables et l'asymptote.
\end{enumerate}
\end{methode}

\begin{exoresolu}[Étude complète type Probatoire]
\textbf{Énoncé.} Soit $f$ la fonction définie sur $\mathbb{R}$ par $f(x) = (x - 2)\mathrm{e}^{x} + 3$ et $\mathcal{C}_f$ sa courbe représentative.
\begin{enumerate}
\item Calculer les limites de $f$ en $-\infty$ et en $+\infty$. Interpréter graphiquement le résultat en $-\infty$.
\item Calculer $f'(x)$ et dresser le tableau de variations de $f$.
\item Calculer $f(0)$ et $f(2)$.
\end{enumerate}
\tcblower
\textbf{Solution détaillée.}

\textbf{1. Limites.}

\emph{En $+\infty$ :} $\lim\limits_{x \to +\infty} (x-2) = +\infty$ et $\lim\limits_{x \to +\infty} \mathrm{e}^{x} = +\infty$, donc par produit $\lim\limits_{x \to +\infty} (x-2)\mathrm{e}^{x} = +\infty$, puis par somme :
\[ \lim_{x \to +\infty} f(x) = +\infty \]

\emph{En $-\infty$ :} on développe : $f(x) = x\,\mathrm{e}^{x} - 2\mathrm{e}^{x} + 3$. Par croissance comparée, $\lim\limits_{x \to -\infty} x\,\mathrm{e}^{x} = 0$, et $\lim\limits_{x \to -\infty} \mathrm{e}^{x} = 0$. Donc :
\[ \lim_{x \to -\infty} f(x) = 0 - 0 + 3 = 3 \]
\emph{Interprétation :} la droite d'équation $y = 3$ est \textbf{asymptote horizontale} à $\mathcal{C}_f$ en $-\infty$.

\textbf{2. Dérivée et variations.} $f$ est dérivable sur $\mathbb{R}$ comme somme et produit de fonctions dérivables. Pour dériver $(x-2)\mathrm{e}^x$, on pose $u(x) = x-2$, $v(x) = \mathrm{e}^x$ : $u'(x) = 1$, $v'(x) = \mathrm{e}^x$.
\[ f'(x) = 1 \times \mathrm{e}^{x} + (x-2)\mathrm{e}^{x} + 0 = \mathrm{e}^{x}\left[1 + x - 2\right] = (x-1)\,\mathrm{e}^{x} \]
Comme $\mathrm{e}^{x} > 0$ pour tout $x$, $f'(x)$ a le signe de $x - 1$ : négatif pour $x < 1$, nul en $x = 1$, positif pour $x > 1$.
\[ f(1) = (1-2)\mathrm{e}^{1} + 3 = 3 - \mathrm{e} \approx 0{,}28 \]

\begin{center}
\begin{tabular}{|c|ccccc|}\hline
$x$ & $-\infty$ & & $1$ & & $+\infty$ \\ \hline
$f'(x)$ & & $-$ & $0$ & $+$ & \\ \hline
$f$ & $3$ & $\searrow$ & $3-\mathrm{e}$ & $\nearrow$ & $+\infty$ \\ \hline
\end{tabular}
\end{center}

\textbf{3. Valeurs.} $f(0) = (0-2)\mathrm{e}^{0} + 3 = -2 + 3 = 1$ et $f(2) = (2-2)\mathrm{e}^{2} + 3 = 3$.

\textbf{Conclusion.} $f$ décroît de $3$ vers $3-\mathrm{e}$ sur $]-\infty\,;1]$, puis croît vers $+\infty$ sur $[1\,;+\infty[$ ; la courbe admet l'asymptote $y = 3$ en $-\infty$, passe par $(0\,;1)$ et $(2\,;3)$, et possède un minimum en $(1\,;\,3-\mathrm{e})$.
\end{exoresolu}

\courssub{Méthode 7 : Modéliser une croissance ou décroissance continue}

\begin{methode}[Exploiter un modèle $f(t) = A\,\mathrm{e}^{kt}$ en situation concrète]
\begin{enumerate}
\item \textbf{Identifier les paramètres} : $A = f(0)$ est la valeur initiale ; $k > 0$ traduit une croissance, $k < 0$ une décroissance.
\item \textbf{Calculer une valeur} à l'instant $t$ : remplacer et utiliser la calculatrice (touches $\mathrm{e}^x$).
\item \textbf{Résoudre $f(t) = c$} : isoler l'exponentielle, passer au logarithme népérien : $\mathrm{e}^{kt} = \dfrac{c}{A} \iff kt = \ln\left(\dfrac{c}{A}\right) \iff t = \dfrac{1}{k}\ln\left(\dfrac{c}{A}\right)$.
\item \textbf{Conclure en langage concret} : donner l'unité (heures, jours, FCFA, degrés...) et une phrase d'interprétation.
\end{enumerate}
\end{methode}

\begin{exoresolu}[Refroidissement d'une pièce mécanique]
\textbf{Énoncé.} Après chauffage, une pièce métallique est sortie d'un four. Sa température (en degrés Celsius) est modélisée par $\theta(t) = 25 + 475\,\mathrm{e}^{-0{,}1t}$, où $t$ est le temps en minutes ($t \geq 0$).
\begin{enumerate}
\item Quelle est la température de la pièce à la sortie du four ? Vers quelle température tend-elle lorsque $t$ devient grand ? Interpréter.
\item Un technicien peut manipuler la pièce lorsque sa température descend à $\SI{60}{\celsius}$. Déterminer au bout de combien de minutes cela se produit (arrondir à la minute).
\end{enumerate}
\tcblower
\textbf{Solution détaillée.}

\textbf{1.} À la sortie du four, $t = 0$ :
\[ \theta(0) = 25 + 475\,\mathrm{e}^{0} = 25 + 475 \times 1 = 500 \]
La pièce sort à $\SI{500}{\celsius}$.

Lorsque $t \to +\infty$ : $\lim\limits_{t \to +\infty} (-0{,}1t) = -\infty$, donc par composition $\lim\limits_{t \to +\infty} \mathrm{e}^{-0{,}1t} = 0$. Ainsi :
\[ \lim_{t \to +\infty} \theta(t) = 25 + 475 \times 0 = 25 \]
\emph{Interprétation :} la pièce refroidit et sa température se rapproche de $\SI{25}{\celsius}$, la température ambiante de l'atelier, sans jamais descendre en dessous.

\textbf{2.} On résout $\theta(t) = 60$ :
\begin{align*}
25 + 475\,\mathrm{e}^{-0{,}1t} = 60 &\iff 475\,\mathrm{e}^{-0{,}1t} = 35 \\
&\iff \mathrm{e}^{-0{,}1t} = \dfrac{35}{475} = \dfrac{7}{95} \\
&\iff -0{,}1t = \ln\left(\dfrac{7}{95}\right) \\
&\iff t = \dfrac{\ln\left(\dfrac{7}{95}\right)}{-0{,}1} = -10\ln\left(\dfrac{7}{95}\right) = 10\ln\left(\dfrac{95}{7}\right)
\end{align*}
Numériquement : $\dfrac{95}{7} \approx 13{,}57$, $\ln(13{,}57) \approx 2{,}608$, donc $t \approx 26{,}08$ minutes.

\textbf{Conclusion.} La pièce sort à $\SI{500}{\celsius}$ et tend vers la température ambiante de $\SI{25}{\celsius}$ ; le technicien pourra la manipuler au bout d'environ \textbf{27 minutes} (on arrondit à la minute supérieure pour garantir la sécurité : à $26$ min, la température est encore légèrement supérieure à $\SI{60}{\celsius}$).
\end{exoresolu}

\begin{savaistu}[Éuler, complexes et suites : regards croisés]
\begin{itemize}
\item \textbf{Leonhard Euler (1707--1783)} a introduit la notation $\mathrm{e}$ pour la base des logarithmes népériens (vers 1727--1728). Il a démontré que $\mathrm{e} = \sum_{n=0}^{+\infty} \frac{1}{n!}$ et établi la formule célèbre $\mathrm{e}^{i\pi} + 1 = 0$, liant les cinq constantes fondamentales : $0, 1, \mathrm{e}, i, \pi$.
\item \textbf{Nombres complexes :} la fonction exponentielle s'étend à $\mathbb{C}$ par $\mathrm{e}^{x+iy} = \mathrm{e}^x(\cos y + i\sin y)$. La formule d'Euler $\mathrm{e}^{i\theta} = \cos\theta + i\sin\theta$ permet de passer de la forme algébrique à la forme exponentielle (trigonométrique) d'un complexe, simplifiant produits, puissances et racines.
\item \textbf{Suites :} la suite $\left(1+\frac{1}{n}\right)^n$ converge vers $\mathrm{e}$. Plus généralement, $\left(1+\frac{x}{n}\right)^n \xrightarrow[n\to+\infty]{} \mathrm{e}^x$. C'est une porte d'entrée vers la définition de $\exp$ comme limite de suites, utile en analyse numérique (calcul approché de $\mathrm{e}^x$).
\item \textbf{Applications techniques :} transformée de Laplace (résolution d'équations différentielles en automatique), transformée de Fourier (traitement du signal, électricité), modélisation de la décharge d'un condensateur ($u(t) = U_0 \mathrm{e}^{-t/RC}$).
\end{itemize}
\end{savaistu}

\begin{attention}[Les 5 erreurs qui coûtent des points au Probatoire]
\begin{enumerate}
\item \textbf{Inventer une formule fausse :} $\mathrm{e}^{a} + \mathrm{e}^{b} \neq \mathrm{e}^{a+b}$ et $\mathrm{e}^{a} \times \mathrm{e}^{b} \neq \mathrm{e}^{ab}$. Seules les formules $\mathrm{e}^{a}\mathrm{e}^{b} = \mathrm{e}^{a+b}$ et $\left(\mathrm{e}^{a}\right)^n = \mathrm{e}^{na}$ sont correctes.
\item \textbf{Oublier la dérivée de l'exposant :} la dérivée de $\mathrm{e}^{3x+2}$ n'est \emph{pas} $\mathrm{e}^{3x+2}$ mais $3\,\mathrm{e}^{3x+2}$. Symétriquement, une primitive de $\mathrm{e}^{3x+2}$ est $\dfrac{1}{3}\mathrm{e}^{3x+2}$, pas $\mathrm{e}^{3x+2}$.
\item \textbf{Changer le sens d'une inéquation :} la fonction exponentielle est \emph{strictement croissante}, donc $\mathrm{e}^{A} \leq \mathrm{e}^{B} \iff A \leq B$ : on ne change \textbf{jamais} le sens de l'inégalité en « descendant » l'exposant.
\item \textbf{Chercher des solutions impossibles :} $\mathrm{e}^{x} = 0$, $\mathrm{e}^{x} = -3$ n'ont aucune solution car $\mathrm{e}^{x} > 0$ pour tout réel $x$. Le dire clairement dans la rédaction rapporte des points.
\item \textbf{Bâcler la rédaction des limites et de la conclusion :} annoncer la forme indéterminée, factoriser par $\mathrm{e}^{x}$ avant de conclure, et toujours terminer par une phrase (« l'ensemble des solutions est... », « la droite $y = 3$ est asymptote... », « la machine produit environ 137 pièces... »). Une réponse sans justification ni unité perd des points même si le calcul est juste.
\end{enumerate}
\end{attention}

\newpage

\coursec{Exercices}

\coursec{Fonction exponentielle népérienne}

\courssub{Je vérifie mes connaissances}

\begin{exercice}[QCM : propriétés de l'exponentielle]
Pour chaque question, une seule réponse est exacte. Entourer la bonne réponse.
\begin{enumerate}
\item Pour tous réels $a$ et $b$, $\mathrm{e}^{a+b}$ est égal à :
\begin{itemize}
\item[a)] $\mathrm{e}^{a}+\mathrm{e}^{b}$ \qquad b) $\mathrm{e}^{a}\times \mathrm{e}^{b}$ \qquad c) $\mathrm{e}^{ab}$ \qquad d) $(\mathrm{e}^{a})^{b}$
\end{itemize}
\item $\dfrac{\mathrm{e}^{5}}{\mathrm{e}^{2}}$ est égal à :
\begin{itemize}
\item[a)] $\mathrm{e}^{3}$ \qquad b) $\mathrm{e}^{2,5}$ \qquad c) $\mathrm{e}^{7}$ \qquad d) $\mathrm{e}^{10}$
\end{itemize}
\item La dérivée de la fonction $x\longmapsto \mathrm{e}^{3x}$ est :
\begin{itemize}
\item[a)] $\mathrm{e}^{3x}$ \qquad b) $3\mathrm{e}^{x}$ \qquad c) $3\mathrm{e}^{3x}$ \qquad d) $\mathrm{e}^{3}$
\end{itemize}
\item $\displaystyle\lim_{x\to -\infty}\mathrm{e}^{x}$ est égal à :
\begin{itemize}
\item[a)] $+\infty$ \qquad b) $-\infty$ \qquad c) $1$ \qquad d) $0$
\end{itemize}
\end{enumerate}
\end{exercice}

\begin{exercice}[Vrai ou faux, en justifiant]
Répondre par \textbf{vrai} ou \textbf{faux} en justifiant chaque réponse.
\begin{enumerate}
\item Pour tout réel $x$, $\mathrm{e}^{x}>0$.
\item $\mathrm{e}^{0}=0$.
\item Pour tout réel $x$, $\mathrm{e}^{-x}=\dfrac{1}{\mathrm{e}^{x}}$.
\item L'équation $\mathrm{e}^{x}=-1$ admet une solution dans $\mathbb{R}$.
\item La fonction exponentielle est strictement croissante sur $\mathbb{R}$.
\end{enumerate}
\end{exercice}

\begin{exercice}[Calculs directs]
Simplifier les expressions suivantes, où $x$ désigne un réel quelconque :
\begin{enumerate}
\item $A=\mathrm{e}^{2}\times \mathrm{e}^{3}\times \mathrm{e}^{-4}$ ;
\item $B=\dfrac{\mathrm{e}^{x+1}}{\mathrm{e}^{x-2}}$ ;
\item $C=\left(\mathrm{e}^{x}\right)^{2}\times \mathrm{e}^{-x}$ ;
\item $D=\mathrm{e}^{x}\left(\mathrm{e}^{x}+\mathrm{e}^{-x}\right)$.
\end{enumerate}
\end{exercice}

\begin{exercice}[Dérivées immédiates]
Déterminer la fonction dérivée de chacune des fonctions suivantes, définies sur $\mathbb{R}$ :
\begin{enumerate}
\item $f(x)=\mathrm{e}^{x}+3x-2$ ;
\item $g(x)=x\,\mathrm{e}^{x}$ ;
\item $h(x)=\mathrm{e}^{-2x}$ ;
\item $k(x)=\mathrm{e}^{x^{2}+1}$.
\end{enumerate}
\end{exercice}

\courssub{Je m'entraîne}

\begin{exercice}[Équations et inéquations exponentielles]
Résoudre dans $\mathbb{R}$ les équations et inéquations suivantes :
\begin{enumerate}
\item $\mathrm{e}^{x}=5$ ;
\item $\mathrm{e}^{2x-1}=\mathrm{e}^{x+3}$ ;
\item $\mathrm{e}^{x}\times \mathrm{e}^{2x}=\mathrm{e}^{6}$ ;
\item $\mathrm{e}^{x}>3$ ;
\item $\mathrm{e}^{2x}-3\mathrm{e}^{x}+2=0$ (on pourra poser $X=\mathrm{e}^{x}$).
\end{enumerate}
\end{exercice}

\begin{exercice}[Limites]
Déterminer les limites suivantes :
\begin{enumerate}
\item $\displaystyle\lim_{x\to +\infty}\left(\mathrm{e}^{x}-3x+1\right)$ ;
\item $\displaystyle\lim_{x\to -\infty}\left(x\,\mathrm{e}^{x}\right)$ ;
\item $\displaystyle\lim_{x\to +\infty}\dfrac{\mathrm{e}^{x}}{x}$ ;
\item $\displaystyle\lim_{x\to -\infty}\left(2x+1\right)\mathrm{e}^{x}$.
\end{enumerate}
\end{exercice}

\begin{exercice}[Primitives et intégrales]
\begin{enumerate}
\item Déterminer une primitive sur $\mathbb{R}$ de chacune des fonctions :
\[ f(x)=\mathrm{e}^{x}+2 \quad ; \quad g(x)=\mathrm{e}^{3x} \quad ; \quad h(x)=2x\,\mathrm{e}^{x^{2}}. \]
\item Calculer les intégrales :
\[ I=\int_{0}^{1}\mathrm{e}^{x}\,\mathrm{d}x \quad ; \quad J=\int_{0}^{2}\mathrm{e}^{-x}\,\mathrm{d}x \quad ; \quad K=\int_{0}^{1}x\,\mathrm{e}^{x^{2}}\,\mathrm{d}x. \]
\end{enumerate}
\end{exercice}

\begin{exercice}[Tangente à la courbe de l'exponentielle]
Soit $f$ la fonction définie sur $\mathbb{R}$ par $f(x)=\mathrm{e}^{x}$ et $\mathcal{C}_{f}$ sa courbe représentative dans un repère orthonormé.
\begin{enumerate}
\item Déterminer l'équation réduite de la tangente $T$ à $\mathcal{C}_{f}$ au point d'abscisse $0$.
\item Déterminer l'équation réduite de la tangente à $\mathcal{C}_{f}$ au point d'abscisse $1$.
\item Tracer $\mathcal{C}_{f}$ et la tangente $T$ dans un repère orthonormé (unité : 2 cm).
\end{enumerate}
\end{exercice}

\courssub{Je me prépare au Probatoire}

\begin{exercice}[Étude complète d'une fonction]
On considère la fonction $f$ définie sur $\mathbb{R}$ par :
\[ f(x)=(2x-3)\,\mathrm{e}^{x}. \]
On note $\mathcal{C}_{f}$ sa courbe représentative dans un repère orthonormé $(O\,;\vec{\imath},\vec{\jmath})$ (unité : 2 cm).
\begin{enumerate}
\item \begin{enumerate}
\item Calculer $\displaystyle\lim_{x\to +\infty}f(x)$.
\item Montrer que pour tout réel $x$, $f(x)=2x\,\mathrm{e}^{x}-3\mathrm{e}^{x}$. En déduire $\displaystyle\lim_{x\to -\infty}f(x)$. Interpréter graphiquement ce résultat.
\end{enumerate}
\item \begin{enumerate}
\item Montrer que pour tout réel $x$, $f'(x)=(2x-1)\,\mathrm{e}^{x}$.
\item Étudier le signe de $f'(x)$ et dresser le tableau de variations de $f$ sur $\mathbb{R}$.
\end{enumerate}
\item Déterminer l'équation réduite de la tangente $T$ à $\mathcal{C}_{f}$ au point d'abscisse $1$.
\item Montrer que l'équation $f(x)=-2\mathrm{e}^{1/2}$ admet une unique solution $\alpha$ dans $\mathbb{R}$, et donner un encadrement de $\alpha$ d'amplitude $0,1$.
\item Tracer $\mathcal{C}_{f}$ et $T$.
\end{enumerate}
\end{exercice}

\begin{exercice}[Concentration d'un médicament]
Après une injection, la concentration $C(t)$ (en mg/L) d'un médicament dans le sang d'un patient est modélisée, pour $t\geq 0$ (en heures), par :
\[ C(t)=50\,\mathrm{e}^{-0,2t}+10\,\mathrm{e}^{-0,5t}. \]
\begin{enumerate}
\item Calculer la concentration initiale $C(0)$.
\item \begin{enumerate}
\item Montrer que pour tout $t\geq 0$ : $C'(t)=-10\,\mathrm{e}^{-0,2t}-5\,\mathrm{e}^{-0,5t}$.
\item En déduire le sens de variation de $C$ sur $[0\,;+\infty[$.
\end{enumerate}
\item Déterminer $\displaystyle\lim_{t\to +\infty}C(t)$. Interpréter ce résultat dans le contexte de l'exercice.
\item Le médicament n'est plus efficace lorsque la concentration passe sous le seuil de $5$ mg/L.
\begin{enumerate}
\item Montrer que l'équation $C(t)=5$ admet une unique solution $t_{0}$ sur $[0\,;+\infty[$.
\item À l'aide de la calculatrice, donner un encadrement de $t_{0}$ d'amplitude $0,5$, puis conclure sur la durée d'efficacité du médicament.
\end{enumerate}
\end{enumerate}
\end{exercice}

\begin{exercice}[Aire sous une courbe et intégration par parties]
On considère la fonction $f$ définie sur $[0\,;2]$ par :
\[ f(x)=x^{2}\,\mathrm{e}^{-x}. \]
On note $\mathcal{C}_{f}$ sa courbe représentative dans un repère orthonormé (unité : 3 cm).
\begin{enumerate}
\item Calculer $f'(x)$ et dresser le tableau de variations de $f$ sur $[0\,;2]$.
\item Tracer $\mathcal{C}_{f}$ sur $[0\,;2]$.
\item On pose $I=\displaystyle\int_{0}^{2}x^{2}\mathrm{e}^{-x}\,\mathrm{d}x$.
\begin{enumerate}
\item À l'aide d'une intégration par parties, montrer que :
\[ I=\left[-x^{2}\mathrm{e}^{-x}\right]_{0}^{2}+2\int_{0}^{2}x\,\mathrm{e}^{-x}\,\mathrm{d}x. \]
\item À l'aide d'une seconde intégration par parties, calculer $\displaystyle\int_{0}^{2}x\,\mathrm{e}^{-x}\,\mathrm{d}x$.
\item En déduire la valeur exacte de $I$.
\end{enumerate}
\item En déduire l'aire, en unités d'aire puis en $\mathrm{cm}^{2}$, du domaine plan délimité par $\mathcal{C}_{f}$, l'axe des abscisses et les droites d'équations $x=0$ et $x=2$. Donner une valeur approchée à $10^{-2}$ près.
\end{enumerate}
\end{exercice}

\courssub{Compléments et ouverture}

\begin{savaistu}[Leonhard Euler et le nombre \cle{e}]
\emph{Leonhard Euler} (1707--1783), mathématicien suisse, a introduit la notation \cle{e} pour la base des logarithmes népériens (vers 1727--1728). Il a démontré que \cle{e} est irrationnel (1737) et a établi la formule célèbre reliant les cinq constantes fondamentales des mathématiques :
\[ \mathrm{e}^{\mathrm{i}\pi} + 1 = 0. \]
Cette formule, appelée \emph{identité d'Euler}, lie le nombre \cle{e}, l'unité imaginaire \cle{i}, le nombre \cle{\pi}, l'unité \cle{1} et le zéro \cle{0}.
En Terminale technique, on rencontre \cle{e} comme limite de la suite $\left(1+\frac{1}{n}\right)^n$ (croissance composée continue) et comme base de la fonction exponentielle dont la dérivée est elle-même : $(\mathrm{e}^x)' = \mathrm{e}^x$.
\end{savaistu}

\begin{savaistu}[Exponentielle complexe et suites]
L'extension de l'exponentielle aux nombres complexes s'écrit $\mathrm{e}^{z} = \mathrm{e}^{x}(\cos y + \mathrm{i}\sin y)$ pour $z=x+\mathrm{i}y$.
Cela permet de définir le cosinus et le sinus via les exponentielles :
\[ \cos x = \frac{\mathrm{e}^{\mathrm{i}x}+\mathrm{e}^{-\mathrm{i}x}}{2}, \quad \sin x = \frac{\mathrm{e}^{\mathrm{i}x}-\mathrm{e}^{-\mathrm{i}x}}{2\mathrm{i}}. \]
On retrouve ainsi la formule de Moivre : $(\cos x + \mathrm{i}\sin x)^n = \cos(nx) + \mathrm{i}\sin(nx)$.
Les suites géométriques de raison $\mathrm{e}^k$ modélisent des croissances/décroissances continues (radioactivité, démographie, finance : intérêts composés continus).
\end{savaistu}

\newpage

\coursec{Corrigés des exercices}

\begin{corrige}[Exercice 1 — QCM : propriétés de l'exponentielle]
\begin{enumerate}
\item \textbf{b)} $\mathrm{e}^{a+b} = \mathrm{e}^{a}\times \mathrm{e}^{b}$ (propriété fondamentale : l'exponentielle transforme la somme en produit).
\item \textbf{a)} $\dfrac{\mathrm{e}^{5}}{\mathrm{e}^{2}} = \mathrm{e}^{5-2} = \mathrm{e}^{3}$ (quotient de puissances même base).
\item \textbf{c)} Dérivée de $x\mapsto \mathrm{e}^{3x}$ : $3\mathrm{e}^{3x}$ (règle de la chaîne : $(u)' u'$ avec $u=3x$).
\item \textbf{d)} $\displaystyle\lim_{x\to -\infty}\mathrm{e}^{x} = 0$ (asymptote horizontale $y=0$ en $-\infty$).
\end{enumerate}
\end{corrige}

\begin{corrige}[Exercice 2 — Vrai ou faux, en justifiant]
\begin{enumerate}
\item \textbf{Vrai.} Pour tout réel $x$, $\mathrm{e}^{x} > 0$. L'ensemble image de la fonction exponentielle est $]0;+\infty[$.
\item \textbf{Faux.} $\mathrm{e}^{0} = 1$ (par définition, $\mathrm{e}^{0}=1$ pour toute exponentielle de base non nulle).
\item \textbf{Vrai.} $\mathrm{e}^{-x} = \dfrac{1}{\mathrm{e}^{x}}$ (propriété $\mathrm{e}^{-a} = 1/\mathrm{e}^{a}$ ou $a=-x$).
\item \textbf{Faux.} L'équation $\mathrm{e}^{x}=-1$ n'admet \textbf{aucune} solution dans $\mathbb{R}$ car $\mathrm{e}^{x}>0$ pour tout $x\in\mathbb{R}$.
\item \textbf{Vrai.} La fonction exponentielle est strictement croissante sur $\mathbb{R}$ car sa dérivée $\mathrm{e}^{x}$ est strictement positive sur $\mathbb{R}$.
\end{enumerate}
\end{corrige}

\begin{corrige}[Exercice 3 — Calculs directs]
\begin{enumerate}
\item $A = \mathrm{e}^{2}\times \mathrm{e}^{3}\times \mathrm{e}^{-4} = \mathrm{e}^{2+3-4} = \mathrm{e}^{1} = \boxed{\mathrm{e}}$.
\item $B = \dfrac{\mathrm{e}^{x+1}}{\mathrm{e}^{x-2}} = \mathrm{e}^{(x+1)-(x-2)} = \mathrm{e}^{3} = \boxed{\mathrm{e}^{3}}$.
\item $C = \left(\mathrm{e}^{x}\right)^{2}\times \mathrm{e}^{-x} = \mathrm{e}^{2x}\times \mathrm{e}^{-x} = \mathrm{e}^{2x-x} = \boxed{\mathrm{e}^{x}}$.
\item $D = \mathrm{e}^{x}\left(\mathrm{e}^{x}+\mathrm{e}^{-x}\right) = \mathrm{e}^{x}\times\mathrm{e}^{x} + \mathrm{e}^{x}\times\mathrm{e}^{-x} = \mathrm{e}^{2x} + \mathrm{e}^{0} = \boxed{\mathrm{e}^{2x} + 1}$.
\end{enumerate}
\end{corrige}

\begin{corrige}[Exercice 4 — Dérivées immédiates]
\begin{enumerate}
\item $f(x)=\mathrm{e}^{x}+3x-2 \quad\Rightarrow\quad f'(x) = \mathrm{e}^{x} + 3$.
\item $g(x)=x\,\mathrm{e}^{x}$ (produit) $\quad\Rightarrow\quad g'(x) = 1\cdot\mathrm{e}^{x} + x\cdot\mathrm{e}^{x} = (x+1)\mathrm{e}^{x}$.
\item $h(x)=\mathrm{e}^{-2x}$ (composée) $\quad\Rightarrow\quad h'(x) = (-2)\mathrm{e}^{-2x} = -2\mathrm{e}^{-2x}$.
\item $k(x)=\mathrm{e}^{x^{2}+1}$ (composée) $\quad\Rightarrow\quad k'(x) = (2x)\mathrm{e}^{x^{2}+1} = 2x\,\mathrm{e}^{x^{2}+1}$.
\end{enumerate}
\end{corrige}

\begin{corrige}[Exercice 5 — Équations et inéquations exponentielles]
\begin{enumerate}
\item $\mathrm{e}^{x}=5 \;\Leftrightarrow\; x = \ln 5$. \quad Solution : $\boxed{\ln 5}$.
\item $\mathrm{e}^{2x-1}=\mathrm{e}^{x+3} \;\Leftrightarrow\; 2x-1 = x+3 \;\Leftrightarrow\; x = 4$. \quad Solution : $\boxed{4}$.
\item $\mathrm{e}^{x}\times \mathrm{e}^{2x}=\mathrm{e}^{6} \;\Leftrightarrow\; \mathrm{e}^{3x}=\mathrm{e}^{6} \;\Leftrightarrow\; 3x=6 \;\Leftrightarrow\; x=2$. \quad Solution : $\boxed{2}$.
\item $\mathrm{e}^{x}>3 \;\Leftrightarrow\; x > \ln 3$. \quad Ensemble solution : $\boxed{]\ln 3; +\infty[}$.
\item Posons $X=\mathrm{e}^{x} > 0$. L'équation devient $X^{2}-3X+2=0 \;\Leftrightarrow\; (X-1)(X-2)=0$.
  $X=1$ ou $X=2$.
  $\mathrm{e}^{x}=1 \Rightarrow x=0$ ; $\mathrm{e}^{x}=2 \Rightarrow x=\ln 2$.
  Solutions : $\boxed{\{0;\ \ln 2\}}$.
\end{enumerate}
\end{corrige}

\begin{corrige}[Exercice 6 — Limites]
\begin{enumerate}
\item $\displaystyle\lim_{x\to +\infty}\left(\mathrm{e}^{x}-3x+1\right)$.
  $\mathrm{e}^{x} \to +\infty$, $-3x \to -\infty$, mais $\mathrm{e}^{x}$ domine tout polynôme.
  $\mathrm{e}^{x}-3x+1 = \mathrm{e}^{x}\left(1 - \frac{3x}{\mathrm{e}^{x}} + \frac{1}{\mathrm{e}^{x}}\right) \to +\infty \times 1 = \boxed{+\infty}$.
\item $\displaystyle\lim_{x\to -\infty}\left(x\,\mathrm{e}^{x}\right)$.
  Forme indéterminée $(-\infty)\times 0$. L'exponentielle tend vers $0$ plus vite que $x$ ne tend vers $-\infty$.
  $\boxed{0}$.
\item $\displaystyle\lim_{x\to +\infty}\dfrac{\mathrm{e}^{x}}{x}$.
  Forme $\frac{+\infty}{+\infty}$. $\mathrm{e}^{x}$ croît plus vite que $x$.
  $\boxed{+\infty}$.
\item $\displaystyle\lim_{x\to -\infty}\left(2x+1\right)\mathrm{e}^{x}$.
  Même raisonnement que 2. : polynôme $\times$ exponentielle $\to 0$.
  $\boxed{0}$.
\end{enumerate}
\end{corrige}

\begin{corrige}[Exercice 7 — Primitives et intégrales]
\begin{enumerate}
\item \textbf{Primitives sur $\mathbb{R}$ (on ajoute $+C$) :}
  \begin{itemize}
  \item $f(x)=\mathrm{e}^{x}+2 \quad\Rightarrow\quad F(x) = \mathrm{e}^{x} + 2x + C$.
  \item $g(x)=\mathrm{e}^{3x} \quad\Rightarrow\quad G(x) = \dfrac{1}{3}\mathrm{e}^{3x} + C$ (vérification : $G'(x)=\frac{1}{3}\cdot 3\mathrm{e}^{3x}=\mathrm{e}^{3x}$).
  \item $h(x)=2x\,\mathrm{e}^{x^{2}}$. On reconnaît la forme $u'(x)\mathrm{e}^{u(x)}$ avec $u(x)=x^{2}$, $u'(x)=2x$.
        $\quad\Rightarrow\quad H(x) = \mathrm{e}^{x^{2}} + C$.
  \end{itemize}
\item \textbf{Intégrales :}
  \begin{align*}
  I &= \int_{0}^{1}\mathrm{e}^{x}\,\mathrm{d}x = \Bigl[\mathrm{e}^{x}\Bigr]_{0}^{1} = \mathrm{e}^{1} - \mathrm{e}^{0} = \boxed{\mathrm{e} - 1}. \\[4pt]
  J &= \int_{0}^{2}\mathrm{e}^{-x}\,\mathrm{d}x = \Bigl[-\mathrm{e}^{-x}\Bigr]_{0}^{2} = -\mathrm{e}^{-2} - (-\mathrm{e}^{0}) = 1 - \mathrm{e}^{-2} = \boxed{1 - \frac{1}{\mathrm{e}^{2}}}. \\[4pt]
  K &= \int_{0}^{1}x\,\mathrm{e}^{x^{2}}\,\mathrm{d}x.
      \text{Posons } u = x^{2},\; \mathrm{d}u = 2x\,\mathrm{d}x \;\Rightarrow\; x\,\mathrm{d}x = \frac{1}{2}\mathrm{d}u.
      \text{Pour } x=0, u=0;\; x=1, u=1. \\
      K &= \frac{1}{2}\int_{0}^{1}\mathrm{e}^{u}\,\mathrm{d}u = \frac{1}{2}\Bigl[\mathrm{e}^{u}\Bigr]_{0}^{1} = \frac{1}{2}(\mathrm{e} - 1) = \boxed{\frac{\mathrm{e}-1}{2}}.
  \end{align*}
\end{enumerate}
\end{corrige}

\begin{corrige}[Exercice 8 — Tangente à la courbe de l'exponentielle]
$f(x)=\mathrm{e}^{x}$, $f'(x)=\mathrm{e}^{x}$.
\begin{enumerate}
\item En $x=0$ : $f(0)=1$, $f'(0)=1$.
  Équation de la tangente $T$ : $y = f(0) + f'(0)(x-0) = 1 + x$.
  $\boxed{y = x + 1}$.
\item En $x=1$ : $f(1)=\mathrm{e}$, $f'(1)=\mathrm{e}$.
  Équation de la tangente : $y = \mathrm{e} + \mathrm{e}(x-1) = \mathrm{e}x$.
  $\boxed{y = \mathrm{e}x}$.
\item \textbf{Trace} (unité 2 cm) :
\begin{popfigure}
\begin{tikzpicture}[scale=0.7, domain=-2:2]
\draw[popline, thin, ->] (-2.5,0) -- (2.5,0) node[right] {$x$};
\draw[popline, thin, ->] (0,-1) -- (0,5.5) node[above] {$y$};
\foreach \x in {-2,-1,1,2} \draw (\x,0.1) -- (\x,-0.1) node[below] {\small $\x$};
\foreach \y in {1,2,3,4,5} \draw (0.1,\y) -- (-0.1,\y) node[left] {\small $\y$};
\draw[popcurve, thick, popblue, samples=200] plot (\x, {exp(\x)});
\draw[popcurve, thick, poporange, dashed] plot (\x, {\x+1});
\draw[popmuted, dashed] (0,1) node[left] {\small $1$} -- (0,0);
\node[popblue] at (1.5, 4) {$\mathcal{C}_f$};
\node[poporange] at (-1.5, 0.5) {$T$};
\end{tikzpicture}
\legende{Courbe $y=\mathrm{e}^x$ et tangente $y=x+1$ au point $(0,1)$.}
\end{popfigure}
\end{enumerate}
\end{corrige}

\begin{corrige}[Exercice 9 — Étude complète d'une fonction]
\begin{enumerate}
\item \begin{enumerate}
\item $\lim\limits_{x\to +\infty}(2x-3)=+\infty$ et $\lim\limits_{x\to +\infty}\mathrm{e}^{x}=+\infty$, donc $\lim\limits_{x\to +\infty}f(x)=+\infty$.
\item $f(x)=(2x-3)\mathrm{e}^{x}=2x\mathrm{e}^{x}-3\mathrm{e}^{x}$ (distributivité).
  $\lim\limits_{x\to -\infty}\mathrm{e}^{x}=0$ et $\lim\limits_{x\to -\infty}x\mathrm{e}^{x}=0$ (l'exponentielle domine le polynôme).
  Donc $\lim\limits_{x\to -\infty}f(x)=0$.
  Graphiquement, l'axe des abscisses (droite d'équation $y=0$) est une asymptote horizontale à la courbe $\mathcal{C}_{f}$ pour $x\to -\infty$.
\end{enumerate}
\item \begin{enumerate}
\item $f'(x)=2\mathrm{e}^{x}+(2x-3)\mathrm{e}^{x}=(2x-1)\mathrm{e}^{x}$ (dérivée du produit).
\item $\mathrm{e}^{x}>0$ sur $\mathbb{R}$. Le signe de $f'(x)$ est celui de $2x-1$.
  $f'(x)=0 \Leftrightarrow x=\frac{1}{2}$.
  $f'(x)<0$ sur $]-\infty;\frac{1}{2}[$ et $f'(x)>0$ sur $]\frac{1}{2};+\infty[$.
  \begin{center}
  \begin{tabular}{|c|ccccc|}\hline
  $x$ & $-\infty$ & & $\frac{1}{2}$ & & $+\infty$ \\ \hline
  $f'(x)$ & & $-$ & $0$ & $+$ & \\ \hline
  $f$ & $0$ & $\searrow$ & $-2\mathrm{e}^{1/2}$ & $\nearrow$ & $+\infty$ \\ \hline
  \end{tabular}
  \end{center}
\end{enumerate}
\item $f(1)=-\mathrm{e}$, $f'(1)=\mathrm{e}$.
  Équation de $T$ : $y = -\mathrm{e} + \mathrm{e}(x-1)$, soit $\boxed{y = \mathrm{e}x - 2\mathrm{e}}$.
\item Le minimum de $f$ sur $\mathbb{R}$ est $f(1/2)=-2\mathrm{e}^{1/2}$.
  L'équation $f(x)=-2\mathrm{e}^{1/2}$ revient à $f(x)=f(1/2)$.
  Comme $f$ est strictement décroissante sur $]-\infty;1/2]$ et strictement croissante sur $[1/2;+\infty[$, l'unique solution est $\alpha=\frac{1}{2}$.
  Encadrement d'amplitude $0,1$ : $\alpha \in [0{,}45\,;\,0{,}55]$.
\item \textbf{Trace} : voir figure ci-dessous.
\end{enumerate}

\begin{popfigure}
\begin{tikzpicture}[scale=0.7, domain=-2.5:2]
\draw[popline, thin, ->] (-3,0) -- (2.5,0) node[right] {$x$};
\draw[popline, thin, ->] (0,-4) -- (0,5) node[above] {$y$};
\foreach \x in {-2,-1,1,2} \draw (\x,0.1) -- (\x,-0.1) node[below] {\small $\x$};
\foreach \y in {-3,-2,-1,1,2,3,4} \draw (0.1,\y) -- (-0.1,\y) node[left] {\small $\y$};
\draw[popcurve, thick, popblue, samples=200] plot (\x, {(2*\x-3)*exp(\x)});
\draw[popcurve, thick, poporange, dashed] plot (\x, {exp(1)*\x - 2*exp(1)});
\draw[popmuted, dashed] (0.5, -3.3) -- (0.5, 0) node[below] {\small $\frac{1}{2}$};
\draw[popmuted, dashed] (0, -3.297) -- (0.5, -3.297) node[left] {\small $-2\mathrm{e}^{1/2}$};
\node[popblue] at (1.8, 3.5) {$\mathcal{C}_f$};
\node[poporange] at (1.8, 1.5) {$T$};
\end{tikzpicture}
\legende{Courbe $\mathcal{C}_f$ de $f(x)=(2x-3)\mathrm{e}^x$ et tangente $T$ en $x=1$.}
\end{popfigure}

\textbf{Barème indicatif (sur 20) :} Limites (3), Développement et limite $-\infty$ (2), Dérivée (3), Tableau de variations (4), Tangente (2), Équation/Encadrement (4), Tracé (2).
\textbf{Points de vigilance :} Justifier $\lim x\mathrm{e}^x=0$ (croissance exponentielle), signe de $\mathrm{e}^x>0$, écriture du tableau de variations (3 lignes, même nb colonnes), forme réduite de la tangente, bijection pour l'unicité de $\alpha$.
\end{corrige}

\begin{corrige}[Exercice 10 — Concentration d'un médicament]
\begin{enumerate}
\item $C(0)=50\mathrm{e}^{0}+10\mathrm{e}^{0}=60$ mg/L.
\item \begin{enumerate}
\item $C'(t)=50(-0{,}2)\mathrm{e}^{-0{,}2t}+10(-0{,}5)\mathrm{e}^{-0{,}5t}=-10\mathrm{e}^{-0{,}2t}-5\mathrm{e}^{-0{,}5t}$.
\item Les exponentielles sont strictement positives, donc $C'(t)<0$ pour tout $t\ge 0$.
  $C$ est strictement décroissante sur $[0;+\infty[$.
\end{enumerate}
\item $\lim\limits_{t\to +\infty}\mathrm{e}^{-0{,}2t}=0$ et $\lim\limits_{t\to +\infty}\mathrm{e}^{-0{,}5t}=0$, donc $\lim\limits_{t\to +\infty}C(t)=0$.
  Interprétation : la concentration tend vers $0$ mg/L, le médicament est totalement éliminé de l'organisme à long terme.
\item \begin{enumerate}
\item $C$ est continue et strictement décroissante sur $[0;+\infty[$.
  $C(0)=60 > 5$ et $\lim\limits_{t\to +\infty}C(t)=0 < 5$.
  D'après le théorème des valeurs intermédiaires (bijection), l'équation $C(t)=5$ admet une unique solution $t_0$ sur $[0;+\infty[$.
\item À la calculatrice (tableau de valeurs ou solveur) :
  $C(10) \approx 50\mathrm{e}^{-2}+10\mathrm{e}^{-5} \approx 6,77+0,07=6,84 > 5$.
  $C(11) \approx 50\mathrm{e}^{-2,2}+10\mathrm{e}^{-5,5} \approx 5,54+0,04=5,58 > 5$.
  $C(12) \approx 50\mathrm{e}^{-2,4}+10\mathrm{e}^{-6} \approx 4,54+0,02=4,56 < 5$.
  $t_0 \in [11; 11{,}5]$ (amplitude $0,5$).
  Le médicament reste efficace pendant environ $11$ heures.
\end{enumerate}
\end{enumerate}

\begin{popfigure}
\begin{tikzpicture}[scale=0.6, domain=0:15]
\draw[popline, thin, ->] (0,0) -- (16,0) node[right] {$t$ (h)};
\draw[popline, thin, ->] (0,0) -- (0,65) node[above] {$C(t)$ (mg/L)};
\foreach \x in {5,10,15} \draw (\x,0.5) -- (\x,-0.5) node[below] {\small $\x$};
\foreach \y in {10,20,30,40,50,60} \draw (0.3,\y) -- (-0.3,\y) node[left] {\small $\y$};
\draw[popcurve, thick, popblue, samples=200] plot (\x, {50*exp(-0.2*\x)+10*exp(-0.5*\x)});
\draw[poporange, dashed] (0,5) -- (15,5) node[right] {Seuil $5$ mg/L};
\draw[popmuted, dashed] (11.2,0) -- (11.2,5) node[below] {\small $t_0\approx 11{,}2$};
\node[popblue] at (12, 40) {$C(t)$};
\end{tikzpicture}
\legende{Évolution de la concentration du médicament.}
\end{popfigure}

\textbf{Barème indicatif (sur 20) :} $C(0)$ (1), Dérivée (2), Variations (2), Limite + interprétation (3), TVI + unicité (4), Encadrement calculatrice (4), Tracé (4).
\textbf{Points de vigilance :} Dérivée de $\mathrm{e}^{kt}$ ($k\mathrm{e}^{kt}$), signe de la dérivée (négatif), conditions du TVI (continuité, monotonie, encadrement), lecture calculatrice (tableau de valeurs), unité du résultat (heures).
\end{corrige}

\begin{corrige}[Exercice 11 — Aire sous une courbe et intégration par parties]
\begin{enumerate}
\item $f'(x)=2x\mathrm{e}^{-x}-x^{2}\mathrm{e}^{-x}=x(2-x)\mathrm{e}^{-x}$.
  Sur $[0;2]$, $\mathrm{e}^{-x}>0$, $x\ge 0$, $2-x\ge 0$, donc $f'(x)\ge 0$.
  $f$ est croissante sur $[0;2]$.
  \begin{center}
  \begin{tabular}{|c|ccc|}\hline
  $x$ & $0$ & & $2$ \\ \hline
  $f'(x)$ & & $+$ & \\ \hline
  $f$ & $0$ & $\nearrow$ & $4\mathrm{e}^{-2}$ \\ \hline
  \end{tabular}
  \end{center}
\item \textbf{Trace} : voir figure ci-dessous.
\item \begin{enumerate}
\item IPP : $u=x^2 \Rightarrow u'=2x$ ; $v'=\mathrm{e}^{-x} \Rightarrow v=-\mathrm{e}^{-x}$.
  $I = \left[-x^2\mathrm{e}^{-x}\right]_0^2 - \int_0^2 (-\mathrm{e}^{-x})2x\,\mathrm{d}x = \left[-x^2\mathrm{e}^{-x}\right]_0^2 + 2\int_0^2 x\mathrm{e}^{-x}\,\mathrm{d}x$.
\item $J = \int_0^2 x\mathrm{e}^{-x}\,\mathrm{d}x$.
  IPP : $u=x \Rightarrow u'=1$ ; $v'=\mathrm{e}^{-x} \Rightarrow v=-\mathrm{e}^{-x}$.
  $J = \left[-x\mathrm{e}^{-x}\right]_0^2 - \int_0^2 (-\mathrm{e}^{-x})\,\mathrm{d}x = -2\mathrm{e}^{-2} + \left[-\mathrm{e}^{-x}\right]_0^2 = -2\mathrm{e}^{-2} -\mathrm{e}^{-2} + 1 = 1 - 3\mathrm{e}^{-2}$.
\item $I = \left(-4\mathrm{e}^{-2} - 0\right) + 2\left(1 - 3\mathrm{e}^{-2}\right) = -4\mathrm{e}^{-2} + 2 - 6\mathrm{e}^{-2} = \boxed{2 - 10\mathrm{e}^{-2}}$.
\end{enumerate}
\item Aire $\mathcal{A} = I = 2 - \frac{10}{\mathrm{e}^2}$ unités d'aire.
  Unité graphique : $3$ cm $\Rightarrow$ $1$ u.a. $= 9$ cm$^2$.
  $\mathcal{A}_{\text{cm}^2} = 9\left(2 - \frac{10}{\mathrm{e}^2}\right) \approx 9(2 - 1{,}353) = 9 \times 0{,}647 = \boxed{5{,}82 \text{ cm}^2}$ (au $10^{-2}$ près).
\end{enumerate}

\begin{popfigure}
\begin{tikzpicture}[scale=1.2, domain=0:2]
\draw[popline, thin, ->] (0,0) -- (2.3,0) node[right] {$x$};
\draw[popline, thin, ->] (0,0) -- (0,1.6) node[above] {$y$};
\foreach \x in {0.5,1,1.5,2} \draw (\x,0.03) -- (\x,-0.03) node[below] {\small $\x$};
\foreach \y in {0.5,1,1.5} \draw (0.03,\y) -- (-0.03,\y) node[left] {\small $\y$};
\draw[popcurve, thick, popblue, samples=200] plot (\x, {\x*\x*exp(-\x)});
\draw[popblue!30, fill=popblue!10] (0,0) -- plot[domain=0:2, samples=100] (\x, {\x*\x*exp(-\x)}) -- (2,0) -- cycle;
\node[popblue] at (1.5, 1.2) {$\mathcal{C}_f$};
\end{tikzpicture}
\legende{Courbe $f(x)=x^2\mathrm{e}^{-x}$ et aire $\mathcal{A}$ sur $[0;2]$.}
\end{popfigure}

\textbf{Barème indicatif (sur 20) :} Dérivée + variations (4), IPP 1 (3), IPP 2 (3), Valeur exacte $I$ (2), Aire u.a. + cm$^2$ (3), Tracé + hachure (5).
\textbf{Points de vigilance :} Choix judicieux $u=x^2$ puis $u=x$ pour IPP, signes lors de l'IPP ($v=-\mathrm{e}^{-x}$), calcul exact sans approximation prématurée, conversion u.a. $\to$ cm$^2$ (échelle 3 cm $\Rightarrow$ facteur 9).
\end{corrige}

\newpage

\coursec{Fiche bilan}

\begin{popfigure}
\begin{tikzpicture}[mindmap, grow cyclic, every node/.style={concept, font=\small},
root concept/.append style={concept color=popblue, font=\bfseries},
level 1/.append style={level distance=3.4cm, sibling angle=51.4},
level 1 concept/.append style={font=\scriptsize\bfseries},
level 2/.append style={level distance=2.1cm, sibling angle=45},
level 2 concept/.append style={font=\tiny}]
\node[root concept]{Fonction exponentielle $\mathrm{e}^x$}
child[concept color=poporange]{ node{Définition et propriétés}
  child{ node{$\mathrm{e}^{a+b}=\mathrm{e}^a\mathrm{e}^b$} }
  child{ node{$\mathrm{e}^0=1$ ; $\mathrm{e}^x>0$} }
}
child[concept color=popgreen]{ node{Limites}
  child{ node{$\lim_{+\infty}\mathrm{e}^x=+\infty$} }
  child{ node{$\lim_{-\infty}\mathrm{e}^x=0$} }
}
child[concept color=poppurple]{ node{Dérivation}
  child{ node{$(\mathrm{e}^x)'=\mathrm{e}^x$} }
  child{ node{$(\mathrm{e}^u)'=u'\mathrm{e}^u$} }
}
child[concept color=poppink]{ node{Primitives}
  child{ node{$\int \mathrm{e}^x\,\mathrm{d}x=\mathrm{e}^x$} }
  child{ node{$\int u'\mathrm{e}^u=\mathrm{e}^u$} }
}
child[concept color=popgold]{ node{Équations, inéquations}
  child{ node{$\mathrm{e}^a=\mathrm{e}^b \Leftrightarrow a=b$} }
  child{ node{$\mathrm{e}^x=k \Rightarrow x=\ln k$} }
}
child[concept color=popblue!60]{ node{Étude de fonctions}
  child{ node{Signe, variations, tangente} }
}
child[concept color=popgreen!70!popblue]{ node{Modélisation}
  child{ node{Croissance, décroissance continues} }
};
\end{tikzpicture}
\legende{Carte mentale de la leçon : la fonction exponentielle népérienne}
\end{popfigure}

\begin{bilan}
\textbf{Définitions clés}
\begin{itemize}
\item La fonction \cle{exponentielle népérienne}, notée $\exp$ ou $x\mapsto \mathrm{e}^x$, est l'unique fonction dérivable sur $\mathbb{R}$ telle que $f'=f$ et $f(0)=1$.
\item $\mathrm{e}=\mathrm{e}^1\approx \num{2,718}$ ; pour tout réel $x$ : $\mathrm{e}^x>0$.
\item La fonction $\exp$ est \textbf{strictement croissante} sur $\mathbb{R}$.
\end{itemize}

\textbf{Formules et théorèmes à connaître par cœur}
\begin{center}
\begin{tabularx}{\linewidth}{|l|X|}\hline
\rowcolor{popblueL} \textbf{Notion} & \textbf{Résultat} \\ \hline
Relation fonctionnelle & $\mathrm{e}^{a+b}=\mathrm{e}^a\times \mathrm{e}^b$ ; \quad $\mathrm{e}^{a-b}=\dfrac{\mathrm{e}^a}{\mathrm{e}^b}$ ; \quad $\mathrm{e}^{-a}=\dfrac{1}{\mathrm{e}^a}$ ; \quad $(\mathrm{e}^a)^n=\mathrm{e}^{na}$ \\ \hline
Limites & $\displaystyle\lim_{x\to +\infty}\mathrm{e}^x=+\infty$ ; \quad $\displaystyle\lim_{x\to -\infty}\mathrm{e}^x=0$ (asymptote $y=0$) ; \quad $\displaystyle\lim_{x\to 0}\frac{\mathrm{e}^x-1}{x}=1$ \\ \hline
Dérivées & $(\mathrm{e}^x)'=\mathrm{e}^x$ ; \quad $(\mathrm{e}^{u})'=u'\,\mathrm{e}^{u}$ ; \quad $(\mathrm{e}^{ax+b})'=a\,\mathrm{e}^{ax+b}$ \\ \hline
Primitives & $\displaystyle\int \mathrm{e}^x\,\mathrm{d}x=\mathrm{e}^x+C$ ; \quad $\displaystyle\int u'\,\mathrm{e}^{u}\,\mathrm{d}x=\mathrm{e}^{u}+C$ ; \quad $\displaystyle\int \mathrm{e}^{ax+b}\,\mathrm{d}x=\frac{1}{a}\mathrm{e}^{ax+b}+C$ \\ \hline
Équations / inéquations & $\mathrm{e}^a=\mathrm{e}^b \Leftrightarrow a=b$ ; \quad $\mathrm{e}^a<\mathrm{e}^b \Leftrightarrow a<b$ ; \quad $\mathrm{e}^x=k \;(k>0) \Leftrightarrow x=\ln k$ \\ \hline
Tangente en 0 & $y=x+1$ (car $\mathrm{e}^0=1$ et $(\mathrm{e}^x)'(0)=1$) \\ \hline
\end{tabularx}
\end{center}

\textbf{Méthodes express}
\begin{itemize}
\item \cle{Résoudre} $\mathrm{e}^{u(x)}=\mathrm{e}^{v(x)}$ : égaliser les exposants $u(x)=v(x)$. Pour $\mathrm{e}^{u(x)}=k$ avec $k>0$ : $u(x)=\ln k$ ; si $k\leq 0$ : pas de solution.
\item \cle{Inéquation} $\mathrm{e}^{u(x)}>k$ ($k>0$) : $u(x)>\ln k$ (le sens est conservé car $\exp$ est croissante).
\item \cle{Dériver} $f(x)=\mathrm{e}^{u(x)}$ : identifier $u$, calculer $u'$, écrire $f'(x)=u'(x)\,\mathrm{e}^{u(x)}$.
\item \cle{Étudier le signe} de $f'$ : comme $\mathrm{e}^{u}>0$, le signe de $f'$ est celui de $u'$.
\item \cle{Primitive} de $\mathrm{e}^{ax+b}$ : diviser par $a$ ; vérifier en redérivant.
\item \cle{Modélisation} : une quantité évoluant à taux constant proportionnel à elle-même suit $f(t)=f(0)\,\mathrm{e}^{kt}$ ($k>0$ croissance, $k<0$ décroissance).
\end{itemize}
\end{bilan}

\begin{aretenir}[Checklist avant le Probatoire]
\begin{itemize}
\item $\square$ Je sais que $\mathrm{e}^0=1$, $\mathrm{e}^1=\mathrm{e}\approx \num{2,718}$ et $\mathrm{e}^x>0$ pour tout $x$.
\item $\square$ Je maîtrise les règles de calcul : $\mathrm{e}^{a+b}$, $\mathrm{e}^{a-b}$, $\mathrm{e}^{-a}$, $(\mathrm{e}^a)^n$.
\item $\square$ Je connais les limites en $+\infty$ et $-\infty$ et j'identifie l'asymptote horizontale $y=0$.
\item $\square$ Je sais dériver $\mathrm{e}^x$ et $\mathrm{e}^{u(x)}$ sans oublier le facteur $u'(x)$.
\item $\square$ Je sais déterminer une primitive de $\mathrm{e}^{ax+b}$ et de $u'\,\mathrm{e}^{u}$.
\item $\square$ Je résous les équations $\mathrm{e}^{u(x)}=k$ en distinguant les cas $k>0$ et $k\leq 0$.
\item $\square$ Je résous les inéquations exponentielles en conservant le sens des inégalités.
\item $\square$ Je dresse le tableau de variations complet d'une fonction avec $\exp$ (domaine, limites, dérivée, signe).
\item $\square$ Je sais écrire l'équation de la tangente : $y=f'(a)(x-a)+f(a)$.
\item $\square$ Je vérifie toujours mes primitives en les redérivant.
\item $\square$ Je sais modéliser une croissance ou décroissance continue par $f(t)=A\,\mathrm{e}^{kt}$.
\item $\square$ Je rédige mes justifications : théorème cité, conditions vérifiées, conclusion claire.
\end{itemize}
\end{aretenir}

\findecours

\end{document}
